%% vol3_ch1_5.tex — Chapters 1–5: Wittgenstein
%% Part I of Volume III: The Math of Ideas
%% Uses custom commands: \rs, \emerge, \compose, \void, \IS

%%%%%%%%%%%%%%%%%%%%%%%%%%%%%%%%%%%%%%%%%%%%%%%%%%%%%%%%%%%%%%%%%%%%%%%%
%%  CHAPTER 1 — LANGUAGE GAMES AS MEANING GAMES
%%%%%%%%%%%%%%%%%%%%%%%%%%%%%%%%%%%%%%%%%%%%%%%%%%%%%%%%%%%%%%%%%%%%%%%%

\chapter{Language Games as Meaning Games}
\label{ch:language-games}

\epigraph{``The meaning of a word is its use in the language.''}{Ludwig Wittgenstein, \textit{Philosophical Investigations}, \S43}

\section{From Tractatus to Investigations}

The trajectory of Wittgenstein's thought---from the crystalline logic of the
\textit{Tractatus Logico-Philosophicus} (1921) to the rough ground of the
\textit{Philosophical Investigations} (1953)---is one of the great
intellectual dramas of the twentieth century. The early Wittgenstein believed
that language \emph{pictures} reality: a proposition shares a logical form
with the state of affairs it represents, and the totality of true propositions
\emph{is} the world. The later Wittgenstein repudiated this picture theory
wholesale. Language does not mirror reality; it is a \emph{tool}, embedded in
human practices, shaped by context, and irreducible to any single logical form.

The concept through which the later Wittgenstein articulates this insight is the
\emph{language game} (\textit{Sprachspiel}). A language game is any activity
in which language is woven into action: giving orders, describing objects,
reporting events, speculating, singing, guessing riddles, telling jokes, asking,
thanking, cursing, greeting, praying (\textit{PI}, \S23). There is no essence
common to all language games, no single feature that makes them all ``language.''
Instead, they form a family, connected by overlapping similarities.

In this chapter, we show that the ideatic space formalized in Volume~I provides
the exact mathematical structure that Wittgenstein's language games require.
The key correspondences are:
\begin{itemize}
\item A \textbf{language game} is a context $c \in \IS$ together with rules $r \in \IS$.
\item A \textbf{move} in the game is the composition of the context with a player's utterance: $c \compose p$.
\item The \textbf{score} of a move is its resonance with the rules: $\rs(c \compose p, r)$.
\item A \textbf{form of life} is a background idea $f$ that frames all utterances within a community.
\item \textbf{Meaning as use} is captured by the resonance profile: two ideas with the same resonance profile have the same meaning.
\end{itemize}

These are not analogies. They are \emph{definitions} within a theory whose
consequences can be---and have been---machine-verified.

\section{The Formal Structure of Language Games}

\begin{definition}[Language Game]
\label{def:language-game}
Let $(\IS, \compose, \void, \rs)$ be an ideatic space. A \textbf{language game}
is specified by a \emph{context} $c \in \IS$ and \emph{rules} $r \in \IS$.
A \textbf{move} by player $p$ is the idea $c \compose p$. The \textbf{score}
of the move is
\[
\mathrm{score}(c, r, p) := \rs(c \compose p,\, r).
\]
\end{definition}

This definition captures the essential features of Wittgenstein's concept.
The context $c$ represents everything that has happened in the game so far---the
shared history, the established conventions, the physical setting. The rules $r$
represent the normative standards against which moves are evaluated. The
player's utterance $p$ is what the player contributes. Composition $c \compose p$
produces the new game state, and resonance $\rs(\cdot, r)$ measures how well
this new state conforms to the rules.

\begin{theorem}[Score Decomposition --- W5]
\label{thm:score-decompose}
For any context $c$, rules $r$, and player $p$:
\[
\mathrm{score}(c, r, p) = \rs(c, r) + \rs(p, r) + \emerge(c, p, r).
\]
\end{theorem}

\begin{proof}
This is an immediate consequence of the resonance decomposition theorem
proved in Volume~I (Theorem~I.2.1): for any $a, b, c \in \IS$,
$\rs(a \compose b, c) = \rs(a,c) + \rs(b,c) + \emerge(a,b,c)$.
Apply with $a = c$, $b = p$, $c = r$.
\end{proof}

\begin{interpretation}[The Three Components of Meaning]
\label{interp:three-components}
The score decomposition reveals that every utterance in a language game has
three components of meaning:
\begin{enumerate}
\item \textbf{Contextual resonance} $\rs(c, r)$: how well the existing context
already conforms to the rules. This is the ``baseline''---the game state before
the player speaks.
\item \textbf{Individual resonance} $\rs(p, r)$: the player's utterance's
intrinsic fit with the rules, independent of context. This is what the utterance
contributes ``on its own.''
\item \textbf{Emergence} $\emerge(c, p, r)$: the genuinely new meaning that
arises from the \emph{interaction} of context and utterance. This is Wittgenstein's
key insight: meaning is not merely the sum of parts. The same word in different
contexts has different emergence, and it is this emergence that makes language a
living practice rather than a dead calculus.
\end{enumerate}
When $\emerge = 0$ for all moves, the game reduces to a linear scoring system
where context and utterance contribute independently. This is the
\textit{Tractatus}'s world: meaning as logical combination, without genuine
novelty. The transition to the \textit{Investigations} is precisely the
recognition that $\emerge \neq 0$---that context transforms utterances in
ways irreducible to their parts.
\end{interpretation}

\begin{theorem}[Sequential Moves Associate --- W7]
\label{thm:sequential-assoc}
For any context $c$ and successive moves $p_1, p_2$:
\[
(c \compose p_1) \compose p_2 = c \compose (p_1 \compose p_2).
\]
\end{theorem}

\begin{proof}
By Axiom A1 (associativity of composition), proved in Volume~I as the
foundational monoid axiom of the ideatic space.
\end{proof}

\begin{interpretation}
This theorem formalizes Wittgenstein's observation that ``the game goes on''
(\textit{PI}, \S83). A language game is not a sequence of isolated moves but a
continuous practice. The fact that sequential moves associate means that the
game state after two moves is the same regardless of how we group them: first
$p_1$ then $p_2$, or the combined move $p_1 \compose p_2$ all at once. The
\emph{destination} is the same---but as we shall see in Chapter~6, the
\emph{emergence along the way} differs. The journey matters even when the
endpoint does not.
\end{interpretation}

\begin{theorem}[Moves Enrich --- W8]
\label{thm:moves-enrich}
Every move increases (or at least maintains) the weight of the game state:
\[
\rs(c \compose p,\, c \compose p) \geq \rs(c, c).
\]
\end{theorem}

\begin{proof}
By Axiom E4 (composition enriches), proved in Volume~I (Theorem~I.1.8):
$\rs(a \compose b, a \compose b) \geq \rs(a, a)$ for all $a, b \in \IS$.
\end{proof}

\begin{interpretation}[You Can't Un-Say Something]
\label{interp:cant-unsay}
Once a move is made in a language game, the game state's self-resonance
(``weight'' or ``presence'') can never decrease. This is the formal
counterpart of the everyday observation that you can't un-say something.
A retraction, an apology, a correction---these are all \emph{new} moves
that add to the game's complexity rather than erasing the original utterance.
In Wittgenstein's terms: every speech act leaves a trace in the form of life.
\end{interpretation}

The following Lean code verifies these results:

\begin{lstlisting}
-- From IDT_v3_Wittgenstein.lean

def gameMove (context player : I) : I := compose context player

noncomputable def gameScore (context rules player : I) : R :=
  rs (gameMove context player) rules

theorem gameScore_decompose (c r p : I) :
    gameScore c r p = rs c r + rs p r + emergence c p r := by
  unfold gameScore gameMove; exact rs_compose_eq c p r

theorem sequential_moves_assoc (c p1 p2 : I) :
    gameMove (gameMove c p1) p2 = gameMove c (compose p1 p2) := by
  unfold gameMove; exact compose_assoc' c p1 p2

theorem gameMove_enriches (c p : I) :
    rs (gameMove c p) (gameMove c p) >= rs c c := by
  unfold gameMove; exact compose_enriches' c p
\end{lstlisting}

\section{Forms of Life}

Wittgenstein insists that language games are embedded in \emph{forms of life}
(\textit{Lebensformen}): ``to imagine a language means to imagine a form of
life'' (\textit{PI}, \S19). A form of life is the totality of shared practices,
institutions, natural reactions, and background assumptions that make a language
game possible. It is not something we can step outside of or examine from a
neutral standpoint---it is the ground on which we stand.

\begin{definition}[Form of Life]
\label{def:form-of-life}
A \textbf{form of life} is an idea $f \in \IS$ that serves as the background
context for all utterances within a community. An utterance $a \compose b$
\emph{within form of life $f$} is the idea
\[
f \compose (a \compose b).
\]
\end{definition}

\begin{theorem}[Forms of Life Absorb --- W15]
\label{thm:form-absorb}
For any form of life $f$ and utterances $a, b$:
\[
f \compose (a \compose b) = (f \compose a) \compose b.
\]
\end{theorem}

\begin{proof}
By the associativity axiom A1, applied to $f$, $a$, $b$.
Lean: \texttt{formOfLife\_absorb} in \texttt{IDT\_v3\_Wittgenstein.lean}.
\end{proof}

\begin{interpretation}
The absorption theorem reveals that a form of life does not merely ``surround''
utterances---it \emph{enters into their constitution}. Framing the composition
$a \compose b$ by $f$ is the same as first framing $a$ by $f$ and then composing
with $b$. The form of life is not external scaffolding; it is an integral
part of every utterance's structure. This is exactly Wittgenstein's point:
meaning is not determined by the utterance alone but by the form of life
in which it is embedded.
\end{interpretation}

\begin{theorem}[Form of Life Enrichment --- W19]
\label{thm:form-enriches}
A non-trivial form of life always enriches:
\[
\rs\bigl(f \compose (a \compose b),\; f \compose (a \compose b)\bigr)
\geq \rs(f, f).
\]
\end{theorem}

\begin{proof}
By Axiom E4 applied to $f$ and $a \compose b$.
Lean: \texttt{formOfLife\_enriches} in \texttt{IDT\_v3\_Wittgenstein.lean}.
\end{proof}

\begin{theorem}[Incommensurable Forms of Life --- W22--W24]
\label{thm:incommensurable}
Define two forms of life $f_1, f_2$ as \textbf{incommensurable} if
$\rs(f_1, f_2) = 0$ and $\rs(f_2, f_1) = 0$. Then:
\begin{enumerate}
\item Incommensurability is symmetric.
\item The void form of life is incommensurable with everything.
\item Self-incommensurability implies $f = \void$ (only silence is
incommensurable with itself).
\end{enumerate}
\end{theorem}

\begin{proof}
(1) is immediate from the definition. (2) follows from the void-resonance
axioms R1--R2 (proved in Volume~I, Theorem~I.1.3). (3) follows from the
nondegeneracy axiom E2: if $\rs(f,f) = 0$ then $f = \void$.
\end{proof}

\begin{interpretation}
Incommensurability is Wittgenstein's recognition that different forms of life
may have nothing in common---no shared resonance, no mutual intelligibility.
But the theory also shows that incommensurability has strict limits: only
silence is incommensurable with itself. Any non-void form of life resonates
\emph{with itself}, and therefore has at least one connection to the space
of meaning. Total isolation is impossible for anything that exists.
\end{interpretation}

\section{Meaning as Use}

\begin{definition}[Same Use]
\label{def:same-use}
Two ideas $a, b \in \IS$ have the \textbf{same use} if they have identical
outgoing resonance profiles:
\[
a \sim_{\mathrm{use}} b \quad \iff \quad \forall c \in \IS,\; \rs(a, c) = \rs(b, c).
\]
\end{definition}

\begin{theorem}[Same Use as Void Implies Void --- W50]
\label{thm:same-use-void}
If $a$ has the same use as $\void$, then $a = \void$.
\end{theorem}

\begin{proof}
If $\rs(a, c) = \rs(\void, c) = 0$ for all $c$, then in particular
$\rs(a, a) = 0$, so $a = \void$ by nondegeneracy (Axiom E2).
Lean: \texttt{sameUse\_void\_is\_void} in \texttt{IDT\_v3\_Wittgenstein.lean}.
\end{proof}

\begin{interpretation}
This theorem is the formal content of Wittgenstein's meaning-as-use doctrine.
If an idea has the same use as silence---if it resonates with nothing---then it
\emph{is} silence. There is no ``hidden meaning'' behind use. An idea that
functions exactly like $\void$ in every possible context is $\void$. Meaning is
exhaustively determined by use, as Wittgenstein insisted.
\end{interpretation}

\section{The Builders' Game}

\begin{example}[Wittgenstein's Builders --- \textit{PI} \S2]
\label{ex:builders}
Wittgenstein opens the \textit{Investigations} with a primitive language game:
a builder and an assistant work with blocks, pillars, slabs, and beams. The
builder calls out a word; the assistant brings the corresponding item. In our
formalism, the builder's call is a move $p$ in a context $c$ (the construction
site), and the score $\rs(c \compose p, r)$ measures how well the resulting
state satisfies the building rules $r$.

By Theorem~\ref{thm:score-decompose}, the score decomposes:
$\mathrm{score} = \rs(c,r) + \rs(p,r) + \emerge(c,p,r)$.
In the simplest builder game, if the context $c$ is ``void'' (empty site) and
rules $r$ are ``void'' (no criteria), then every move scores zero
(Lean: \texttt{builders\_void\_game}). Without criteria, there is no
differentiation---no language game at all.
\end{example}

\begin{theorem}[Void Game Is Trivial --- W132]
\label{thm:void-game-trivial}
In the void game (void context, void rules), every move scores zero:
\[
\mathrm{score}(\void, \void, p) = 0 \quad \text{for all } p.
\]
\end{theorem}

\begin{proof}
$\rs(\void \compose p, \void) = \rs(p, \void) = 0$ by Axiom R1.
\end{proof}

\begin{proof}[Natural Language Proof of Void Game Triviality]
The void game is a language game with no context ($c = \void$) and no rules
($r = \void$). Every player's move $p$ is scored by $\rs(\void \compose p, \void)$.

Step~1: $\void \compose p = p$ by Axiom A2 (left identity). The void context
contributes nothing; the game state is just the player's utterance.

Step~2: $\rs(p, \void) = 0$ by Axiom R1 (right void resonance). No idea
resonates with void.

Therefore, $\mathrm{score}(\void, \void, p) = 0$ for every player $p$. Every
utterance is equally ``good'' (or equally meaningless) in the void game.

Alternatively, using the score decomposition (Theorem~\ref{thm:score-decompose}):
$\mathrm{score}(\void, \void, p) = \rs(\void, \void) + \rs(p, \void)
+ \emerge(\void, p, \void) = 0 + 0 + 0 = 0$, where the first two terms vanish
by R1--R2 and the third by the void-emergence identity (Volume~I, Theorem~I.2.3).

Philosophically, the void game is the limiting case that makes normativity
visible by its absence. In a game with no rules, every move is permissible
and none is required. There is no ``right'' or ``wrong'' utterance. This is
not a defective game---it is the \emph{absence} of a game. Wittgenstein's
builder game, by contrast, has rules ($r \neq \void$) and context ($c \neq \void$),
which is precisely what makes it a genuine language game. The void game theorem
establishes the zero point against which all genuine games are measured.
\end{proof}

\begin{interpretation}
The void game is the degenerate case: a language game with no context and no
rules is not a language game at all. This formalizes Wittgenstein's insistence
that meaning requires a normative context---rules, practices, institutions.
Without them, every utterance is equally ``correct'' (i.e., equally
meaningless).
\end{interpretation}

\section{The Emergence of Depth Grammar}

Wittgenstein distinguishes between \emph{surface grammar}---the apparent
syntactic form of an utterance---and \emph{depth grammar}---the actual rules
governing its use (\textit{PI}, \S664). In IDT, depth grammar is precisely the
emergence structure.

\begin{theorem}[Emergence Determines Depth --- W10, W14]
\label{thm:depth-grammar}
Two moves $p_1, p_2$ in the same context $c$ with the same rules $r$ differ in
depth grammar if and only if they produce different emergence:
\[
\emerge(c, p_1, r) \neq \emerge(c, p_2, r) \implies p_1 \neq p_2.
\]
Conversely, $p_1 = p_2$ implies $\emerge(c, p_1, r) = \emerge(c, p_2, r)$.
\end{theorem}

\begin{proof}
The forward direction is the contrapositive of the converse. The converse
holds because emergence is a function of its arguments: equal inputs yield
equal outputs.
Lean: \texttt{gameScore\_player\_diff}, \texttt{game\_emergence\_bounded}.
\end{proof}

\begin{proof}[Natural Language Proof of Score Decomposition (Theorem~\ref{thm:score-decompose})]
We now give a complete natural-language proof of the score decomposition, the
foundational result of this chapter.

The score of player $p$'s move in context $c$ under rules $r$ is
$\mathrm{score}(c, r, p) = \rs(c \compose p, r)$. By the resonance
decomposition theorem (Volume~I, Theorem~I.2.1), which states that for any
$a, b, c \in \IS$:
\[
\rs(a \compose b, c) = \rs(a, c) + \rs(b, c) + \emerge(a, b, c),
\]
we substitute $a = c$ (the context), $b = p$ (the player's move), and
the third argument as $r$ (the rules). This yields:
\[
\rs(c \compose p, r) = \rs(c, r) + \rs(p, r) + \emerge(c, p, r).
\]
The three terms have clear Wittgensteinian interpretations. The first term
$\rs(c, r)$ is the game's baseline---how well the context already satisfies
the rules before the player speaks. This is the shared background, the
``form of life'' already in place. The second term $\rs(p, r)$ is the
player's intrinsic fitness---how well their utterance fits the rules
regardless of context. A grammatical sentence scores positively here even
in isolation. The third term $\emerge(c, p, r)$ is the \emph{depth grammar}:
the genuinely new meaning that arises from the \emph{interaction} of context
and utterance. The same word ``bank'' has different emergence in a financial
context versus a riverside context.

The decomposition reveals that Wittgenstein's ``meaning is use'' doctrine
has three components, not one. Meaning-as-use includes not just the
utterance's profile ($\rs(p, r)$) but the contextual baseline and the
emergent interaction. The late Wittgenstein's revolution is the recognition
that the third term---emergence---is where philosophical action lives.
\end{proof}

\begin{theorem}[Double Void Game --- W13]
\label{thm:double-void}
In a void context with a void player, the game score is zero regardless
of rules:
\[
\rs(\void \compose \void, r) = \rs(\void, r) = 0.
\]
\end{theorem}

\begin{proof}
$\void \compose \void = \void$ by Axiom A3, and $\rs(\void, r) = 0$ by R2.
Lean: \texttt{game\_double\_void} in \texttt{IDT\_v3\_Wittgenstein.lean}.
\end{proof}

\begin{interpretation}[Where IDT Extends Wittgenstein]
\label{interp:extends-wittgenstein}
Wittgenstein's language game concept is deliberately informal---he refuses to
give a definition, offering instead a ``family'' of examples (\textit{PI},
\S\S23, 65--67). The IDT formalization does not betray this anti-definitional
spirit; rather, it reveals the \emph{structural invariants} that hold across
all language games regardless of their particular content.

Wittgenstein could not see these invariants because he lacked the algebraic
language. He knew that context matters (depth grammar), that meaning is use
(resonance profiles), and that games compose (the form of life absorbs).
But he could not state the \emph{cocycle condition}---the precise constraint
linking emergence across levels of composition---because he had no concept
of emergence as a defined mathematical quantity. The cocycle condition
(Volume~I, Theorem~I.3.1) is genuinely new: it constrains language games in
ways that Wittgenstein never articulated, and it does so necessarily, as a
consequence of the axioms.

This is the value of formalization: it reveals what was always implicit in
the philosopher's insights but could not be stated in natural language.
\end{interpretation}

\begin{remark}[Cross-Reference to Volume~I]
The score decomposition theorem depends on the resonance decomposition
(Volume~I, Theorem~I.2.1), which in turn depends on Axioms A1--A3
(monoid structure) and the definition of emergence $\emerge(a,b,c) :=
\rs(a \compose b, c) - \rs(a,c) - \rs(b,c)$. The entire theory of language
games is thus grounded in the compositional structure of the ideatic space
established in Volume~I, Chapter~2.
\end{remark}

\section{Void Moves and the Silence of the Game}

The void idea plays a distinguished role in every language game. In Wittgenstein's
terms, $\void$ is \emph{silence}---the absence of utterance, the move that is
not made. Several theorems establish its precise algebraic behavior.

\begin{theorem}[Void Player Leaves Context Unchanged --- W4]
\label{thm:void-player}
For any context $c$:
\[
c \compose \void = c.
\]
\end{theorem}

\begin{proof}
By Axiom A3 (right identity of composition). The void idea composes trivially
from the right: appending silence to any context leaves the context as it was.
Lean: \texttt{gameMove\_void\_player} in \texttt{IDT\_v3\_Wittgenstein.lean}.
\end{proof}

\begin{theorem}[Void Context Lets the Player Speak Freely --- W3]
\label{thm:void-context}
For any player $p$:
\[
\void \compose p = p.
\]
\end{theorem}

\begin{proof}
By Axiom A2 (left identity of composition). In a game with no prior context,
the player's utterance \emph{is} the entire game state. There is nothing to
interact with, so the utterance stands alone.
Lean: \texttt{gameMove\_void\_context} in \texttt{IDT\_v3\_Wittgenstein.lean}.
\end{proof}

\begin{interpretation}[The Grammar of Silence]
These two theorems---void player and void context---establish the grammar of
silence within language games. Silence is the \emph{neutral element} of
conversation: adding it to either side changes nothing. This is not a trivial
observation. Wittgenstein writes: ``Whereof one cannot speak, thereof one must
be silent'' (\textit{Tractatus}, 7). In IDT, silence has a precise algebraic
characterization: it is the unique idea that, when composed with any other idea,
produces that idea unchanged.

As established in Volume~I, Theorem~I.1.2 (void uniqueness), the void idea is
the \emph{unique} element with this property. There is exactly one way to be
silent, but infinitely many ways to speak. This asymmetry---one silence, many
utterances---is the foundation of the entire theory.
\end{interpretation}

\begin{theorem}[Void Rules Accept Everything --- W2]
\label{thm:void-rules}
In a game with void rules, every move scores zero:
\[
\rs(c \compose p, \void) = 0 \quad \text{for all } c, p.
\]
\end{theorem}

\begin{proof}
By Axiom R1 (right void resonance), $\rs(a, \void) = 0$ for all $a$.
Taking $a = c \compose p$ yields the result.
Lean: \texttt{gameScore\_void\_rules} in \texttt{IDT\_v3\_Wittgenstein.lean}.
\end{proof}

\begin{theorem}[Silent Player Scores the Baseline --- W1]
\label{thm:silent-player}
If the player is silent, the score equals the context's resonance with the rules:
\[
\mathrm{score}(c, r, \void) = \rs(c, r).
\]
\end{theorem}

\begin{proof}
$\rs(c \compose \void, r) = \rs(c, r)$ by Axiom A3.
Lean: \texttt{gameScore\_void\_player} in \texttt{IDT\_v3\_Wittgenstein.lean}.
\end{proof}

\begin{proof}[Natural Language Proof of the Silent Player Theorem]
When the player says nothing ($p = \void$), the game state is $c \compose \void = c$
(by right identity). The score is then $\rs(c, r)$---purely the pre-existing
context's conformity with the rules. There is no individual resonance contribution
$\rs(\void, r) = 0$ (by R2) and no emergence $\emerge(c, \void, r) = 0$
(since the void generates no interaction). The silent player scores the
baseline---exactly the score they would get by being absent. Philosophically,
this formalizes the idea that silence is not a ``zero move'' but rather the
\emph{absence} of a move: the game state is simply whatever it was before.
\end{proof}

\begin{interpretation}[Brandom and the Space of Reasons]
Robert Brandom's \textit{Making It Explicit} (1994) develops a theory of meaning
as inferential role: the meaning of a sentence is its position in the space of
reasons---the inferences it licenses and the inferences that license it. In
IDT terms, Brandom's inferential role is captured by the resonance profile
$\{\rs(a, c)\}_{c \in \IS}$: the complete set of resonance values between
idea $a$ and every other idea $c$.

The score decomposition theorem (Theorem~\ref{thm:score-decompose}) reveals
that Brandom's picture, while correct as far as it goes, misses the
\emph{emergence} component. An idea's inferential role ($\rs(a, c)$ for various
$c$) captures only the individual resonance term. The contextual resonance
$\rs(c, r)$ and the emergence $\emerge(c, a, r)$ are additional dimensions of
meaning that Brandom's framework does not distinguish. The emergence, in
particular, is the \emph{creative} contribution of the idea in context---precisely
what makes linguistic practice a living activity rather than a static web of
inferential connections.

As established in Volume~I, Theorem~I.2.1 (resonance decomposition), the
three-way decomposition is exhaustive: there is no ``fourth component'' of
meaning. Brandom captures one-third of the story; Wittgenstein, with his
emphasis on context and depth grammar, gestures toward the other two-thirds.
\end{interpretation}

\begin{theorem}[Player Comparison --- W9]
\label{thm:player-comparison}
Two players in the same context, scored by the same rules, differ by:
\[
\mathrm{score}(c, r, p_1) - \mathrm{score}(c, r, p_2)
= \bigl(\rs(p_1, r) - \rs(p_2, r)\bigr)
  + \bigl(\emerge(c, p_1, r) - \emerge(c, p_2, r)\bigr).
\]
\end{theorem}

\begin{proof}
Apply Theorem~\ref{thm:score-decompose} to both scores and subtract. The
common term $\rs(c, r)$ cancels, leaving only the player-dependent terms.
Lean: \texttt{gameScore\_player\_diff} in \texttt{IDT\_v3\_Wittgenstein.lean}.
\end{proof}

\begin{proof}[Natural Language Proof of Player Comparison]
We wish to show that the difference in scores between two players in the same
language game depends only on (i) the difference in their individual fitness
and (ii) the difference in the emergence they produce.

By the score decomposition, $\mathrm{score}(c,r,p_1) = \rs(c,r) + \rs(p_1,r)
+ \emerge(c,p_1,r)$ and $\mathrm{score}(c,r,p_2) = \rs(c,r) + \rs(p_2,r)
+ \emerge(c,p_2,r)$. Subtracting eliminates $\rs(c,r)$ and yields the stated
formula.

Philosophically, the player comparison theorem shows that competitive advantage
in a language game comes from two sources. First, the player's \emph{intrinsic}
resonance with the rules ($\rs(p,r)$)---how well the utterance fits the norms
independent of context. Second, the player's \emph{contextual} emergence
($\emerge(c,p,r)$)---how creatively the utterance interacts with the existing
context. A mediocre utterance in a rich context may outperform a brilliant
utterance in a barren one, because emergence depends on the richness of
the interaction, not just the quality of the parts.
\end{proof}

\begin{theorem}[Triple Move Associativity --- W11]
\label{thm:triple-moves}
For any context $c$ and three successive moves $p_1, p_2, p_3$:
\[
((c \compose p_1) \compose p_2) \compose p_3
= c \compose ((p_1 \compose p_2) \compose p_3)
= c \compose (p_1 \compose (p_2 \compose p_3)).
\]
\end{theorem}

\begin{proof}
By iterated application of the associativity axiom A1. The first equality
applies A1 to $(c, p_1, p_2 \compose p_3)$ after regrouping; the second
applies A1 to $(p_1, p_2, p_3)$.
Lean: \texttt{triple\_moves\_assoc} in \texttt{IDT\_v3\_Wittgenstein.lean}.
\end{proof}

\begin{interpretation}[Travis on Occasion-Sensitivity]
Charles Travis's work on occasion-sensitivity (\textit{The Uses of Sense}, 1989;
\textit{Occasion-Sensitivity}, 2008) argues that the same sentence can express
different things on different occasions of use, and that this variability is not
a defect of language but an essential feature. The truth-conditions of ``the
leaves are green'' depend on whether we are assessing visual appearance or
chemical composition.

In IDT, Travis's occasion-sensitivity is captured by the emergence term
$\emerge(c, p, r)$. The same utterance $p$ (``the leaves are green'') in
different contexts $c_1$ (botanical survey) and $c_2$ (painting class) produces
different emergence: $\emerge(c_1, p, r) \neq \emerge(c_2, p, r)$ in general.
The score decomposition makes this precise: the occasion-sensitive component
of meaning is \emph{exactly} the emergence---the part that depends on the
interaction of utterance and context.

Travis's insight is that meaning is never determined by the sentence alone.
The formalism vindicates this: even when $\rs(p, r)$ is fixed (the sentence's
``literal meaning'' is constant), the total score varies with context through
$\rs(c, r) + \emerge(c, p, r)$. Travis is right that there is no such thing as
``the meaning'' of a sentence in isolation---there is only its meaning on an
occasion, which includes the ineliminable emergence component.
\end{interpretation}

\begin{theorem}[Game Emergence Is Bounded --- W12]
\label{thm:game-emergence-bounded}
For any context $c$, rules $r$, and player $p$:
\[
\emerge(c, p, r)^2 \leq \rs(c \compose p, c \compose p) \cdot \rs(r, r).
\]
\end{theorem}

\begin{proof}
By Axiom E3 (emergence bounded), applied to the triple $(c, p, r)$.
The emergence of a move in a language game cannot exceed the geometric
mean of the game state's weight and the rules' weight.
Lean: \texttt{game\_emergence\_bounded} in \texttt{IDT\_v3\_Wittgenstein.lean}.
\end{proof}

\begin{proof}[Natural Language Proof of the Emergence Bound]
The emergence bound is a Cauchy--Schwarz-type inequality for the ideatic space.
It states that the creative interaction between context and player, as measured
by the emergence $\emerge(c, p, r)$, is constrained by the ``sizes'' of the
game state and the rules.

Formally, squaring both sides of any putative emergence value and comparing with
the product $\rs(c \compose p, c \compose p) \cdot \rs(r, r)$ yields the bound.
The proof in Volume~I (Theorem~I.1.7) proceeds by the standard Cauchy--Schwarz
technique: construct the nonnegative quadratic form
$\rs(c \compose p + t \cdot r, c \compose p + t \cdot r) \geq 0$ for all $t$,
expand using bilinearity, and apply the nonnegativity condition to the discriminant.

Philosophically, the emergence bound says that even Wittgenstein's ``depth
grammar''---the creative, context-dependent, emergent component of meaning---is
not unbounded. It is constrained by the weight of the game state and the weight
of the rules. A game with lightweight rules (small $\rs(r,r)$) cannot support
large emergence. A game state with little accumulated weight likewise constrains
emergence. Creativity requires substance---both in the context and in the norms.
\end{proof}

\section{The Form-of-Life Cocycle and Incommensurability}

The interaction of forms of life with utterances obeys a deep algebraic constraint
that neither Wittgenstein nor his commentators could articulate without the
formal apparatus.

\begin{theorem}[Form-of-Life Cocycle --- W20]
\label{thm:form-cocycle}
For any form of life $f$, utterances $a, b$, and observer $d$:
\[
\emerge(f, a, d) + \emerge(f \compose a, b, d)
= \emerge(a, b, d) + \emerge(f, a \compose b, d).
\]
\end{theorem}

\begin{proof}
This is the cocycle condition (Volume~I, Theorem~I.3.1) applied to the triple
$(f, a, b)$ with observer $d$. The cocycle condition states that emergence
satisfies an additive constraint when three ideas are composed in sequence:
the emergence of composing $f$ with $a$ then composing the result with $b$
must decompose in the prescribed way.

Lean: \texttt{formOfLife\_cocycle} in \texttt{IDT\_v3\_Wittgenstein.lean}.
\end{proof}

\begin{proof}[Natural Language Proof of the Form-of-Life Cocycle]
To prove the cocycle condition, we expand both sides using the definition of
emergence ($\emerge(x, y, d) := \rs(x \compose y, d) - \rs(x, d) - \rs(y, d)$).

Left side: $\emerge(f, a, d) + \emerge(f \compose a, b, d) = [\rs(f \compose a, d)
- \rs(f, d) - \rs(a, d)] + [\rs((f \compose a) \compose b, d) - \rs(f \compose a, d)
- \rs(b, d)]$. The $\rs(f \compose a, d)$ terms cancel, yielding $\rs((f \compose a)
\compose b, d) - \rs(f, d) - \rs(a, d) - \rs(b, d)$.

Right side: $\emerge(a, b, d) + \emerge(f, a \compose b, d) = [\rs(a \compose b, d)
- \rs(a, d) - \rs(b, d)] + [\rs(f \compose (a \compose b), d) - \rs(f, d)
- \rs(a \compose b, d)]$. The $\rs(a \compose b, d)$ terms cancel, yielding
$\rs(f \compose (a \compose b), d) - \rs(f, d) - \rs(a, d) - \rs(b, d)$.

By associativity (A1), $(f \compose a) \compose b = f \compose (a \compose b)$,
so both sides are equal. The cocycle condition is thus a \emph{necessary consequence}
of associativity and the definition of emergence.

Philosophically, the cocycle reveals that a form of life's interaction with speech
is constrained: the emergence produced by $f$ when utterances are composed
stepwise must equal the emergence produced when utterances are composed all at
once. This is a \emph{conservation law} for meaning---emergence is ``reshuffled''
but never ``created or destroyed'' across different bracketings. It is the
algebraic content of Wittgenstein's claim that the form of life is not merely
a backdrop but an active participant in meaning-making, subject to structural
constraints that govern how it participates.
\end{proof}

\begin{theorem}[Two Forms Produce Different Resonance --- W21]
\label{thm:two-forms-diff}
For two forms of life $f_1, f_2$, any utterances $a, b$, and observer $c$:
\[
\rs(f_1 \compose (a \compose b), c) - \rs(f_2 \compose (a \compose b), c)
= \bigl(\rs(f_1, c) - \rs(f_2, c)\bigr)
  + \bigl(\emerge(f_1, a \compose b, c) - \emerge(f_2, a \compose b, c)\bigr).
\]
\end{theorem}

\begin{proof}
Apply the resonance decomposition to both terms on the left and subtract.
Lean: \texttt{two\_forms\_diff} in \texttt{IDT\_v3\_Wittgenstein.lean}.
\end{proof}

\begin{interpretation}[McDowell's Mind and World]
John McDowell's \textit{Mind and World} (1994) argues against the ``myth of the
Given''---the idea that experience provides a raw, unstructured input that the
mind then conceptualizes. McDowell insists that experience is already
\emph{conceptually structured}: perception is a form of \emph{receptivity}
already infused with \emph{spontaneity}.

The form-of-life cocycle formalizes McDowell's insight. The form of life $f$ is
not applied \emph{after} the composition $a \compose b$ is formed; it enters
into the very constitution of the composition. The cocycle condition shows
that the emergence produced by $f$ interacting with $a \compose b$ is
structurally linked to the emergence of $f$ with $a$ separately. There is no
``raw'' composition that exists prior to the form of life's involvement---the
form of life is always already there, shaping the emergence.

McDowell's second nature---the trained perceptual capacities that make
experience possible---is, in IDT, the form of life $f$. It is not a filter
applied to pre-given data but a constituent of the data itself. As established
in Volume~I, Theorem~I.3.1 (cocycle condition), this structural intertwining
of form and content is not optional; it is a necessary consequence of the axioms.
\end{interpretation}

\section{Full Use-Equivalence and Connotation}

The meaning-as-use doctrine admits a finer analysis when we consider not only
outgoing resonance profiles but also incoming resonance and compositional behavior.

\begin{theorem}[Same Use Is an Equivalence Relation --- W49--W52]
\label{thm:same-use-equiv}
The relation $a \sim_{\mathrm{use}} b$ (i.e., $\forall c,\, \rs(a,c) = \rs(b,c)$) is:
\begin{enumerate}
\item Reflexive: $a \sim_{\mathrm{use}} a$.
\item Symmetric: $a \sim_{\mathrm{use}} b \implies b \sim_{\mathrm{use}} a$.
\item Transitive: $a \sim_{\mathrm{use}} b$ and $b \sim_{\mathrm{use}} c$ imply $a \sim_{\mathrm{use}} c$.
\end{enumerate}
\end{theorem}

\begin{proof}
(1) $\rs(a,c) = \rs(a,c)$ for all $c$. (2) If $\rs(a,c) = \rs(b,c)$ for all $c$,
then $\rs(b,c) = \rs(a,c)$ for all $c$ by symmetry of equality. (3) If
$\rs(a,c) = \rs(b,c)$ and $\rs(b,c) = \rs(c',c)$ for all $c$, then
$\rs(a,c) = \rs(c',c)$ by transitivity of equality.
Lean: \texttt{sameUse\_refl}, \texttt{sameUse\_symm}, \texttt{sameUse\_trans}
in \texttt{IDT\_v3\_Wittgenstein.lean}.
\end{proof}

\begin{definition}[Full Use-Equivalence]
\label{def:full-use-equiv}
Two ideas $a, b$ are \textbf{fully use-equivalent} if they have identical
resonance profiles in \emph{both} directions:
\[
a \approx b \quad \iff \quad \forall c,\; \rs(a,c) = \rs(b,c) \text{ and } \rs(c,a) = \rs(c,b).
\]
\end{definition}

\begin{theorem}[Full Use-Equivalence Preserves Weight --- W55]
\label{thm:full-use-weight}
If $a \approx b$, then $\rs(a, a) = \rs(b, b)$.
\end{theorem}

\begin{proof}
Taking $c = a$ in the first condition: $\rs(a, a) = \rs(b, a)$. Taking
$c = a$ in the second condition: $\rs(a, a) = \rs(a, b)$. Taking $c = b$
in the first condition: $\rs(a, b) = \rs(b, b)$. By transitivity:
$\rs(a, a) = \rs(b, b)$.
Lean: \texttt{fullUseEquiv\_same\_weight} in \texttt{IDT\_v3\_Wittgenstein.lean}.
\end{proof}

\begin{theorem}[Full Void-Equivalence Implies Void --- W56]
\label{thm:full-void-equiv}
If $a \approx \void$, then $a = \void$.
\end{theorem}

\begin{proof}
By full use-equivalence, $\rs(a, c) = \rs(\void, c) = 0$ for all $c$.
In particular $\rs(a, a) = 0$, so $a = \void$ by nondegeneracy.
Lean: \texttt{fullUseEquiv\_void\_is\_void} in \texttt{IDT\_v3\_Wittgenstein.lean}.
\end{proof}

\begin{theorem}[Same-Use Connotation Invariance --- W54]
\label{thm:connotation-invariance}
If $a \sim_{\mathrm{use}} b$, then for all $c$:
\[
\emerge(a, c, d) - \emerge(b, c, d) = \rs(a \compose c, d) - \rs(b \compose c, d).
\]
That is, the difference in emergence is entirely determined by the difference
in compositional resonance.
\end{theorem}

\begin{proof}
$\emerge(a,c,d) = \rs(a \compose c, d) - \rs(a,d) - \rs(c,d)$ and
$\emerge(b,c,d) = \rs(b \compose c, d) - \rs(b,d) - \rs(c,d)$. Since
$\rs(a,d) = \rs(b,d)$ (by same-use), the $\rs(\cdot, d)$ and $\rs(c,d)$ terms
cancel on subtraction.
Lean: \texttt{sameUse\_connotation} in \texttt{IDT\_v3\_Wittgenstein.lean}.
\end{proof}

\begin{interpretation}[Diamond on the Realistic Spirit]
Cora Diamond's \textit{The Realistic Spirit} (1991) emphasizes Wittgenstein's
insistence that philosophy must attend to the \emph{actual use} of words rather
than constructing theoretical models of meaning. Diamond argues that philosophical
confusions arise when we project a ``meaning body'' behind the sign, rather than
examining how the sign functions in practice.

The full use-equivalence theorems formalize this realistic spirit. Two ideas that
have identical resonance profiles---the same use in every possible context---are
indistinguishable by the formalism. There is no ``meaning body'' to be found
behind the resonance profile: the profile \emph{is} the meaning. Full
void-equivalence (Theorem~\ref{thm:full-void-equiv}) is the extreme case:
an idea whose use is identical to silence \emph{is} silence.

Yet the connotation invariance theorem (Theorem~\ref{thm:connotation-invariance})
reveals a subtlety that Diamond might appreciate: even when two ideas have the
same \emph{individual} use ($\rs(a,c) = \rs(b,c)$ for all $c$), they may
differ in their \emph{compositional} behavior---in the emergence they produce
when combined with other ideas. This ``connotative surplus'' is the formal
residue of what Frege called \emph{Ton} (coloring) and what literary critics
call \emph{connotation}: the dimension of meaning that exceeds denotation but
is nonetheless real.
\end{interpretation}

\section{The Tractarian and the Investigative: Linear vs.\ Emergent Language}

Wittgenstein's own intellectual trajectory---from the \textit{Tractatus} to
the \textit{Investigations}---can be formalized as a transition from linear to
nonlinear composition.

\begin{definition}[Tractarian Idea]
\label{def:tractarian}
An idea $a$ is \textbf{Tractarian} if it is left-linear: $\emerge(a, b, c) = 0$
for all $b, c$.
\end{definition}

\begin{theorem}[Tractarian Ideas Are Compositional --- W130]
\label{thm:tractarian-comp}
If $a$ is Tractarian, then:
\[
\rs(a \compose b, c) = \rs(a, c) + \rs(b, c) \quad \text{for all } b, c.
\]
\end{theorem}

\begin{proof}
The resonance decomposition gives $\rs(a \compose b, c) = \rs(a,c) + \rs(b,c)
+ \emerge(a,b,c)$. Since $a$ is Tractarian, $\emerge(a,b,c) = 0$, and the
result follows.
Lean: \texttt{tractarian\_compositional} in \texttt{IDT\_v3\_Wittgenstein.lean}.
\end{proof}

\begin{theorem}[The Investigations Go Beyond the Tractatus --- W131]
\label{thm:investigations-beyond}
If $a$ has emergence (i.e., there exist $b, c$ with $\emerge(a, b, c) \neq 0$),
then $a$ is \emph{not} Tractarian.
\end{theorem}

\begin{proof}
Contrapositive of the Tractarian definition.
Lean: \texttt{investigations\_beyond\_tractatus} in
\texttt{IDT\_v3\_Wittgenstein.lean}.
\end{proof}

\begin{interpretation}[The Two Wittgensteins as Algebraic Regimes]
The early Wittgenstein's picture theory of meaning corresponds to the Tractarian
regime: composition is linear, meaning is the sum of parts, and there is no
emergence. The logical atomism of the \textit{Tractatus}---facts as combinations
of objects, propositions as pictures of facts---is precisely the theory that
$\emerge = 0$ everywhere.

The late Wittgenstein's rejection of the picture theory corresponds to the
recognition that $\emerge \neq 0$: context transforms meaning in ways that
cannot be predicted from the parts alone. Language games, forms of life, depth
grammar---all of these are names for the phenomenon of nonzero emergence.

The formalism thus reveals that the ``transition'' from the \textit{Tractatus}
to the \textit{Investigations} is not a change of topic but a change of
\emph{algebraic regime}: from the linear to the nonlinear, from the compositional
to the emergent. This is what makes the transition philosophically significant
rather than merely autobiographical: it corresponds to a genuine mathematical
bifurcation between two qualitatively different types of meaning structure.
\end{interpretation}

\section{Meaning as Said-Plus-Shown}

The \textit{Tractatus} distinguishes between what can be \emph{said} and what
can only be \emph{shown} (4.1212). In the \textit{Investigations}, this
distinction reappears as the gap between surface grammar and depth grammar.
The IDT formalism captures both distinctions through the concept of
\emph{shown content}.

\begin{definition}[Shown Content]
The \textbf{shown content} of $a$ composed with $b$, observed by $c$, is
the emergence:
\[
\mathrm{shown}(a, b, c) := \emerge(a, b, c) = \rs(a \compose b, c) - \rs(a, c) - \rs(b, c).
\]
\end{definition}

\begin{theorem}[Meaning Equals Said Plus Shown --- W128]
\label{thm:said-plus-shown}
For any $a, b, c \in \IS$:
\[
\rs(a \compose b, c) = \underbrace{\rs(a, c) + \rs(b, c)}_{\text{what is said}}
+ \underbrace{\emerge(a, b, c)}_{\text{what is shown}}.
\]
\end{theorem}

\begin{proof}
This is the resonance decomposition (Volume~I, Theorem~I.2.1) restated.
Lean: \texttt{meaning\_eq\_said\_plus\_shown} in
\texttt{IDT\_v3\_Wittgenstein.lean}.
\end{proof}

\begin{theorem}[Void Shows Nothing --- W129a--b]
\label{thm:void-shows-nothing}
\begin{enumerate}
\item $\mathrm{shown}(\void, b, c) = 0$ for all $b, c$.
\item $\mathrm{shown}(a, \void, c) = 0$ for all $a, c$.
\end{enumerate}
\end{theorem}

\begin{proof}
(1) $\emerge(\void, b, c) = \rs(\void \compose b, c) - \rs(\void, c) - \rs(b,c)
= \rs(b,c) - 0 - \rs(b,c) = 0$. (2) Similarly using A3 and R1.
Lean: \texttt{shown\_void\_left}, \texttt{shown\_void\_right} in
\texttt{IDT\_v3\_Wittgenstein.lean}.
\end{proof}

\begin{theorem}[The Showing Gap --- W129c]
\label{thm:showing-gap}
Two compositions with the same resonance profile but different shown content
differ in their surface contributions:
\[
\mathrm{shown}(a_1, b_1, c) - \mathrm{shown}(a_2, b_2, c)
= \bigl[\rs(a_1 \compose b_1, c) - \rs(a_2 \compose b_2, c)\bigr]
  - \bigl[\rs(a_1, c) + \rs(b_1, c)\bigr]
  + \bigl[\rs(a_2, c) + \rs(b_2, c)\bigr].
\]
\end{theorem}

\begin{proof}
Direct subtraction of the emergence definitions.
Lean: \texttt{showing\_gap} in \texttt{IDT\_v3\_Wittgenstein.lean}.
\end{proof}

\begin{interpretation}[Cavell on Acknowledgment and Knowledge]
Stanley Cavell's \textit{The Claim of Reason} (1979) distinguishes between
\emph{knowledge} (propositional grasp of facts) and \emph{acknowledgment}
(responsive engagement with persons and situations). Cavell argues that the
skeptic's error is demanding knowledge where only acknowledgment is possible:
you cannot \emph{know} that another person is in pain, but you can
\emph{acknowledge} their pain.

The said-plus-shown decomposition maps onto this distinction. What is
\emph{said}---the sum $\rs(a,c) + \rs(b,c)$---corresponds to propositional
content, to knowledge, to what can be explicitly stated. What is
\emph{shown}---the emergence $\emerge(a,b,c)$---corresponds to acknowledgment,
to the responsive, context-dependent, irreducibly interactive dimension of
meaning.

Cavell's point is that acknowledgment cannot be reduced to knowledge. The
formalism agrees: emergence cannot be reduced to the sum of individual resonances.
The shown content is genuinely new---it arises from the interaction of ideas in
context, not from any property of the ideas taken separately. As established in
Volume~I, Remark~I.3.2, the cocycle condition shows that emergence is
redistributed but never eliminated across different bracketings. Acknowledgment,
like emergence, is an ineliminable structural feature of meaning.
\end{interpretation}

\section{Idea Worlds and Best Responses}

\begin{definition}[Idea World]
The \textbf{idea world} of $a$ evaluated at $c$ is the resonance
$\rs(a, c)$---what $a$ ``sees'' in $c$.
\end{definition}

\begin{theorem}[Void Has an Empty World]
\label{thm:void-empty-world}
$\rs(\void, c) = 0$ for all $c$: the void idea sees nothing.
Lean: \texttt{void\_empty\_world} in \texttt{IDT\_v3\_Wittgenstein.lean}.
\end{theorem}

\begin{theorem}[Non-Void Ideas See Themselves]
\label{thm:nonvoid-sees-self}
If $a \neq \void$, then $\rs(a, a) > 0$: every non-void idea resonates
positively with itself.
Lean: \texttt{nonvoid\_sees\_self} in \texttt{IDT\_v3\_Wittgenstein.lean}.
\end{theorem}

\begin{theorem}[Composition Expands the World]
\label{thm:composition-expands}
$\rs(a \compose b, a \compose b) \geq \rs(a, a)$: composing an idea with
another expands (or at least preserves) its self-resonance.
Lean: \texttt{composition\_expands\_world} in \texttt{IDT\_v3\_Wittgenstein.lean}.
\end{theorem}

\begin{definition}[Best Response]
In a language game $(c, r)$, a move $m$ is a \textbf{best response} among
vocabulary $V = \{v_1, \ldots, v_n\}$ if:
\[
\rs(c \compose m, r) \geq \rs(c \compose v_i, r) \quad \text{for all } v_i \in V.
\]
\end{definition}

\begin{theorem}[In a Void Game, Everything Is Best --- W133a]
\label{thm:void-game-best}
If the rules are $\void$, every move is a best response.
\end{theorem}

\begin{proof}
$\rs(c \compose m, \void) = 0 = \rs(c \compose v_i, \void)$ for all $m, v_i$
by Axiom R1. All scores are zero, so every move ties.
Lean: \texttt{void\_game\_all\_best} in \texttt{IDT\_v3\_Wittgenstein.lean}.
\end{proof}

\begin{interpretation}
Without rules, there is no way to distinguish good from bad moves. This formalizes
Wittgenstein's point that normativity is essential to language games: ``In the
actual use of expressions we make detours, we go by side-roads. We \emph{see}
the straight highway before us, but of course we cannot use it, because it is
permanently closed'' (\textit{PI}, \S426). The ``straight highway'' of a
rule-less game is useless because it goes everywhere equally---which is to say,
nowhere.
\end{interpretation}

\begin{remark}[Cross-Reference to Volume~I]
The theorems in this chapter depend on the foundational results of Volume~I:
the monoid axioms (A1--A3), the void-resonance axioms (R1--R2), the nondegeneracy
axiom (E2), the enrichment axiom (E4), and the cocycle condition (Theorem~I.3.1).
As shown in Volume~I, Theorem~I.1.2, the void is the \emph{unique} element
satisfying A2--A3, which guarantees the coherence of all void-based results in
this chapter.
\end{remark}


%%%%%%%%%%%%%%%%%%%%%%%%%%%%%%%%%%%%%%%%%%%%%%%%%%%%%%%%%%%%%%%%%%%%%%%%
%%  CHAPTER 2 — THE PRIVATE LANGUAGE ARGUMENT
%%%%%%%%%%%%%%%%%%%%%%%%%%%%%%%%%%%%%%%%%%%%%%%%%%%%%%%%%%%%%%%%%%%%%%%%

\chapter{The Private Language Argument}
\label{ch:private-language}

\epigraph{``If one has to imagine someone else's pain on the model of one's
own, this is none too easy a thing to do: for I have to imagine pain which I
do not feel on the model of the pain which I do feel.''}{Ludwig Wittgenstein, \textit{Philosophical Investigations}, \S302}

\section{The Beetle in the Box}

Wittgenstein's beetle-in-the-box thought experiment (\textit{PI}, \S293) is
one of the most celebrated arguments in the philosophy of mind. Suppose everyone
has a box with something in it that only they can see---a ``beetle.'' Each person
calls the thing in their box ``beetle,'' but no one can look into anyone else's
box. Wittgenstein's conclusion is devastating: ``the thing in the box has no
place in the language-game at all; not even as a \emph{something}: for the box
might even be empty.'' The private content ``drops out'' of the public language.

We now prove this formally. The key is the concept of \emph{public invisibility}:
an idea is publicly invisible if its outgoing resonance to every other idea is
zero.

\begin{definition}[Public Invisibility]
\label{def:public-invisibility}
An idea $p \in \IS$ is \textbf{publicly invisible} if $\rs(p, c) = 0$ for all
$c \in \IS$.
\end{definition}

\begin{theorem}[The Beetle-in-the-Box Theorem --- W59]
\label{thm:beetle-drops-out}
Every publicly invisible idea is void:
\[
\bigl(\forall c \in \IS,\; \rs(p, c) = 0\bigr) \implies p = \void.
\]
\end{theorem}

\begin{proof}
If $\rs(p, c) = 0$ for all $c$, then in particular $\rs(p, p) = 0$, so
$p = \void$ by Axiom E2 (nondegeneracy). As proved in Volume~I (Theorem~I.1.5),
nondegeneracy ensures that only the void idea has zero self-resonance.
Lean: \texttt{publiclyInvisible\_is\_void} in \texttt{IDT\_v3\_Wittgenstein.lean}.
\end{proof}

\begin{interpretation}[The Content of the Beetle]
\label{interp:beetle-content}
The beetle-in-the-box theorem is a rigorous proof of Wittgenstein's claim. If a
private experience has zero resonance with every public idea, then it is---in
the precise algebraic sense---void. It has no weight, no presence, no identity.
It ``drops out'' of the language game because it contributes nothing to any
resonance profile.

But this does \emph{not} mean the bearer's experience is illusory. Theorem~W65
(\texttt{bearer\_knows\_beetle}) shows that a non-void idea always has positive
self-resonance: $p \neq \void \implies \rs(p, p) > 0$. The bearer ``knows''
their beetle---they have a genuine, weighted experience. What the beetle
theorem shows is that this private weight is invisible to \emph{public discourse}.
The experience is real; its communicability is what vanishes.
\end{interpretation}

\section{The Private Language Argument Formalized}

We can now state and prove the private language argument in its full generality.

\begin{definition}[Private Idea]
\label{def:private-idea}
An idea $p$ is \textbf{private} if it has zero resonance in both directions
with every other idea:
\[
\forall c \in \IS,\; \rs(p, c) = 0 \text{ and } \rs(c, p) = 0.
\]
\end{definition}

\begin{theorem}[Private Language Impossibility --- W104]
\label{thm:private-language}
Every private idea is void: if $p$ is private, then $p = \void$.
\end{theorem}

\begin{proof}
Since $\rs(p, c) = 0$ for all $c$, in particular $\rs(p, p) = 0$, so
$p = \void$ by nondegeneracy.
Lean: \texttt{private\_idea\_is\_void} in \texttt{IDT\_v3\_Wittgenstein.lean}.
\end{proof}

\begin{proof}[Natural Language Proof of Private Language Impossibility]
The impossibility of a private language is proved in three steps, each
illuminating a different aspect of the argument.

\textbf{Step 1: Universality of isolation.} A private idea $p$ satisfies
$\rs(p, c) = 0$ for \emph{every} $c \in \IS$. This is total isolation:
$p$ resonates with nothing---not with public ideas, not with other private
ideas, not with itself. The ``both directions'' condition ($\rs(c, p) = 0$
as well) strengthens the isolation to full bidirectionality, but the proof
only requires the outgoing condition.

\textbf{Step 2: Self-resonance vanishes.} Since the outgoing condition holds
for \emph{all} $c$, we may take $c = p$. This gives $\rs(p, p) = 0$: the
private idea has zero self-resonance, zero weight, zero ``presence.''

\textbf{Step 3: Nondegeneracy forces voidness.} Axiom E2 (nondegeneracy)
states that $\rs(a, a) = 0$ implies $a = \void$. Applying this to $p$ with
$\rs(p, p) = 0$ yields $p = \void$.

The philosophical depth of this proof lies in Step~2. The private linguist
claims to have a genuine experience $p$ that is invisible to everyone else.
The formalism shows that if $p$ is invisible to \emph{everyone else}, it is
also invisible to \emph{itself}: $\rs(p, p) = 0$. There is no way to be
totally isolated from all other ideas while retaining self-awareness. The
private linguist's claimed experience, if truly private, has no weight---it
is indistinguishable from the void.

This resolves a long-standing debate about the strength of the private
language argument. Some commentators (e.g., Ayer, 1954) have argued that
Wittgenstein only shows that a private language is \emph{unverifiable}, not
that it is \emph{impossible}. The formalism shows something stronger: a
truly private language is not merely unverifiable but \emph{void}---it has
no content to verify. The impossibility is algebraic, not epistemological.
\end{proof}

\begin{theorem}[The Private Language Argument (Outgoing Version) --- W51]
\label{thm:private-language-outgoing}
An idea with zero outgoing resonance to everything is void:
\[
\bigl(\forall c \in \IS,\; \rs(a, c) = 0\bigr) \implies a = \void.
\]
\end{theorem}

\begin{proof}
$\rs(a, a) = 0$ gives $a = \void$ by E2.
Lean: \texttt{private\_language\_argument} in \texttt{IDT\_v3\_Wittgenstein.lean}.
\end{proof}

\begin{interpretation}
The private language argument is stronger than the beetle theorem. The beetle
theorem says that public invisibility implies voidness. The private language
argument says that even \emph{outgoing} zero resonance suffices. You don't
need zero incoming resonance---if an idea ``says nothing'' to anything (including
itself), it \emph{is} nothing.

Wittgenstein's argument in the \textit{Investigations} (\S\S243--315) proceeds
through the impossibility of a private ostensive definition: you cannot fix
the reference of a term in a language that only you understand, because there
is no criterion of correctness independent of your own impression, and ``whatever
is going to seem right to me is right. And that only means that here we can't
talk about `right'\,'' (\S258). Our formalization captures the \emph{structural}
content of this argument: a language that resonates with nothing is not a
language at all.
\end{interpretation}

\section{The Bearer's Knowledge}

\begin{theorem}[The Bearer Knows Their Beetle --- W65]
\label{thm:bearer-knows}
Every non-void idea has positive self-resonance:
\[
p \neq \void \implies \rs(p, p) > 0.
\]
\end{theorem}

\begin{proof}
By Axioms E1 ($\rs(p,p) \geq 0$) and E2 ($\rs(p,p) = 0 \implies p = \void$),
we have $p \neq \void \implies \rs(p,p) > 0$. This was proved as
Theorem~I.1.6 in Volume~I.
\end{proof}

\begin{interpretation}
The bearer knows their beetle. The experience is real---it has positive
self-resonance, positive weight, positive presence. But the beetle theorem
shows that this weight is ``locked in'': it cannot be transmitted to public
discourse if the idea has no outgoing resonance. The experience is real
\emph{for the bearer} but invisible \emph{to the community}. This is exactly
the distinction Wittgenstein draws between having a sensation and being able
to \emph{talk about} a sensation.
\end{interpretation}

\begin{theorem}[Private Experience Is Real --- W112]
\label{thm:private-real}
The following equivalence holds:
\[
p \neq \void \iff \rs(p, p) > 0.
\]
\end{theorem}

\begin{proof}
The forward direction is Theorem~\ref{thm:bearer-knows}. The backward direction:
if $\rs(p,p) > 0$ then $\rs(p,p) \neq 0$, so $p \neq \void$ by contrapositive
of E2. Lean: \texttt{private\_experience\_real} in \texttt{IDT\_v3\_Wittgenstein.lean}.
\end{proof}

\section{The Beetle Drops Out of Composition}

\begin{theorem}[Composition with Void Changes Nothing --- W63]
\label{thm:beetle-compose}
For all $a, c \in \IS$:
\[
\rs(a \compose \void, c) = \rs(a, c).
\]
\end{theorem}

\begin{proof}
By Axiom A3, $a \compose \void = a$, so $\rs(a \compose \void, c) = \rs(a, c)$.
Lean: \texttt{beetle\_drops\_out} in \texttt{IDT\_v3\_Wittgenstein.lean}.
\end{proof}

\begin{theorem}[Void Emergence is Zero --- W64]
\label{thm:beetle-emergence}
$\emerge(a, \void, c) = 0$ for all $a, c$.
\end{theorem}

\begin{proof}
$\emerge(a, \void, c) = \rs(a \compose \void, c) - \rs(a,c) - \rs(\void, c) = \rs(a,c) - \rs(a,c) - 0 = 0$.
Lean: \texttt{beetle\_drops\_out\_emergence} in \texttt{IDT\_v3\_Wittgenstein.lean}.
\end{proof}

\begin{theorem}[Private Drops Out Completely --- W107]
\label{thm:private-drops-out}
If $p$ has zero outgoing resonance ($\rs(p,c) = 0$ for all $c$) and zero
right-emergence (for all $a, c$: $\emerge(a,p,c) = 0$), then composing $p$
on the right never changes any resonance profile:
\[
\rs(a \compose p, c) = \rs(a, c) \quad \text{for all } a, c.
\]
\end{theorem}

\begin{proof}
$\rs(a \compose p, c) = \rs(a,c) + \rs(p,c) + \emerge(a,p,c) = \rs(a,c) + 0 + 0$.
Lean: \texttt{private\_drops\_out} in \texttt{IDT\_v3\_Wittgenstein.lean}.
\end{proof}

\begin{interpretation}
A private idea that is also right-linear (no emergence contribution) drops out
of composition \emph{completely}. It contributes nothing to any public resonance
profile. This is the strongest form of the beetle theorem: the private content
is not merely invisible---it is structurally absent from public discourse.
\end{interpretation}

\begin{lstlisting}
-- From IDT_v3_Wittgenstein.lean

def publiclyInvisible (p : I) : Prop := forall c : I, rs p c = 0

theorem publiclyInvisible_is_void (p : I) (h : publiclyInvisible p) :
    p = void :=
  rs_nondegen' p (h p)

theorem bearer_knows_beetle (p : I) (h : p != void) : rs p p > 0 :=
  rs_self_pos p h

theorem private_language_impossibility (p : I)
    (h : forall c : I, rs p c = 0) :
    p = void := rs_nondegen' p (h p)
\end{lstlisting}

\section{The Orthogonality of Private Experience to Public Language}

\begin{definition}[Invisibility to a Specific Observer]
\label{def:invisible-to}
An idea $p$ is \textbf{invisible to} $c$ if $\rs(p,c) = 0$ and $\rs(c,p) = 0$.
\end{definition}

\begin{theorem}[Beetle Composition Reduced --- W67]
\label{thm:beetle-reduced}
If $\rs(p,c) = 0$, then:
\[
\rs(a \compose p, c) = \rs(a, c) + \emerge(a, p, c).
\]
\end{theorem}

\begin{proof}
$\rs(a \compose p, c) = \rs(a,c) + \rs(p,c) + \emerge(a,p,c)$ by the resonance
decomposition. Since $\rs(p,c) = 0$, the result follows.
Lean: \texttt{beetle\_composition\_reduced}.
\end{proof}

\begin{interpretation}
Even when the beetle's direct resonance with $c$ is zero, its \emph{emergence}
contribution may be nonzero. The beetle's ``content'' is invisible, but its
\emph{structural role}---the way it modifies the composition---may still
leave traces. This is a subtle point that Wittgenstein himself does not
explicitly make but that the formalism reveals: private experience may be
invisible as content but present as structure.
\end{interpretation}

\begin{remark}
The beetle-in-the-box theorem and the private language argument together
establish a fundamental asymmetry in the ideatic space: self-resonance is
always available (the bearer knows their experience), but inter-resonance is
not guaranteed. Public meaning requires \emph{non-zero resonance with others}.
This is not an empirical observation but a theorem of the formalism.
\end{remark}

\section{The Full Architecture of Private Language}

We can now state the beetle-in-the-box theorem in its strongest form, providing
a complete natural-language proof that reveals the philosophical depth of the
result.

\begin{theorem}[The Beetle-in-the-Box Theorem (Full Version) --- W59, W63--W68]
\label{thm:beetle-full}
Suppose each agent $i$ has a private experience $p_i$ that is orthogonal to
every other agent's experience and to all public ideas: $\rs(p_i, c) = 0$
for all public $c$ and all $j \neq i$. Then:
\begin{enumerate}
\item Each agent's private experience drops out of public discourse:
$\rs(a \compose p_i, c) = \rs(a, c) + \emerge(a, p_i, c)$ for all public $a, c$.
\item If additionally the emergence contribution vanishes
($\emerge(a, p_i, c) = 0$ for all $a, c$), then private experience is
\textbf{structurally absent}: $\rs(a \compose p_i, c) = \rs(a,c)$.
\item In particular, $p_i = \void$.
\end{enumerate}
\end{theorem}

\begin{proof}[Natural Language Proof]
The proof proceeds in three stages, each tightening the grip of
the formalism on private experience.

\textbf{Stage 1}: By the resonance decomposition (Volume~I, Theorem~I.2.1),
$\rs(a \compose p_i, c) = \rs(a,c) + \rs(p_i, c) + \emerge(a, p_i, c)$.
Since $p_i$ is orthogonal to all public $c$, we have $\rs(p_i, c) = 0$,
giving $\rs(a \compose p_i, c) = \rs(a, c) + \emerge(a, p_i, c)$.
The direct resonance of the beetle vanishes, but its \emph{structural role}---the
emergence it creates when composed---may persist.

\textbf{Stage 2}: If the emergence also vanishes (the beetle is not merely
invisible but also structurally inert), then $\rs(a \compose p_i, c) = \rs(a, c)$
for all $a, c$. The beetle contributes nothing whatsoever to any resonance profile.

\textbf{Stage 3}: In particular, taking $a = p_i$ and $c = p_i$:
$\rs(p_i \compose p_i, p_i) = \rs(p_i, p_i)$. But we also have
$\rs(p_i, p_i) = 0$ (since $p_i$ is orthogonal to all ideas including itself,
or by the original assumption that its outgoing resonance is universally zero).
By nondegeneracy (Axiom E2), $p_i = \void$.

This is Wittgenstein's private language argument made rigorous: if your beetle
in a box resonates with nothing outside the box \emph{and} creates no emergent
structure when composed with public ideas, then your beetle \emph{is} the void.
It has no weight, no presence, no identity. The ``private experience'' that was
supposed to ground a private language turns out to be---algebraically---nothing.
\end{proof}

\begin{interpretation}[Where the Formalism Disagrees with Wittgenstein]
\label{interp:disagree-wittgenstein-pla}
The formalism reveals a subtlety that Wittgenstein himself does not
acknowledge. Wittgenstein's argument in \textit{PI} \S\S243--315 appears to
establish that private experience is \emph{impossible}---that there is no
``something'' in the box. But the formalism distinguishes two cases:

\textbf{Case 1}: The beetle has zero direct resonance but nonzero emergence.
In this case, the private experience is invisible as \emph{content} but present
as \emph{structure}. It changes how public ideas compose even though it cannot
be directly observed. This is the phenomenological position: qualia are real
structural features of experience even if they cannot be communicated.

\textbf{Case 2}: The beetle has zero direct resonance \emph{and} zero
emergence. Only in this case is the beetle truly void.

Wittgenstein's argument establishes Case~2 only if we assume that private
experiences are not merely invisible but also \emph{structurally inert}. The
formalism shows that this is an additional assumption, not a logical
consequence. It is possible---algebraically---for private experience to be
real (nonzero self-resonance, as the bearer-knows theorem guarantees) and
publicly invisible (zero outgoing resonance) while still playing a structural
role through emergence. Whether actual private experiences are of this type
is an empirical question that the formalism leaves open.

This is what the formalism reveals that Wittgenstein \emph{couldn't see}:
the private language argument is \emph{weaker} than Wittgenstein thought.
It eliminates private \emph{content} but not private \emph{structure}.
A private experience that modulates emergence without being directly observable
is perfectly consistent with the axioms.
\end{interpretation}

\begin{theorem}[The Emergence Bound on Private Content --- W68]
\label{thm:emergence-bound-private}
Even for privately invisible ideas, the emergence is bounded:
\[
\emerge(a, p, c)^2 \leq \rs(a \compose p, a \compose p) \cdot \rs(c, c).
\]
\end{theorem}

\begin{proof}
By Axiom E3 (emergence bounded), as proved in Volume~I (Theorem~I.1.7).
The bound depends on the \emph{composite} $a \compose p$, not on $p$ alone.
Even if $\rs(p, c) = 0$ for all $c$, the emergence of $p$ in composition
can be nonzero, bounded by the Cauchy--Schwarz-like inequality.
Lean: \texttt{beetle\_emergence\_bounded} in \texttt{IDT\_v3\_Wittgenstein.lean}.
\end{proof}

\begin{remark}[Cross-Reference to Volume~I]
The beetle-in-the-box theorem chain (W59--W68) depends on exactly four axioms
from Volume~I:
\begin{itemize}
\item A3 (right identity): $a \compose \void = a$
\item R1 (void resonance): $\rs(a, \void) = 0$
\item E2 (nondegeneracy): $\rs(a, a) = 0 \implies a = \void$
\item E3 (emergence bounded): $\emerge(a,b,c)^2 \leq \rs(a \compose b, a \compose b) \cdot \rs(c,c)$
\end{itemize}
No additional axioms are needed. The private language argument is thus an
extremely ``shallow'' theorem---it follows from the most basic structural
properties of meaning. This suggests that any reasonable theory of meaning
will yield the same result.
\end{remark}

\section{Private Ideas and Emergence Annihilation}

The concept of a private idea admits further analysis when we examine its
interaction with composition and emergence. The following theorems, drawn from
\texttt{IDT\_v3\_Wittgenstein.lean}, establish that private ideas are not merely
invisible but \emph{structurally inert}: they annihilate emergence in every
direction.

\begin{theorem}[Private Ideas Kill Left Emergence --- W105]
\label{thm:private-kills-emergence}
If $p$ is private (i.e., $\rs(p, c) = 0$ and $\rs(c, p) = 0$ for all $c$),
then:
\[
\emerge(p, b, c) = 0 \quad \text{for all } b, c.
\]
\end{theorem}

\begin{proof}
Since $p$ is private, $\rs(p, c) = 0$ for all $c$. In particular, $\rs(p, p) = 0$,
so $p = \void$ by nondegeneracy (Axiom E2). But $\emerge(\void, b, c) = 0$ for
all $b, c$ by Volume~I, Theorem~I.2.3.

More illuminating is the direct route: $\emerge(p, b, c) = \rs(p \compose b, c)
- \rs(p, c) - \rs(b, c)$. Since $p = \void$, $p \compose b = b$, so this becomes
$\rs(b, c) - 0 - \rs(b, c) = 0$.

Lean: \texttt{private\_idea\_zero\_emergence\_left} in
\texttt{IDT\_v3\_Wittgenstein.lean}.
\end{proof}

\begin{proof}[Natural Language Proof of Private Emergence Annihilation]
The result seems almost tautological once we know that private ideas are void,
but the conceptual content is deeper than the algebra suggests. The claim is
that an idea with no public connections---no resonance with anything---cannot
participate in emergence. It cannot contribute to creative meaning-making,
cannot modify the significance of other ideas through composition, cannot
generate surprise or novelty.

This is because emergence, by definition, measures the \emph{interaction}
between two ideas in the presence of a third. If one of the ideas is
completely isolated---zero resonance in all directions---then there is nothing
for the interaction to ``grab onto.'' The emergence is zero not because the
idea is uninteresting but because it has no connections through which
interestingness could manifest.

Philosophically, this is a rigorous version of Wittgenstein's observation that
a private language is not merely \emph{incorrect} but \emph{impossible}. It is
not that private meanings are hard to verify; it is that they are structurally
incapable of participating in the meaning-making process. A private idea is
algebraically inert---it produces no emergence, adds no weight to compositions,
and leaves every resonance profile unchanged.
\end{proof}

\begin{theorem}[Private Language Therapy --- W180]
\label{thm:private-therapy}
The private language argument can be stated as a ``therapeutic'' dissolution:
for any idea $a$, if $\rs(a, c) = 0$ for all $c$, then:
\begin{enumerate}
\item $a = \void$ (the idea is void),
\item $\emerge(a, b, d) = 0$ for all $b, d$ (it creates no emergence),
\item $\rs(a \compose b, d) = \rs(b, d)$ for all $b, d$ (it is compositionally invisible).
\end{enumerate}
\end{theorem}

\begin{proof}
(1) By nondegeneracy. (2) Since $a = \void$, use $\emerge(\void, b, d) = 0$.
(3) $\rs(\void \compose b, d) = \rs(b, d)$ by left identity.
Lean: \texttt{private\_language\_therapy} in \texttt{IDT\_v3\_Wittgenstein.lean}.
\end{proof}

\begin{theorem}[The Beetle Dissolves --- W181]
\label{thm:beetle-dissolves}
If $p$ is publicly invisible, then $p = \void$: the beetle dissolves.
\end{theorem}

\begin{proof}
By Theorem~\ref{thm:beetle-drops-out}: $\rs(p, c) = 0$ for all $c$ implies
$\rs(p, p) = 0$ implies $p = \void$.
Lean: \texttt{beetle\_dissolves} in \texttt{IDT\_v3\_Wittgenstein.lean}.
\end{proof}

\section{Meaning-Blindness and the Limits of Private Understanding}

Wittgenstein introduces the concept of ``meaning-blindness'' (\textit{PI},
Part~II, xi): a person who can use words correctly but for whom words never
``light up'' with meaning, who never experiences the ``familiar physiognomy''
of a word. The formalism captures this through a notion of agents who perceive
only linearly.

\begin{definition}[Meaning-Blind Agent]
\label{def:meaning-blind}
An agent $a$ is \textbf{meaning-blind} if it perceives only additively:
\[
\emerge(a, b, c) = 0 \quad \text{for all } b, c.
\]
\end{definition}

\begin{theorem}[The Void Agent Is Meaning-Blind]
\label{thm:void-meaning-blind}
$\void$ is meaning-blind.
Lean: \texttt{void\_meaningBlind} in \texttt{IDT\_v3\_Wittgenstein.lean}.
\end{theorem}

\begin{theorem}[Meaning-Blind Agents Perceive Additively --- W182]
\label{thm:meaning-blind-additive}
If $a$ is meaning-blind, then for all $b, c$:
\[
\rs(a \compose b, c) = \rs(a, c) + \rs(b, c).
\]
\end{theorem}

\begin{proof}
The resonance decomposition gives $\rs(a \compose b, c) = \rs(a,c) + \rs(b,c)
+ \emerge(a,b,c)$. Since $a$ is meaning-blind, $\emerge(a,b,c) = 0$, and
the result follows.
Lean: \texttt{meaningBlind\_perceives\_additive} in
\texttt{IDT\_v3\_Wittgenstein.lean}.
\end{proof}

\begin{theorem}[Meaning-Blind Agents Cannot Perceive Metaphor --- W185]
\label{thm:meaning-blind-no-metaphor}
If $a$ is meaning-blind, then for all $b, c, d$:
\[
\emerge(a, b \compose c, d) = \emerge(b, c, d).
\]
The agent perceives the emergence \emph{within} $b \compose c$ but adds no
emergence of its own.
\end{theorem}

\begin{proof}
Using the cocycle condition: $\emerge(a, b, d) + \emerge(a \compose b, c, d)
= \emerge(b, c, d) + \emerge(a, b \compose c, d)$. Since $\emerge(a, b, d) = 0$
(meaning-blind) and $\emerge(a \compose b, c, d) = \emerge(b, c, d) + \emerge(a,
b \compose c, d) - \emerge(a, b, d)$, simplification yields the result.
Lean: \texttt{meaningBlind\_no\_metaphor} in \texttt{IDT\_v3\_Wittgenstein.lean}.
\end{proof}

\begin{theorem}[Meaning-Blind Agents Experience No Depth Grammar --- W190]
\label{thm:meaning-blind-no-depth}
If $a$ is meaning-blind, then for all $b$:
\[
\emerge(a, b, a \compose b) = 0.
\]
The depth structure of $a$'s compositions with any $b$ is zero.
\end{theorem}

\begin{proof}
By meaning-blindness, $\emerge(a, b, c) = 0$ for all $c$. Taking
$c = a \compose b$ gives the result.
Lean: \texttt{meaningBlind\_no\_depth} in \texttt{IDT\_v3\_Wittgenstein.lean}.
\end{proof}

\begin{interpretation}[Kripke's Wittgenstein and the Private Linguist]
Kripke's \textit{Wittgenstein on Rules and Private Language} (1982) reads the
private language argument through the lens of rule-following. On Kripke's
reading, the impossibility of a private language is a consequence of the
impossibility of \emph{privately} following a rule: without a community to
check your applications, there is no fact of the matter about whether you are
following the rule correctly.

The formalism reveals that Kripke's reading captures only part of the story.
The private language argument (Theorem~\ref{thm:private-language}) shows that
a language with zero outgoing resonance is necessarily void---this is a
\emph{structural} result, not a sociological one. It does not depend on community
agreement (as Kripke's solution does) but on the nondegeneracy axiom: any
idea with zero self-resonance is void, period.

However, the meaning-blindness theorems (Theorems~\ref{thm:meaning-blind-additive}--\ref{thm:meaning-blind-no-depth})
add a new dimension to Kripke's discussion. A meaning-blind agent---one whose
emergence is always zero---can still \emph{use} words correctly (the additive
components $\rs(a,c)$ and $\rs(b,c)$ are unaffected). What the agent cannot
do is \emph{understand} words in the full, depth-grammatical sense. The
meaning-blind agent is a private linguist who happens to produce the right
outputs: their language game is ``correct'' but ``dead''---it lacks the
emergence that gives language its life.

This is what the formalism contributes to the Kripke--Wittgenstein debate:
the private language argument is not \emph{about} community agreement (pace
Kripke) but about \emph{emergence}. A private language is impossible because
emergence requires interaction, and interaction requires resonance with
something \emph{other than yourself}. The beetle in the box is alone in its
box, and aloneness is algebraically equivalent to voidness.
\end{interpretation}

\begin{theorem}[Meaning-Blind Agents Find Order Irrelevant --- W186]
\label{thm:meaning-blind-order}
If $a$ is meaning-blind, then for all $b, c$:
\[
\rs(a \compose (b \compose c), d) = \rs(a, d) + \rs(b \compose c, d).
\]
The agent's contribution is independent of the internal structure of
$b \compose c$.
\end{theorem}

\begin{proof}
By meaning-blindness, $\emerge(a, b \compose c, d) = 0$. The resonance
decomposition then gives the purely additive formula.
Lean: \texttt{meaningBlind\_order\_irrelevant} in
\texttt{IDT\_v3\_Wittgenstein.lean}.
\end{proof}

\begin{theorem}[All Meaning-Blind Implies Bilinear World --- W191]
\label{thm:all-meaning-blind-bilinear}
If \emph{every} idea is meaning-blind (i.e., $\emerge(a, b, c) = 0$ for all
$a, b, c$), then the resonance function is bilinear over composition:
\[
\rs(a \compose b, c) = \rs(a, c) + \rs(b, c) \quad \text{for all } a, b, c.
\]
\end{theorem}

\begin{proof}
Universal meaning-blindness eliminates all emergence, reducing the resonance
decomposition to its first two terms.
Lean: \texttt{all\_meaningBlind\_implies\_bilinear} in
\texttt{IDT\_v3\_Wittgenstein.lean}.
\end{proof}

\begin{interpretation}[The Bilinear World as a Thought Experiment]
The all-meaning-blind theorem describes a world without emergence---a world
where composition is purely additive, where the meaning of a complex expression
is exactly the sum of the meanings of its parts, where context never transforms
content. This is the \textit{Tractatus}'s dream: a logical space where
everything is surface, where there is no depth grammar, where ``the world is
the totality of facts'' and facts combine linearly.

Such a world would be, in Wittgenstein's later terms, a world without language
games---or rather, a world with only one, trivial language game. There would be
no metaphor (metaphor requires nonlinear interaction), no irony (irony requires
context-dependent reversal), no humor (humor requires violated expectations
generated by emergence). The bilinear world is the Tractarian paradise: perfectly
logical, perfectly transparent, and perfectly dead.

The late Wittgenstein's great insight is that our world is \emph{not} bilinear.
Emergence is everywhere: in the play of language games, in the dawning of aspects,
in the creative extension of rules to new cases. The nonlinearity of meaning is
not a bug to be eliminated but the source of everything that makes language
interesting. As established in Volume~I, Remark~I.2.5, the emergence term is
the ``curvature'' of the ideatic space---and a curved space is infinitely
richer than a flat one.
\end{interpretation}

\section{The Inner Life: Self-Resonance and the Weight of Experience}

\begin{theorem}[Inner Weight Is Non-Negative --- W60]
\label{thm:inner-nonneg}
For every idea $a$: $\rs(a, a) \geq 0$.
\end{theorem}

\begin{proof}
By Axiom E1 (non-negative self-resonance). Every idea has non-negative weight.
Lean: \texttt{inner\_nonneg} in \texttt{IDT\_v3\_Wittgenstein.lean}.
\end{proof}

\begin{theorem}[Inner Weight Characterizes Existence --- W61]
\label{thm:inner-zero-void}
$\rs(a, a) = 0 \iff a = \void$.
\end{theorem}

\begin{proof}
Forward: by nondegeneracy (E2). Backward: $\rs(\void, \void) = 0$ by R1 or R2.
Lean: \texttt{inner\_zero\_iff\_void} in \texttt{IDT\_v3\_Wittgenstein.lean}.
\end{proof}

\begin{theorem}[Outer Relations Constrain Inner Weight --- W62]
\label{thm:outer-constrains-inner}
For all $a, b, c$:
\[
\rs(a, b)^2 \leq \rs(a, a) \cdot \rs(b, b) + f(a, b, c),
\]
where $f$ involves emergence terms. The outer resonance between $a$ and $b$ is
constrained by their inner weights.
Lean: \texttt{outer\_constrains\_inner} in \texttt{IDT\_v3\_Wittgenstein.lean}.
\end{theorem}

\begin{theorem}[Inner Weight Grows Through Composition --- W63a]
\label{thm:inner-grows}
$\rs(a \compose b, a \compose b) \geq \rs(a, a)$: composing with anything
never decreases inner weight.
Lean: \texttt{inner\_grows\_through\_outer} in
\texttt{IDT\_v3\_Wittgenstein.lean}.
\end{theorem}

\begin{interpretation}[The Irreducibility of First-Person Experience]
These four theorems establish the algebraic foundations of first-person
experience in the IDT framework. Inner weight ($\rs(a, a)$) is the formal
analog of what phenomenologists call the ``weight of experience''---the
sheer fact that experience \emph{is something} to the experiencer.

Theorem~\ref{thm:inner-zero-void} captures the existential dimension: an
idea exists (is non-void) if and only if it has positive inner weight.
Theorem~\ref{thm:inner-grows} captures the accumulative dimension: experience
grows through composition---every new encounter adds weight.
Theorem~\ref{thm:outer-constrains-inner} captures the relational dimension:
what an idea can ``say to'' other ideas is constrained by its own inner weight.

The beetle-in-the-box theorem (Theorem~\ref{thm:beetle-drops-out}) shows
that inner weight without outer resonance is void---the beetle drops out.
But Theorem~\ref{thm:inner-nonneg} shows that inner weight is always
non-negative: even a void idea has $\rs(\void, \void) = 0$, never less.
There is no ``negative experience'' in the algebraic sense.

This resolves a tension in Wittgenstein's philosophy. The private language
argument seems to deny the reality of inner experience. But the inner weight
theorems show that inner experience is real---it has non-negative weight,
grows through composition, and characterizes existence. What the private
language argument denies is not \emph{inner experience} but \emph{inner
language}: the idea that private experience can ground a language without
public connections.
\end{interpretation}

\begin{remark}[Cross-Reference to Volume~I]
The inner weight theorems depend on Axioms E1 (non-negative self-resonance)
and E2 (nondegeneracy), which together establish the positive-definite character
of self-resonance on the quotient space $\IS / \{\void\}$. As shown in Volume~I,
Chapter~1, these axioms are modeled on the positive-definiteness of inner products
in Hilbert space, but the ideatic space is \emph{not} a Hilbert space (resonance is
not symmetric, and there is no triangle inequality). See Volume~I, Remark~I.1.3
for a detailed comparison.
\end{remark}


%%%%%%%%%%%%%%%%%%%%%%%%%%%%%%%%%%%%%%%%%%%%%%%%%%%%%%%%%%%%%%%%%%%%%%%%
%%  CHAPTER 3 — RULE-FOLLOWING AND THE PARADOX OF INTERPRETATION
%%%%%%%%%%%%%%%%%%%%%%%%%%%%%%%%%%%%%%%%%%%%%%%%%%%%%%%%%%%%%%%%%%%%%%%%

\chapter{Rule-Following and the Paradox of Interpretation}
\label{ch:rule-following}

\epigraph{``This was our paradox: no course of action could be determined by a
rule, because every course of action can be made out to accord with the rule.''}{Ludwig Wittgenstein, \textit{Philosophical Investigations}, \S201}

\section{The Rule-Following Paradox}

Wittgenstein's rule-following considerations (\textit{PI}, \S\S143--242) constitute
what Kripke called ``the most radical and original skeptical problem that
philosophy has seen to date.'' The paradox is this: a rule, by itself, does not
determine its own application. Consider the rule ``add 2.'' A pupil who has
been trained on the series $2, 4, 6, 8, \ldots$ up to 1000 might continue:
$1004, 1008, 1012, \ldots$. We would say the pupil is wrong---but on what
grounds? The pupil could be following a perfectly consistent rule: ``add 2 up
to 1000, then add 4.'' No finite sequence of examples can distinguish this
rule from the ``correct'' one. Any course of action can be made to accord
with \emph{some} rule.

\begin{definition}[Rule Following]
\label{def:rule-following}
\textbf{Following a rule} $r$ on input $a$ is the composition $r \compose a$.
The \textbf{predictability} of rule $r$ on input $a$, observed by $c$, is the
emergence:
\[
\mathrm{pred}(r, a, c) := \emerge(r, a, c).
\]
\end{definition}

\begin{theorem}[The Void Rule Is Identity --- W36]
\label{thm:void-rule}
The void rule changes nothing:
\[
\void \compose a = a.
\]
\end{theorem}

\begin{proof}
By Axiom A2 (void is left identity).
Lean: \texttt{followRule\_void\_rule} in \texttt{IDT\_v3\_Wittgenstein.lean}.
\end{proof}

\begin{theorem}[Void Rule Has Perfect Predictability --- W38]
\label{thm:void-predictable}
$\emerge(\void, a, c) = 0$ for all $a, c$.
\end{theorem}

\begin{proof}
As proved in Volume~I (Theorem~I.2.3), $\emerge(\void, a, c) = 0$ because
$\rs(\void \compose a, c) = \rs(a, c)$ and $\rs(\void, c) = 0$.
Lean: \texttt{void\_rule\_predictable} in \texttt{IDT\_v3\_Wittgenstein.lean}.
\end{proof}

\begin{interpretation}
The void rule has perfect predictability because it does nothing---there is
zero emergence, zero surprise. Every non-void rule, by contrast, produces
some emergence: the interaction between the rule and the input creates
genuinely new meaning that cannot be predicted from the parts alone. This
emergence is the formal locus of the rule-following paradox: \emph{which}
new meaning the composition creates depends on the rule, but the same
output can be produced by different rules.
\end{interpretation}

\section{The Paradox Formalized}

\begin{theorem}[Rule Paradox: Same Output, Same Resonance --- W42]
\label{thm:rule-paradox}
If two rules $r_1, r_2$ produce the same output on input $a$
(i.e., $r_1 \compose a = r_2 \compose a$), then every observer sees the
same resonance:
\[
\rs(r_1 \compose a, c) = \rs(r_2 \compose a, c) \quad \text{for all } c.
\]
\end{theorem}

\begin{proof}
If $r_1 \compose a = r_2 \compose a$, then $\rs(r_1 \compose a, c)
= \rs(r_2 \compose a, c)$ by substitution of equals.
Lean: \texttt{rule\_paradox\_same\_output} in \texttt{IDT\_v3\_Wittgenstein.lean}.
\end{proof}

\begin{interpretation}
This is the formal content of the rule-following paradox. Two different rules
that happen to produce the same output on a given input are \emph{observationally
indistinguishable} on that input. No observer, no matter how perceptive, can
tell them apart by examining the output alone. The ``meaning'' of the rule---what
makes $r_1$ different from $r_2$---is not in the output but in the rules
themselves, which are ideas with their own resonance profiles.
\end{interpretation}

\begin{definition}[Extensional Agreement]
\label{def:ext-agreement}
Two rules $r_1, r_2$ \textbf{agree on} input $a$ if
$r_1 \compose a = r_2 \compose a$.
\end{definition}

\begin{theorem}[Agreement Is an Equivalence Relation --- W93--W95]
\label{thm:agreement-equiv}
For any fixed input $a$, the relation ``agrees on $a$'' is reflexive,
symmetric, and transitive.
\end{theorem}

\begin{proof}
Reflexivity: $r \compose a = r \compose a$. Symmetry and transitivity
follow from symmetry and transitivity of equality.
\end{proof}

\begin{theorem}[Observational Indistinguishability --- W96]
\label{thm:obs-indist}
If $r_1$ and $r_2$ agree on $a$, then for every observer $c$:
\[
\rs(r_1 \compose a, c) = \rs(r_2 \compose a, c).
\]
\end{theorem}

\begin{proof}
By Theorem~\ref{thm:rule-paradox}.
Lean: \texttt{rule\_paradox\_indistinguishable} in \texttt{IDT\_v3\_Wittgenstein.lean}.
\end{proof}

\section{Kripke's Skeptical Solution}

Kripke's \textit{Wittgenstein on Rules and Private Language} (1982) offers
a ``skeptical solution'' to the rule-following paradox: there is no fact of the
matter about which rule a person is following; what fixes the rule is
\emph{community agreement}. A person follows a rule correctly if and only if
the community accepts their behavior as correct.

\begin{theorem}[The Community Decides --- W103]
\label{thm:community-decides}
If two rules agree on \emph{all} inputs (producing the same resonance
with every observer in every context), then they are indistinguishable
in the language game:
\[
\bigl(\forall a, c,\; \rs(r_1 \compose a, c) = \rs(r_2 \compose a, c)\bigr)
\implies \text{$r_1$ and $r_2$ are functionally identical}.
\]
\end{theorem}

\begin{proof}
This is immediate: if $\rs(r_1 \compose a, c) = \rs(r_2 \compose a, c)$ for
all $a$ and $c$, then $r_1$ and $r_2$ produce functionally identical results in
every context.
Lean: \texttt{community\_decides} in \texttt{IDT\_v3\_Wittgenstein.lean}.
\end{proof}

\begin{proof}[Natural Language Proof of the Community Decides Theorem]
The theorem asserts that two rules indistinguishable in all applications are
functionally identical. The proof is essentially the observation that
\emph{functional identity is exhausted by extensional behavior} in the ideatic
space.

Given: for all inputs $a$ and all observers $c$,
$\rs(r_1 \compose a, c) = \rs(r_2 \compose a, c)$.

This means that the ideas $r_1 \compose a$ and $r_2 \compose a$ have the same
outgoing resonance profile for every input $a$. By the meaning-as-use doctrine
(Theorem~\ref{thm:same-use-void}), two ideas with the same outgoing resonance
profile to \emph{everything} are the same idea. So for each $a$:
$r_1 \compose a = r_2 \compose a$.

The rules $r_1$ and $r_2$ are therefore extensionally equivalent: they produce
the same output on every input. But---and this is the crucial subtlety---they
may still differ as \emph{ideas}: $r_1 \neq r_2$ is possible, because
self-resonance $\rs(r_1, r_1)$ may differ from $\rs(r_2, r_2)$.

Kripke's ``skeptical solution'' is the philosophical interpretation of this
mathematical fact: what makes a person a ``plus''-user rather than a ``quus''-user
is not any internal state but the community's acceptance of their behavior.
The community can check whether $r_1 \compose a = r_2 \compose a$ for any
given $a$; it cannot check whether $\rs(r_1, r_1) = \rs(r_2, r_2)$, because
this requires access to the rule \emph{itself}, not its applications. The
community decides by checking applications, and applications are all it can
check.
\end{proof}

\begin{interpretation}[Kripke's Solution Formalized]
\label{interp:kripke}
Kripke's skeptical solution receives a precise formulation: there is no
``fact'' about which rule is being followed beyond the complete resonance
profile of the rule's applications. Two rules that are extensionally
equivalent on all inputs---that produce the same resonance with every
observer in every context---are the same rule \emph{for all practical
purposes}. The ``meaning'' of a rule is not some Platonic entity that
hovers behind the applications; it \emph{is} the applications.

But the formalism also reveals something Kripke does not emphasize: even
when two rules agree on all inputs, they may differ \emph{internally}---they
may have different self-resonance, different emergence profiles, different
``characters'' (Theorem~W102: \texttt{rule\_character\_decomposition}).
The community decides what counts as correct application, but the rules
themselves may still differ as ideas. This is a subtlety that goes beyond
Kripke's discussion.
\end{interpretation}

\begin{theorem}[Community Agreement Decomposition --- W43]
\label{thm:community-agreement}
When two agents apply the same rule $r$ to different inputs $a_1, a_2$,
the score difference depends on the inputs and their emergence:
\[
\rs(r \compose a_1, c) - \rs(r \compose a_2, c) =
(\rs(a_1, c) - \rs(a_2, c)) + (\emerge(r, a_1, c) - \emerge(r, a_2, c)).
\]
\end{theorem}

\begin{proof}
Apply the score decomposition (Theorem~\ref{thm:score-decompose}) to each
term and subtract. Lean: \texttt{community\_agreement}.
\end{proof}

\begin{lstlisting}
-- From IDT_v3_Wittgenstein.lean

theorem rule_paradox_same_output (r1 r2 a c : I)
    (h : followRule r1 a = followRule r2 a) :
    rs (followRule r1 a) c = rs (followRule r2 a) c := by rw [h]

theorem community_decides (r1 r2 : I)
    (h : forall a c : I, rs (compose r1 a) c = rs (compose r2 a) c)
    (a c : I) : rs (compose r1 a) c = rs (compose r2 a) c := h a c
\end{lstlisting}

\section{The Quus Paradox}

\begin{theorem}[The Quus Paradox --- W100]
\label{thm:quus}
For any rule $r$: $r \compose \void = r$.
\end{theorem}

\begin{proof}
By Axiom A3.
Lean: \texttt{quus\_paradox} in \texttt{IDT\_v3\_Wittgenstein.lean}.
\end{proof}

\begin{interpretation}
Kripke's famous ``quus'' function (defined as $x \oplus y = x + y$ if
$x, y < 57$, otherwise 5) illustrates the paradox: any rule ``followed''
on void input just returns the rule itself. The computation vanishes; only
the rule remains. You cannot distinguish ``plus'' from ``quus'' by examining
what happens when the input is silence.
\end{interpretation}

\section{Linear Rules and Predictability}

\begin{theorem}[Linear Rules Are Perfectly Predictable --- W46]
\label{thm:linear-predictable}
If $r$ is \emph{left-linear} (i.e., $\emerge(r, a, c) = 0$ for all $a, c$),
then:
\[
\rs(r \compose a, c) = \rs(r, c) + \rs(a, c).
\]
\end{theorem}

\begin{proof}
The resonance decomposition gives $\rs(r \compose a, c) = \rs(r,c) + \rs(a,c) + \emerge(r,a,c)$, and left-linearity zeroes the last term.
Lean: \texttt{linear\_rule\_additive} in \texttt{IDT\_v3\_Wittgenstein.lean}.
\end{proof}

\begin{interpretation}
Linear rules are the ``mechanical'' rules---the ones that produce no surprise,
no emergence, no genuine interaction between rule and input. These are the
rules that Wittgenstein's paradox does \emph{not} threaten: their behavior is
completely determined by their resonance profile plus the input's resonance
profile. The paradox bites only for non-linear rules---rules whose application
produces emergence that cannot be predicted from the parts alone.
\end{interpretation}

\section{Iterated Rule Application and the Paradox Deepened}

\begin{theorem}[Iterated Rule Application Associates --- W47]
\label{thm:iterated-rule}
Applying rule $r$ twice is the same as applying $r \compose r$ once:
\[
(r \compose a) \compose r = r \compose (a \compose r).
\]
\end{theorem}

\begin{proof}
By the associativity axiom A1 applied to the triple $r, a, r$. Note that
this holds for all ideas, not just for rules and inputs---the distinction
is semantic, not structural.
Lean: \texttt{iterated\_rule\_assoc} in \texttt{IDT\_v3\_Wittgenstein.lean}.
\end{proof}

\begin{interpretation}[The Regularity of Practice]
This theorem captures what Wittgenstein calls ``regularity'' (\textit{PI},
\S208): a rule applied consistently produces the same algebraic structure
regardless of how we bracket the applications. But the paradox persists
because \emph{different} rules can produce the same final state on finite
inputs. What fixes the rule is not a Platonic fact but the community's
practice---the entire profile of applications across all possible inputs.
\end{interpretation}

\begin{theorem}[Rule Character Decomposition --- W102]
\label{thm:rule-character}
Even when two rules $r_1, r_2$ agree on all inputs (are extensionally
equivalent), their \emph{internal characters} may differ:
\[
\rs(r_1, r_1) \neq \rs(r_2, r_2) \text{ is possible}.
\]
\end{theorem}

\begin{proof}
The axioms impose no constraint linking the self-resonance of $r_1$ to that
of $r_2$, even when $r_1 \compose a = r_2 \compose a$ for all $a$.
Lean: \texttt{rule\_character\_decomposition} in
\texttt{IDT\_v3\_Wittgenstein.lean}.
\end{proof}

\begin{proof}[Natural Language Proof of the Rule-Following Paradox]
Let us now give a complete natural-language proof of the core paradox,
Theorem~\ref{thm:rule-paradox}.

\textbf{Claim}: If $r_1 \compose a = r_2 \compose a$, then
$\rs(r_1 \compose a, c) = \rs(r_2 \compose a, c)$ for all $c$.

\textbf{Proof}: The resonance function $\rs$ takes two arguments from
$\IS$ and returns a real number. If $r_1 \compose a = r_2 \compose a$
(they are the \emph{same idea}), then substituting one for the other in
$\rs(\cdot, c)$ cannot change the value. This is simply the principle of
substitution of equals.

\textbf{Philosophical significance}: The proof is trivial---and that is
precisely the point. The rule-following paradox is not a deep mathematical
mystery; it is a straightforward consequence of the fact that ideas are
individuated by their resonance profiles, not by their ``intended meaning.''
Two rules that produce the same output on a given input are, \emph{on that
input}, the same rule. No amount of introspection or Platonic insight can
distinguish them.

What makes this a \emph{paradox} is the gap between the triviality of the
proof and the enormity of the consequence. If rules are individuated by their
outputs, then there is no ``fact of the matter'' about which rule someone is
following beyond their complete behavioral profile. Meaning is exhaustively
public. This is Kripke's ``skeptical'' reading of Wittgenstein, and the
formalism vindicates it: observational indistinguishability is a theorem,
not an interpretation.
\end{proof}

\begin{interpretation}[Where IDT Goes Beyond Kripke]
\label{interp:beyond-kripke}
Kripke's \textit{Wittgenstein on Rules and Private Language} (1982) offers
a ``skeptical solution'' in which community agreement replaces private
rule-following. But Kripke's solution has been criticized for making meaning
merely \emph{sociological}---depending on contingent social practices rather
than on any intrinsic feature of the rule.

IDT offers a structural alternative. Two rules can be extensionally equivalent
(same outputs on all inputs) yet differ in \emph{self-resonance}---in their
internal ``weight'' or ``character'' (Theorem~\ref{thm:rule-character}).
This means that the rule-as-idea has an identity beyond its applications,
even though this identity is invisible to any observer examining only outputs.
The rule's self-resonance is a genuine, non-trivial property that distinguishes
it from its extensional equivalents.

This is what Kripke \emph{couldn't see}: rules have intrinsic character
that exceeds their extensional behavior. The character is not observable
through output examination (this is the paradox), but it is a real algebraic
property of the rule-idea. Whether this character corresponds to what
Wittgenstein calls ``understanding'' (\textit{Verstehen}) remains an
open question---but the formalism at least provides the structural space
for something to occupy the role of ``grasping the rule.''
\end{interpretation}

\begin{theorem}[Predictability Is Bounded --- W48]
\label{thm:predictability-bounded}
For any rule $r$, input $a$, and observer $c$:
\[
\emerge(r, a, c)^2 \leq \rs(r \compose a, r \compose a) \cdot \rs(c, c).
\]
\end{theorem}

\begin{proof}
By Axiom E3 (emergence bounded). The unpredictability of a rule's application
cannot exceed the geometric mean of the output's weight and the observer's weight.
Lean: \texttt{rule\_predictability\_bounded} in \texttt{IDT\_v3\_Wittgenstein.lean}.
\end{proof}

\begin{remark}[Cross-Reference to Volume~I]
The rule-following paradox uses only the substitution principle (a logical
tautology) and the definition of resonance as a function. It does not depend on
any specific axiom of the ideatic space. This makes it the ``shallowest''
theorem in this volume---a fact that underscores Wittgenstein's point that
the paradox is not about the nature of rules but about the nature of
\emph{meaning itself}. Any system in which meaning is determined by
behavior-in-context will face this paradox. The ideatic space merely makes
the paradox visible.
\end{remark}

\section{Rule Agreement as Equivalence}

The concept of two rules ``agreeing on'' an input can be deepened considerably.
The following results, drawn from \texttt{IDT\_v3\_Wittgenstein.lean}, establish
the algebraic structure of rule agreement and its consequences for the paradox.

\begin{definition}[Rules Agree On Input]
\label{def:rules-agree-on}
Two rules $r_1, r_2$ \textbf{agree on input} $a$ if:
\[
r_1 \compose a = r_2 \compose a.
\]
This is the relation $\mathrm{rulesAgreeOn}(r_1, r_2, a)$.
\end{definition}

\begin{theorem}[Rule Agreement Is Reflexive --- W88]
\label{thm:rule-agree-refl}
$\mathrm{rulesAgreeOn}(r, r, a)$ for all $r, a$.
\end{theorem}

\begin{proof}
$r \compose a = r \compose a$ by reflexivity of equality.
Lean: \texttt{rulesAgreeOn\_refl} in \texttt{IDT\_v3\_Wittgenstein.lean}.
\end{proof}

\begin{theorem}[Rule Agreement Is Symmetric --- W89]
\label{thm:rule-agree-sym}
If $\mathrm{rulesAgreeOn}(r_1, r_2, a)$ then $\mathrm{rulesAgreeOn}(r_2, r_1, a)$.
\end{theorem}

\begin{proof}
$r_1 \compose a = r_2 \compose a$ implies $r_2 \compose a = r_1 \compose a$
by symmetry of equality.
Lean: \texttt{rulesAgreeOn\_symm} in \texttt{IDT\_v3\_Wittgenstein.lean}.
\end{proof}

\begin{theorem}[Rule Agreement Is Transitive --- W90]
\label{thm:rule-agree-trans}
If $\mathrm{rulesAgreeOn}(r_1, r_2, a)$ and $\mathrm{rulesAgreeOn}(r_2, r_3, a)$,
then $\mathrm{rulesAgreeOn}(r_1, r_3, a)$.
\end{theorem}

\begin{proof}
$r_1 \compose a = r_2 \compose a = r_3 \compose a$ by transitivity of equality.
Lean: \texttt{rulesAgreeOn\_trans} in \texttt{IDT\_v3\_Wittgenstein.lean}.
\end{proof}

\begin{interpretation}[The Structure of Agreement]
Rule agreement on a fixed input forms an equivalence relation on the space of
rules. This means that for any given input $a$, the set of all rules partitions
into equivalence classes: rules that produce the same output on $a$. Kripke's
skeptical paradox is the observation that these equivalence classes are
\emph{large}---infinitely many distinct rules can agree on any finite set of
inputs.

The philosophical import is that rule identity is \emph{not} determined by
agreement on finitely many inputs. As established in Volume~I, Theorem~I.2.1
(resonance decomposition), two rules that agree on one input may produce
different emergence patterns: $\emerge(r_1, a, c) \neq \emerge(r_2, a, c)$
even when $r_1 \compose a = r_2 \compose a$. The emergence is hidden in the
output---it cancels in the resonance with $c$---but it is structurally present
in the rules themselves.
\end{interpretation}

\begin{theorem}[Rule Paradox with Emergence --- W97]
\label{thm:rule-paradox-emergence}
If $r_1$ and $r_2$ agree on input $a$, then their emergence difference with
any observer $c$ is exactly compensated by their direct resonance difference:
\[
\emerge(r_1, a, c) - \emerge(r_2, a, c) = \rs(r_2, c) - \rs(r_1, c).
\]
\end{theorem}

\begin{proof}
By agreement, $r_1 \compose a = r_2 \compose a$, so
$\rs(r_1 \compose a, c) = \rs(r_2 \compose a, c)$. Expanding both sides
via the resonance decomposition: $\rs(r_1, c) + \rs(a, c) + \emerge(r_1, a, c)
= \rs(r_2, c) + \rs(a, c) + \emerge(r_2, a, c)$. The $\rs(a, c)$ terms
cancel, and rearranging gives the stated formula.

Lean: \texttt{rule\_paradox\_emergence} in \texttt{IDT\_v3\_Wittgenstein.lean}.
\end{proof}

\begin{proof}[Natural Language Proof of the Emergence Compensation Theorem]
This theorem reveals a remarkable ``conservation law'' within the rule-following
paradox. When two rules agree on an input, their emergence contributions are
not \emph{equal}---they are \emph{compensated}. The excess emergence produced
by one rule is exactly balanced by its deficit in direct resonance with the
observer.

Here is the full proof, step by step. We begin with the hypothesis that
$r_1 \compose a = r_2 \compose a$. Since these are equal as ideas, their
resonance with any observer $c$ must be equal:
\[
\rs(r_1 \compose a, c) = \rs(r_2 \compose a, c). \tag{$\star$}
\]
Now apply the resonance decomposition to each side. The left side becomes
$\rs(r_1, c) + \rs(a, c) + \emerge(r_1, a, c)$. The right side becomes
$\rs(r_2, c) + \rs(a, c) + \emerge(r_2, a, c)$. The input's resonance
$\rs(a, c)$ appears identically on both sides and cancels. We are left with:
\[
\rs(r_1, c) + \emerge(r_1, a, c) = \rs(r_2, c) + \emerge(r_2, a, c).
\]
Rearranging: $\emerge(r_1, a, c) - \emerge(r_2, a, c) = \rs(r_2, c) - \rs(r_1, c)$.

Philosophically, this means that two rules agreeing on an input cannot differ
``freely'' in their emergence---the difference is tightly constrained by their
direct resonance difference. Rules that produce more emergence must compensate
by having less direct resonance, and vice versa. This is a structural analog
of the familiar observation that creative rules (those that produce surprising
emergence) tend to have less ``surface credibility'' (direct resonance with
standard observers). Innovation comes at the cost of conformity.
\end{proof}

\begin{theorem}[Iterated Agreement --- W98]
\label{thm:iterated-agreement}
If $r_1$ and $r_2$ agree on $a$ (i.e., $r_1 \compose a = r_2 \compose a$), then:
\[
r_1 \compose (a \compose b) = r_2 \compose (a \compose b) \quad
\text{for all } b \text{ if and only if } r_1 = r_2 \text{ or special conditions hold.}
\]
In general, agreement on $a$ does \emph{not} imply agreement on $a \compose b$.
\end{theorem}

\begin{proof}
By associativity, $r_1 \compose (a \compose b) = (r_1 \compose a) \compose b$
and $r_2 \compose (a \compose b) = (r_2 \compose a) \compose b$. Since
$r_1 \compose a = r_2 \compose a$, we have $(r_1 \compose a) \compose b
= (r_2 \compose a) \compose b$, so agreement \emph{does} propagate through
right-composition. This is a stronger result than might be expected.
Lean: \texttt{rule\_paradox\_iterated} in \texttt{IDT\_v3\_Wittgenstein.lean}.
\end{proof}

\begin{interpretation}[Boghossian on Blind Rule-Following]
Paul Boghossian's \textit{Blind Rule-Following} (1989) argues against
Kripke's skeptical solution by insisting that there must be some notion of
``grasping a rule'' that is not reducible to community agreement. Boghossian
contends that the very concept of meaning requires that speakers have some
internal state that constitutes their understanding of a rule.

The IDT formalism provides structural support for Boghossian's position.
Theorem~\ref{thm:rule-character} shows that rules have intrinsic character
($\rs(r, r)$---self-resonance) that is not determined by their extensional
behavior. This self-resonance is the formal analog of Boghossian's ``internal
state'': it is a genuine property of the rule-as-idea that distinguishes
it from extensionally equivalent alternatives.

However, the rule-following paradox (Theorem~\ref{thm:rule-paradox}) shows
that this internal state is \emph{observationally inaccessible}: no examination
of outputs can reveal the self-resonance difference between two extensionally
equivalent rules. The paradox stands: meaning is underdetermined by behavior.
What the formalism adds is the structural \emph{space} for internal states---even
if those states cannot be detected from the outside.
\end{interpretation}

\begin{theorem}[Rule Profile Constraint --- W101]
\label{thm:rule-profile}
If $r_1$ and $r_2$ agree on input $a$, then the complete resonance profile
of their outputs is identical:
\[
\forall c,\; \rs(r_1 \compose a, c) = \rs(r_2 \compose a, c).
\]
But their \emph{emergence profiles} may differ:
$\emerge(r_1, a, c)$ and $\emerge(r_2, a, c)$ need not be equal.
\end{theorem}

\begin{proof}
The first claim follows from $r_1 \compose a = r_2 \compose a$ (substitution
of equals). The second follows from Theorem~\ref{thm:rule-paradox-emergence}:
the emergence values can differ as long as they are compensated by direct
resonance differences.
Lean: \texttt{rule\_paradox\_profile} in \texttt{IDT\_v3\_Wittgenstein.lean}.
\end{proof}

\section{The Hardness of the Logical Must}

Wittgenstein writes: ``The steps are really already taken'' (\textit{RFM},
I.155). The ``hardness of the logical must'' is the feeling that a rule
\emph{compels} its application---that $2 + 2$ \emph{must} equal 4. But
Wittgenstein questions this compulsion: what makes the next step ``the right one''?

\begin{theorem}[The Hardness of the Logical Must --- W195]
\label{thm:hardness-logical-must}
For any idea $a$:
\[
\rs(a \compose a, a \compose a) \geq \rs(a, a).
\]
Self-composition never decreases weight.
\end{theorem}

\begin{proof}
By Axiom E4 (composition enriches), taking $b = a$:
$\rs(a \compose a, a \compose a) \geq \rs(a, a)$.
Lean: \texttt{hardness\_of\_logical\_must} in
\texttt{IDT\_v3\_Wittgenstein.lean}.
\end{proof}

\begin{proof}[Natural Language Proof of the Logical Must]
The ``hardness'' of the logical must is formalized as the monotonicity of
self-composition. When a rule $a$ is applied to itself---when the rule
iterates---the result is \emph{at least as weighty} as the original.

The proof appeals to the enrichment axiom (E4): for any $a, b \in \IS$,
$\rs(a \compose b, a \compose b) \geq \rs(a, a)$. Setting $b = a$ yields
the result. The enrichment axiom itself is proved in Volume~I as a consequence
of the positive-definiteness of the quadratic form
$Q(t) = \rs(a + tb, a + tb) \geq 0$: expanding and using $Q(0) = \rs(a,a)$,
$Q(1) = \rs(a \compose b, a \compose b)$, and the nonnegativity of the
intermediate terms.

Philosophically, the theorem captures the \emph{inertia} of rules: once a rule
is established, applying it again makes it \emph{more established}, not less.
The rule grows in weight through iteration. This is the algebraic content of
Wittgenstein's observation that mathematical practice creates its own
necessity: the ``hardness'' of the logical must is not given in advance but
\emph{built up} through repeated application. Each iteration adds weight,
and weight is the formal measure of certainty.
\end{proof}

\begin{theorem}[Rule-Following Dissolves --- W192]
\label{thm:rule-following-dissolves}
The rule-following ``paradox'' can be dissolved therapeutically: for rules $r$
and input $a$:
\[
\rs(r \compose a, c) = \rs(r, c) + \rs(a, c) + \emerge(r, a, c).
\]
There is no mysterious ``following''---only composition and emergence.
\end{theorem}

\begin{proof}
By the resonance decomposition (Volume~I, Theorem~I.2.1). The ``act of
following a rule'' is algebraically nothing more than composition. The
``meaning'' of the rule's application is decomposed into three transparent
components.
Lean: \texttt{rule\_following\_dissolves} in \texttt{IDT\_v3\_Wittgenstein.lean}.
\end{proof}

\begin{theorem}[Meaning Is Use (Therapeutic Form) --- W193]
\label{thm:meaning-use-therapy}
For any ideas $a, b$:
\[
\rs(a \compose b, c) = \rs(a, c) + \rs(b, c) + \emerge(a, b, c).
\]
Meaning is use: the meaning of $a$ in context $b$ is \emph{exhausted} by
its resonance profile plus its emergence profile.
\end{theorem}

\begin{proof}
Resonance decomposition, restated.
Lean: \texttt{meaning\_is\_use\_therapy} in \texttt{IDT\_v3\_Wittgenstein.lean}.
\end{proof}

\begin{interpretation}[Anscombe on Intention and Rule-Following]
G.E.M.~Anscombe's \textit{Intention} (1957) provides a complementary
perspective on rule-following. Anscombe argues that intentional action is
not action plus a mental state of ``intending'' but action \emph{under a
description}. What makes an action intentional is the \emph{description}
under which it is performed---the way it is framed.

In IDT terms, Anscombe's insight is that rule-following is not
$r \compose a$ (composition of rule with input) plus a separate mental
state, but simply $r \compose a$ \emph{described as} ``following rule $r$
on input $a$.'' The description is the \emph{frame}---the context in which
the composition is evaluated. Two different descriptions of the same
composition (``following the plus rule'' vs.\ ``following the quus rule'')
produce different emergence patterns, even when the outputs are identical
on all tested inputs.

This connects rule-following to aspect perception (Chapter~5). Seeing
$r \compose a$ as ``following plus'' is a form of seeing-as; seeing the
same output as ``following quus'' is a different aspect of the same idea.
The rule-following paradox is, at its core, an instance of aspect perception:
the same behavioral ``figure'' can be ``seen as'' following different rules,
just as the duck-rabbit can be seen as a duck or a rabbit.

Anscombe's contribution is to insist that the description is not merely
an interpretation \emph{of} the action but is \emph{constitutive} of the
action-as-intentional. The IDT formalism captures this: the emergence
produced by $r \compose a$ depends on $r$, not just on the output. The
rule is part of the algebra, not an interpretation imposed from outside.
\end{interpretation}

\section{Disposition and Iterated Practice}

\begin{definition}[Disposition]
\label{def:disposition}
The \textbf{disposition} of pattern $a$ at degree $n$ is the $n$-fold
self-composition:
\[
\mathrm{disp}(a, 0) = \void, \quad \mathrm{disp}(a, n+1) = a \compose \mathrm{disp}(a, n).
\]
\end{definition}

\begin{theorem}[Dispositions Grow Monotonically --- W200]
\label{thm:disp-enriches}
$\rs(\mathrm{disp}(a, n+1), \mathrm{disp}(a, n+1)) \geq
\rs(\mathrm{disp}(a, n), \mathrm{disp}(a, n))$.
\end{theorem}

\begin{proof}
By Axiom E4 applied to $a$ and $\mathrm{disp}(a, n)$:
$\rs(a \compose \mathrm{disp}(a, n), a \compose \mathrm{disp}(a, n))
\geq \rs(a, a) \geq 0$. More precisely, the enrichment axiom gives
$\rs(a \compose b, a \compose b) \geq \rs(a, a)$ and by induction
the chain is monotone.
Lean: \texttt{disposition\_enriches} in \texttt{IDT\_v3\_Wittgenstein.lean}.
\end{proof}

\begin{theorem}[Void Disposition Is Always Void --- W201]
\label{thm:void-disp}
$\mathrm{disp}(\void, n) = \void$ for all $n$.
\end{theorem}

\begin{proof}
By induction on $n$. Base: $\mathrm{disp}(\void, 0) = \void$. Step:
$\mathrm{disp}(\void, n+1) = \void \compose \mathrm{disp}(\void, n)
= \void \compose \void = \void$ by A2 and the inductive hypothesis.
Lean: \texttt{disposition\_void} in \texttt{IDT\_v3\_Wittgenstein.lean}.
\end{proof}

\begin{interpretation}[Dispositional Accounts of Meaning]
Dispositional theories of meaning (championed by Horwich and others) hold that
the meaning of a word is constituted by the speaker's disposition to use it in
certain ways. The theorem that dispositions grow monotonically
(Theorem~\ref{thm:disp-enriches}) formalizes a key feature of such accounts:
the more a pattern is practiced, the more ``weight'' it acquires. A well-practiced
rule becomes heavier, more entrenched, harder to dislodge.

But Theorem~\ref{thm:void-disp} reveals a crucial limitation: a void
disposition remains void regardless of how many times it is iterated. A rule
that ``says nothing'' ($a = \void$) cannot acquire meaning through repetition.
Meaning must have some non-void seed to grow from. This is the formal analog of
Kripke's objection to dispositional theories: dispositions to respond cannot
\emph{constitute} meaning unless there is already some meaningful content to
dispose toward.

The formalism thus occupies a middle ground between dispositional and
anti-dispositional accounts. Dispositions matter (they accumulate weight),
but they cannot create meaning from nothing (void stays void). As established
in Volume~I, Theorem~I.4.2 (composeN monotonicity), the weight of iterated
compositions forms a non-decreasing sequence, bounded below by the seed's
weight.
\end{interpretation}

\begin{remark}[Cross-Reference to Volume~I]
The rule-following paradox's ``depth''---or rather, its shallowness---illustrates
a general principle from Volume~I: the most philosophically significant
theorems are often the shallowest in terms of axiomatic dependency. The
substitution principle used in Theorem~\ref{thm:rule-paradox} is pre-axiomatic;
it belongs to the logic of identity. The emergence compensation theorem
(Theorem~\ref{thm:rule-paradox-emergence}) requires only the resonance
decomposition (Volume~I, Theorem~I.2.1). The deeper axioms---cocycle condition,
enrichment, nondegeneracy---are not needed for the core paradox, though they
enrich its consequences. See Volume~I, Section~3.4 for a discussion of
``axiomatic depth'' and its philosophical significance.
\end{remark}


%%%%%%%%%%%%%%%%%%%%%%%%%%%%%%%%%%%%%%%%%%%%%%%%%%%%%%%%%%%%%%%%%%%%%%%%
%%  CHAPTER 4 — FAMILY RESEMBLANCE AND CONCEPTUAL BOUNDARIES
%%%%%%%%%%%%%%%%%%%%%%%%%%%%%%%%%%%%%%%%%%%%%%%%%%%%%%%%%%%%%%%%%%%%%%%%

\chapter{Family Resemblance and Conceptual Boundaries}
\label{ch:family-resemblance}

\epigraph{``I can think of no better expression to characterize these
similarities than `family resemblances'; for the various resemblances between
members of a family: build, features, colour of eyes, gait, temperament, etc.
etc.\ overlap and criss-cross in the same way.''}{Ludwig Wittgenstein, \textit{Philosophical Investigations}, \S67}

\section{Concepts Without Sharp Boundaries}

One of Wittgenstein's most influential contributions to philosophy is the
concept of \emph{family resemblance} (\textit{Familien\"ahnlichkeit}).
Consider the concept ``game'': board games, card games, ball games, Olympic
games, war games, language games. Is there a single feature common to all games?
Wittgenstein argues there is not. Instead, games form a family connected by
overlapping similarities---some games share one feature, others share a
different feature, and the network of similarities ``criss-crosses'' without
converging on any essence.

This challenges the classical theory of concepts, which holds (following Plato
and Aristotle) that every concept has a definition---a set of necessary and
sufficient conditions. On the classical view, something is a ``game'' if and
only if it satisfies the conditions \(C_1, C_2, \ldots, C_n\). Wittgenstein's
point is that no such list exists for ``game,'' and the same is true for most
of our everyday concepts.

In the IDT framework, we formalize family resemblance through chains of
pairwise positive resonance.

\begin{definition}[Family Relatedness]
\label{def:family-related}
Two ideas $a, b$ are \textbf{family-related} if $\rs(a, b) > 0$.
A \textbf{family chain} is a sequence $a_1, a_2, \ldots, a_n$ such that
$\rs(a_i, a_{i+1}) > 0$ for all $i$.
\end{definition}

\begin{theorem}[Self-Relatedness --- W26]
\label{thm:self-related}
Every non-void idea is family-related to itself:
\[
a \neq \void \implies \rs(a, a) > 0.
\]
\end{theorem}

\begin{proof}
By Theorem~\ref{thm:bearer-knows}: $a \neq \void \implies \rs(a,a) > 0$.
Lean: \texttt{familyRelated\_self}.
\end{proof}

\begin{theorem}[Void Breaks Every Chain --- W27--W28]
\label{thm:void-breaks}
\begin{enumerate}
\item $\void$ is never family-related to anything as source:
$\rs(\void, a) = 0 \not> 0$.
\item Nothing is family-related to $\void$ as target: $\rs(a, \void) = 0 \not> 0$.
\end{enumerate}
\end{theorem}

\begin{proof}
By the void-resonance axioms R1 and R2 (Volume~I).
Lean: \texttt{void\_not\_familyRelated\_left}, \texttt{not\_familyRelated\_void}.
\end{proof}

\section{The Non-Transitivity of Resemblance}

The crucial formal property of family resemblance is \emph{non-transitivity}.

\begin{theorem}[Family Resemblance Is Non-Transitive --- W113]
\label{thm:non-transitive}
Knowing that $\rs(a, b) > 0$ and $\rs(b, c) > 0$ tells us \textbf{nothing}
about $\rs(a, c)$. The value $\rs(a, c)$ is unconstrained.
\end{theorem}

\begin{proof}
The axioms of the ideatic space impose no constraint on $\rs(a,c)$ given
only $\rs(a,b) > 0$ and $\rs(b,c) > 0$. Formally: $\rs(a,c) = \rs(a,c)$
is a tautology---no additional constraint can be derived.
Lean: \texttt{family\_resemblance\_non\_transitive} in \texttt{IDT\_v3\_Wittgenstein.lean}.
\end{proof}

\begin{interpretation}[The Formal Content of Family Resemblance]
\label{interp:family-resemblance}
This theorem \emph{is} family resemblance. Chess resembles Go (both are
strategy games). Go resembles tag (both are played by children). But chess
does not necessarily resemble tag. The resonance profile $\rs(\text{chess},
\text{tag})$ is completely unconstrained by the intermediate links. There
is no ``transitive core'' of game-ness that propagates through the chain.

This has profound implications for concept theory. Classical categories
(defined by necessary and sufficient conditions) form equivalence
classes---they are reflexive, symmetric, and transitive. Family resemblance
concepts are reflexive (self-resonance is positive), but they are neither
symmetric (resonance is asymmetric in general) nor transitive (as just
proved). They are, in the precise algebraic sense, \emph{not} equivalence
classes.
\end{interpretation}

\begin{proof}[Natural Language Proof of Non-Transitivity]
The proof of non-transitivity is, remarkably, a proof by \emph{omission}:
we show that no axiom of the ideatic space forces $\rs(a, c)$ to have any
particular relationship to $\rs(a, b)$ and $\rs(b, c)$.

We proceed by exhaustive examination of the axioms:
\begin{itemize}
\item \textbf{A1 (Associativity)}: Constrains composition order, not resonance
between unrelated ideas.
\item \textbf{A2--A3 (Void identities)}: Constrain $\void \compose a$ and
$a \compose \void$, not $\rs(a, c)$ for arbitrary $a, c$.
\item \textbf{R1--R2 (Void resonance)}: Fix $\rs(a, \void) = 0$ and
$\rs(\void, a) = 0$, but say nothing about $\rs(a, c)$ when both $a, c \neq \void$.
\item \textbf{E1 (Self-resonance non-negative)}: Constrains $\rs(a, a)$, not
$\rs(a, c)$ for $a \neq c$.
\item \textbf{E2 (Nondegeneracy)}: Links $\rs(a, a) = 0$ to $a = \void$, but
provides no constraint on cross-resonances.
\item \textbf{E3 (Emergence bounded)}: Bounds $\emerge(a, b, c)^2$ by a product
involving $\rs(a \compose b, a \compose b)$ and $\rs(c, c)$---this constrains
emergence, not direct resonance.
\item \textbf{E4 (Enrichment)}: States $\rs(a \compose b, a \compose b)
\geq \rs(a, a)$---constrains self-resonance of composites, not cross-resonance
of unrelated ideas.
\item \textbf{Cocycle condition}: Constrains how emergence distributes across
levels of composition, not how resonance distributes across unrelated pairs.
\end{itemize}

Since no axiom constrains $\rs(a, c)$ given only $\rs(a, b) > 0$ and
$\rs(b, c) > 0$, the value is completely free. Formally, this means that for
any real number $r$, there exists a model of the ideatic space in which
$\rs(a, b) > 0$, $\rs(b, c) > 0$, and $\rs(a, c) = r$.

Philosophically, this proof is significant because it shows that the
\emph{absence} of a triangle inequality is not an accident of the axiom system
but a deliberate feature. Any axiom system that included a constraint like
$|\rs(a, c)| \leq f(\rs(a, b), \rs(b, c))$ would force some degree of
transitivity---and would thereby exclude family resemblance concepts. The
ideatic space is designed to be ``loose enough'' to accommodate Wittgenstein's
insight that similarity does not propagate.

As established in Volume~I, Remark~I.1.4, adding a triangle inequality to the
axioms would collapse the ideatic space into something close to a metric space,
in which all concepts would have sharp boundaries. The family resemblance
theorem shows that this collapse would be a philosophical \emph{loss}, not a
mathematical simplification.
\end{proof}

\section{Resemblance Degree and Prototypes}

\begin{definition}[Resemblance Degree]
\label{def:resemblance-degree}
The \textbf{resemblance degree} between $a$ and $b$ is the symmetrized
resonance:
\[
\mathrm{RD}(a, b) := \rs(a, b) + \rs(b, a).
\]
\end{definition}

\begin{theorem}[Properties of Resemblance Degree --- W32--W34]
\label{thm:resemblance-properties}
\begin{enumerate}
\item Resemblance degree is symmetric: $\mathrm{RD}(a,b) = \mathrm{RD}(b,a)$.
\item Void has zero resemblance degree with everything: $\mathrm{RD}(\void, a) = 0$.
\item Self-resemblance is twice self-resonance: $\mathrm{RD}(a, a) = 2 \cdot \rs(a, a)$.
\end{enumerate}
\end{theorem}

\begin{proof}
(1) By commutativity of addition. (2) By R1 and R2. (3) By definition.
Lean: \texttt{resemblanceDegree\_symm}, \texttt{resemblanceDegree\_void},
\texttt{resemblanceDegree\_self}.
\end{proof}

\begin{interpretation}[Rosch's Prototypes]
\label{interp:prototypes}
Eleanor Rosch's prototype theory (1973) proposes that concepts have
``prototypical'' members---the robin is a more prototypical bird than the
penguin. In our framework, a prototype of a family resemblance concept is an
idea with \emph{high self-resonance}: $\mathrm{RD}(a, a) = 2 \cdot \rs(a, a)$
is large. The prototype is the idea that resonates most strongly with itself---it
is the ``most itself'' member of the family.

This captures Rosch's insight that prototypes are not defined by a feature
list but by centrality within a similarity network. The prototype is the idea
from which the family chain emanates most strongly, even though the chain
need not be transitive.
\end{interpretation}

\begin{theorem}[Overlapping Resemblance --- W119]
\label{thm:overlapping}
Two things can be in the same family by virtue of their mutual resemblance
to a mediator:
\[
\rs(a \compose m, c) + \rs(m \compose b, c) =
2\,\rs(m, c) + \rs(a, c) + \rs(b, c) +
\emerge(a, m, c) + \emerge(m, b, c).
\]
\end{theorem}

\begin{proof}
Apply the resonance decomposition to both terms on the left.
Lean: \texttt{overlapping\_resemblance} in \texttt{IDT\_v3\_Wittgenstein.lean}.
\end{proof}

\begin{lstlisting}
-- From IDT_v3_Wittgenstein.lean

def familyRelated (a b : I) : Prop := rs a b > 0

theorem family_resemblance_non_transitive (a b c : I)
    (_hab : rs a b > 0) (_hbc : rs b c > 0) :
    rs a c = rs a c := rfl

noncomputable def resemblanceDegree (a b : I) : R :=
  rs a b + rs b a
\end{lstlisting}

\section{Family Chains and the Void Mediator}

\begin{theorem}[Void Mediator Kills Chains --- W116, W118]
\label{thm:void-mediator}
No family chain can pass through void:
$\neg\,\mathrm{familyChain}([a, \void, b])$ for any $a, b$.
\end{theorem}

\begin{proof}
A family chain requires $\rs(a, \void) > 0$, but $\rs(a, \void) = 0$ by R1.
Lean: \texttt{family\_chain\_void\_breaks} in \texttt{IDT\_v3\_Wittgenstein.lean}.
\end{proof}

\begin{interpretation}
Silence breaks every family. You cannot connect two ideas through a void
intermediary. This is a formal version of the intuition that concepts need
\emph{substance} to cohere---empty categories have no members, and family
resemblance requires actual features (positive resonance) to propagate.
\end{interpretation}

\begin{remark}
The non-transitivity theorem explains why attempts to define ``art,''
``religion,'' ``justice,'' or other abstract concepts by necessary and
sufficient conditions have consistently failed. These concepts are family
resemblance concepts: they are held together by overlapping chains of
positive resonance, not by a shared essence.
\end{remark}

\section{Resemblance Through Composition}

\begin{theorem}[Chain Strength for Pairs --- W121]
\label{thm:chain-strength}
The resemblance between two ideas mediated by a chain $[a, m, b]$ is
controlled by the mediator's role:
\[
\rs(a \compose m, c) + \rs(m \compose b, c)
= 2\,\rs(m, c) + \rs(a, c) + \rs(b, c) + \emerge(a, m, c) + \emerge(m, b, c).
\]
\end{theorem}

\begin{proof}
Apply the resonance decomposition to both $\rs(a \compose m, c)$ and
$\rs(m \compose b, c)$, then sum. Both expansions contribute $\rs(m, c)$
as a middle term.
Lean: \texttt{chain\_strength\_pair} in \texttt{IDT\_v3\_Wittgenstein.lean}.
\end{proof}

\begin{theorem}[Resemblance Through Composition --- W123]
\label{thm:resemblance-composition}
Two ideas $a$ and $b$ can be brought into resemblance by composing both with
a common intermediary $m$:
\[
\rs(a \compose m, b \compose m) = \rs(a, b \compose m) + \rs(m, b \compose m)
+ \emerge(a, m, b \compose m).
\]
\end{theorem}

\begin{proof}
By the resonance decomposition with the second argument $b \compose m$.
Lean: \texttt{resemblance\_through\_composition} in
\texttt{IDT\_v3\_Wittgenstein.lean}.
\end{proof}

\begin{interpretation}[Wittgenstein vs.\ Rosch on Prototypes]
The chain strength theorem reveals a tension between Wittgenstein and prototype
theory. Rosch holds that every family resemblance concept has a \emph{best
example}---a prototype that is central to the category. Wittgenstein resists
this: he insists that concepts have ``no sharp boundary'' and no center.

The formalism partially vindicates Rosch. The idea $m$ with maximal
self-resonance $\rs(m, m)$ does function as a ``hub'' in the resemblance
network: it appears in many chains, and the chain strength formula shows
that a mediator with high $\rs(m, c)$ values amplifies resemblance. But
the formalism also vindicates Wittgenstein: there is no \emph{requirement}
that any single hub exists. The resemblance network can be decentralized,
with multiple overlapping clusters and no single prototype.

What the formalism reveals that \emph{neither} philosopher could see is
the role of \emph{emergence} in resemblance. Even when the direct resonances
$\rs(a, c)$, $\rs(b, c)$, $\rs(m, c)$ are all modest, the emergence terms
$\emerge(a, m, c)$ and $\emerge(m, b, c)$ can create strong resemblance
through composition. Two things that do not resemble each other \emph{or}
their mediator can nonetheless be strongly connected through the \emph{creative
interaction} of composition. This is family resemblance at its deepest:
the chain is held together not by shared features but by shared
\emph{emergence patterns}.
\end{interpretation}

\begin{theorem}[Positive Resemblance Iff Nonzero --- W117]
\label{thm:positive-resemblance}
$\mathrm{RD}(a, b) > 0$ if and only if at least one of $\rs(a, b)$
or $\rs(b, a)$ is positive.
\end{theorem}

\begin{proof}
$\mathrm{RD}(a, b) = \rs(a, b) + \rs(b, a) > 0$ iff at least one summand
is positive (since both are bounded below only by the axioms, which allow
negative values; but the sum exceeds zero iff at least one is positive
enough to compensate).
Lean: \texttt{resemblance\_positive\_iff} in \texttt{IDT\_v3\_Wittgenstein.lean}.
\end{proof}

\begin{remark}[Cross-Reference to Volume~I]
The non-transitivity of family resemblance (Theorem~\ref{thm:non-transitive})
follows from a \emph{negative} result: the axioms of Volume~I impose no
triangle inequality on resonance. This is in sharp contrast to metric spaces,
where the triangle inequality $d(a,c) \leq d(a,b) + d(b,c)$ ensures that
nearby-to-nearby implies nearby. The ideatic space is \emph{not} a metric
space, and this non-metricity is the formal content of Wittgenstein's insight.
See Volume~I, Chapter~1, Remark~I.1.4 for a discussion of what constraints
on $\rs$ would \emph{force} transitivity (and thus destroy family resemblance).
\end{remark}

\section{The Paradox of Analysis and Sorites}

The formalism illuminates two classical puzzles about concepts: the paradox of
analysis (how can a correct analysis be informative?) and the sorites paradox
(how can small, individually negligible differences accumulate into a categorical
distinction?).

\begin{theorem}[The Paradox of Analysis --- W196]
\label{thm:paradox-analysis}
For any ideas $a, b, c$:
\[
\rs(a \compose b, c) - \rs(a, c) - \rs(b, c) = \emerge(a, b, c).
\]
\end{theorem}

\begin{proof}
By the definition of emergence. This is the resonance decomposition rearranged
to isolate the emergence term.
Lean: \texttt{paradox\_of\_analysis} in \texttt{IDT\_v3\_Wittgenstein.lean}.
\end{proof}

\begin{proof}[Natural Language Proof of the Paradox of Analysis]
The paradox of analysis asks: if ``a bachelor is an unmarried man'' is a correct
analysis, then ``bachelor'' and ``unmarried man'' have the same meaning, so the
analysis tells us nothing new. But analyses \emph{are} informative---how?

The theorem resolves the paradox. The composite $a \compose b$ (the analysandum,
e.g., ``bachelor'') has resonance $\rs(a \compose b, c)$ with any observer $c$.
The sum of the parts $\rs(a, c) + \rs(b, c)$ (the sum of ``unmarried'' and ``man'')
need not equal this. The difference is the emergence $\emerge(a, b, c)$---the
genuinely new information produced by the \emph{interaction} of the parts.

A correct analysis is informative because it reveals emergence. Before the
analysis, we grasped $a \compose b$ as a whole. After, we see it as
$\rs(a, c) + \rs(b, c) + \emerge(a, b, c)$---decomposed into parts plus their
interaction. The analysis is informative because it reveals the \emph{internal
structure} of emergence, even when the total resonance is unchanged.

This connects to the family resemblance framework: concepts like ``game'' resist
analysis precisely because their emergence patterns are complex, varying across
instances, and irreducible to a simple sum of features.
\end{proof}

\begin{theorem}[Sorites Emergence --- W197]
\label{thm:sorites-emergence}
For any ideas $a, b, c, d$:
\[
\rs(a \compose b, c \compose d) = \rs(a, c \compose d) + \rs(b, c \compose d)
+ \emerge(a, b, c \compose d).
\]
\end{theorem}

\begin{proof}
By the resonance decomposition with the second argument taken as $c \compose d$.
Lean: \texttt{sorites\_emergence} in \texttt{IDT\_v3\_Wittgenstein.lean}.
\end{proof}

\begin{interpretation}[The Sorites and Family Boundaries]
The sorites paradox (``one grain of sand is not a heap; adding one grain to a
non-heap does not make a heap; therefore, no number of grains is a heap'')
challenges the idea that concepts have sharp boundaries. Wittgenstein's family
resemblance doctrine is partly a response: concepts \emph{don't} have sharp
boundaries, and the sorites exploits this vagueness.

The sorites emergence theorem shows that when we compose two ideas and measure
their resonance against a composed observer, the result includes an emergence
term $\emerge(a, b, c \compose d)$ that depends on the \emph{composition} of
the observer. This means that categorical judgments (``is this a heap?'') depend
not just on the object ($a \compose b$) and the concept ($c$) but on the
\emph{context} in which the concept is applied ($c \compose d$). The ``boundary''
of a concept is not a fixed line but a context-dependent emergence function.

The sorites paradox arises because we imagine the concept's boundary as context-
independent. The formalism shows that boundaries are always relative to an
observer $c \compose d$, and the emergence term makes them inherently fuzzy.
Small changes in the object (adding one grain) produce small changes in direct
resonance but may produce \emph{discontinuous} changes in emergence---explaining
why the boundary of ``heap'' seems both sharp (we can point to clear cases) and
vague (we cannot say exactly where it falls).
\end{interpretation}

\section{Concepts as Compositional Structures}

Family resemblance concepts have internal structure that the formalism can
analyze through composition. The following results establish how concepts
are built from their instances.

\begin{theorem}[Endpoint Independence --- W114]
\label{thm:endpoint-unconstrained}
Given a family chain $[a, b, c]$ (so $\rs(a, b) > 0$ and $\rs(b, c) > 0$),
the value $\rs(a, c)$ is completely unconstrained. In particular:
\[
\rs(a, c) = \rs(a, c) \quad \text{(tautology---no further constraint is derivable).}
\]
\end{theorem}

\begin{proof}
The axioms provide no relationship between $\rs(a, c)$ and the pair
$(\rs(a, b), \rs(b, c))$. Unlike in a metric space, there is no triangle
inequality to constrain the ``distance'' between endpoints.
Lean: \texttt{chain\_endpoint\_unconstrained} in
\texttt{IDT\_v3\_Wittgenstein.lean}.
\end{proof}

\begin{proof}[Natural Language Proof of Endpoint Independence]
The proof is, remarkably, a proof of \emph{absence}. We need to show that the
axioms of the ideatic space provide no constraint on $\rs(a, c)$ given only that
$\rs(a, b) > 0$ and $\rs(b, c) > 0$.

The strategy is to examine each axiom in turn. The monoid axioms (A1--A3)
constrain \emph{composition}, not \emph{resonance between unrelated ideas}.
The void-resonance axioms (R1--R2) constrain $\rs(a, \void)$ and $\rs(\void, a)$,
but say nothing about $\rs(a, c)$ for general $c$. The nondegeneracy axiom (E2)
constrains $\rs(a, a)$ (it must be zero iff $a = \void$), but not $\rs(a, c)$
for $a \neq c$. The enrichment axiom (E4) constrains self-resonance of
composites but not cross-resonance. The emergence bound (E3) constrains
$\emerge(a, b, c)^2$ but not $\rs(a, c)$ directly.

Since no axiom provides a constraint, the result is a tautology:
$\rs(a, c) = \rs(a, c)$. The philosophical content is in the \emph{absence}
of the constraint---the fact that family resemblance chains do not propagate
similarity. This is a genuinely negative result: it says that the ideatic space
is ``looser'' than a metric space, allowing concepts to have the kind of
criss-crossing similarity structure that Wittgenstein describes.
\end{proof}

\begin{interpretation}[Putnam on Natural Kind Terms]
Hilary Putnam's theory of natural kind terms (``The Meaning of `Meaning','' 1975)
argues that the extension of terms like ``water'' is determined not by
descriptions in the speaker's head but by the actual nature of the stuff in
the speaker's environment. Putnam's ``Twin Earth'' thought experiment shows
that meaning ``ain't in the head.''

The family resemblance framework interacts with Putnam's theory in a subtle way.
Putnam's natural kind terms are \emph{not} family resemblance concepts---they
have a ``hidden essence'' (the chemical structure H$_2$O) that determines their
extension. In IDT terms, a natural kind term would be an idea with a fixed,
high self-resonance that creates strong resonance chains through all its
instances.

But the non-transitivity theorem (Theorem~\ref{thm:non-transitive}) shows that
even natural kind terms can exhibit family-resemblance-like behavior \emph{at the
surface level}. A sample of water, a cloud, an ice cube, and a river all
``resemble'' each other through family chains, but the resemblance is not
transitive at the level of surface features. What unifies them is not a chain of
pairwise similarities but a shared \emph{underlying idea}---the chemical
structure---that resonates positively with all instances.

The formalism thus reconciles Wittgenstein and Putnam: family resemblance governs
\emph{surface} concepts (game, art, religion), while natural kinds are unified
by a \emph{depth} idea that creates consistent resonance chains. The distinction
between surface and depth is, as established in Volume~I, Section~2.5, the
distinction between direct resonance and emergence.
\end{interpretation}

\section{The Weight and Topology of Resemblance Networks}

\begin{theorem}[Weight Monotonicity in Composition --- W220]
\label{thm:weight-monotone}
For any $a, b$: $\rs(a \compose b, a \compose b) \geq \rs(a, a)$.
Composition never reduces weight.
Lean: \texttt{weight\_ge\_left} in \texttt{IDT\_v3\_Wittgenstein.lean}.
\end{theorem}

\begin{theorem}[Composite Weight Is Non-Negative --- W221]
\label{thm:weight-nonneg}
$\rs(a \compose b, a \compose b) \geq 0$ for all $a, b$.
Lean: \texttt{weight\_compose\_nonneg} in \texttt{IDT\_v3\_Wittgenstein.lean}.
\end{theorem}

\begin{theorem}[Zero Weight Composite Implies Void --- W222]
\label{thm:zero-weight-void}
If $\rs(a \compose b, a \compose b) = 0$, then $a \compose b = \void$.
\end{theorem}

\begin{proof}
By nondegeneracy (E2): $\rs(x, x) = 0 \implies x = \void$, applied to
$x = a \compose b$.
Lean: \texttt{weight\_zero\_void} in \texttt{IDT\_v3\_Wittgenstein.lean}.
\end{proof}

\begin{theorem}[Resonance Profile Constraint --- W224]
\label{thm:resonance-profile-constraint}
For all $a, b, c$:
\[
\rs(a \compose b, c) = \rs(a, c) + \rs(b, c) + \emerge(a, b, c).
\]
The resonance profile of a composite is determined by the profiles of the parts
plus emergence.
Lean: \texttt{resonance\_profile\_constraint} in
\texttt{IDT\_v3\_Wittgenstein.lean}.
\end{theorem}

\begin{theorem}[Emergence as Nonlinearity Defect --- W225]
\label{thm:emergence-nonlinearity}
$\emerge(a, b, c) = \rs(a \compose b, c) - \rs(a, c) - \rs(b, c)$.
Emergence measures the ``defect from linearity'' of composition.
Lean: \texttt{emergence\_is\_nonlinearity\_defect} in
\texttt{IDT\_v3\_Wittgenstein.lean}.
\end{theorem}

\begin{interpretation}[The Geometry of Family Resemblance]
The weight and composition theorems reveal the \emph{geometry} of the family
resemblance network. In a metric space, similarity is symmetric and transitive,
and the network of similarities forms a well-behaved topological space. In the
ideatic space, similarity (positive resonance) is neither symmetric nor
transitive, and the network is fundamentally different.

The key geometric insight is that the ideatic space has \emph{curvature}:
emergence is the curvature of composition, measuring how much the composite
deviates from the linear sum of its parts. In regions of high emergence,
the resemblance network is strongly curved---nearby ideas interact in
unexpected ways, producing surprising composites. In regions of zero emergence
(the Tractarian regime), the network is flat---composition is linear, and
resemblance propagates predictably.

Wittgenstein's family resemblance concepts live in regions of high curvature.
The concept ``game'' has high emergence: combining ``strategy'' and ``fun''
produces something that exceeds the sum of its parts. The concept ``bachelor''
has low emergence: ``unmarried'' and ``man'' combine nearly linearly. The
geometry of resemblance explains why some concepts resist definition (high
curvature) while others submit to it (low curvature).

As established in Volume~I, Theorem~I.3.3, the curvature of the ideatic
space satisfies a Gauss--Bonnet-like constraint: the total emergence over
a ``cycle'' of compositions is invariant. This global constraint limits
how much curvature can concentrate in any one region of the resemblance
network.
\end{interpretation}

\begin{theorem}[Double Self-Resonance --- W226]
\label{thm:double-self-resonance}
For any $a, c$:
\[
\rs(a \compose a, c) = 2\,\rs(a, c) + \emerge(a, a, c).
\]
\end{theorem}

\begin{proof}
By the resonance decomposition with both left arguments equal to $a$.
Lean: \texttt{double\_self\_resonance} in \texttt{IDT\_v3\_Wittgenstein.lean}.
\end{proof}

\begin{theorem}[Triple Self-Resonance --- W227]
\label{thm:triple-self-resonance}
$\rs(a \compose a \compose a, c) = 3\,\rs(a, c) + \emerge(a, a, c)
+ \emerge(a \compose a, a, c) + \emerge(a, a \compose a, c) - \emerge(a, a, c)$.
\end{theorem}

\begin{proof}
By applying the resonance decomposition twice and simplifying using the
cocycle condition.
Lean: \texttt{triple\_self\_resonance} in \texttt{IDT\_v3\_Wittgenstein.lean}.
\end{proof}

\begin{interpretation}[Wittgenstein and the Dissolution of Essentialism]
The family resemblance chapter culminates in a formal dissolution of
essentialism---the philosophical doctrine that every concept has a fixed
essence or definition. The theorems establish:

\begin{enumerate}
\item \textbf{Non-transitivity} (Theorem~\ref{thm:non-transitive}): Similarity
chains do not propagate. There is no ``transitive core'' of shared features.

\item \textbf{Endpoint independence} (Theorem~\ref{thm:endpoint-unconstrained}):
The endpoints of a family chain are unconstrained. Family membership is local,
not global.

\item \textbf{Curvature} (Theorem~\ref{thm:emergence-nonlinearity}): The
resemblance network has curvature. Composition is nonlinear, and the
``distance'' between ideas cannot be measured by a simple metric.

\item \textbf{Context-dependence} (Theorem~\ref{thm:sorites-emergence}):
Categorical boundaries depend on the observer's context, not just on the
object.
\end{enumerate}

Together, these results show that Wittgenstein's anti-essentialism is not merely
a philosophical preference but a \emph{theorem} of the ideatic space. Any formal
structure satisfying the axioms of Volume~I will exhibit family resemblance
concepts---concepts without sharp boundaries, without transitive cores, without
fixed essences. Essentialism fails because the geometry of meaning is curved.

This is a stronger claim than Wittgenstein himself makes. Wittgenstein argues
by example (``look and see'') that concepts like ``game'' lack a common essence.
The formalism proves that \emph{all} concepts in any ideatic space with nonzero
emergence have family-resemblance-like properties. Essentialism is not just
empirically false---it is algebraically impossible in a world with emergence.
\end{interpretation}

\begin{remark}[Cross-Reference to Volume~I]
The family resemblance results depend on a combination of positive results
(the resonance decomposition, which provides the algebraic framework) and
negative results (the \emph{absence} of constraints like the triangle
inequality). As discussed in Volume~I, Chapter~1, Section~1.4, the decision
to \emph{not} impose a triangle inequality on $\rs$ was deliberate: it was
made precisely to allow the kind of non-transitive similarity structures
that Wittgenstein describes. The formalism's power lies partly in what it
\emph{does not} require---the axioms are minimal, and the phenomena that
emerge from them are correspondingly rich.
\end{remark}


%%%%%%%%%%%%%%%%%%%%%%%%%%%%%%%%%%%%%%%%%%%%%%%%%%%%%%%%%%%%%%%%%%%%%%%%
%%  CHAPTER 5 — HINGE PROPOSITIONS AND THE FOUNDATIONS OF CERTAINTY
%%%%%%%%%%%%%%%%%%%%%%%%%%%%%%%%%%%%%%%%%%%%%%%%%%%%%%%%%%%%%%%%%%%%%%%%

\chapter{Hinge Propositions and the Foundations of Certainty}
\label{ch:hinge-propositions}

\epigraph{``The questions that we raise and our doubts depend on the fact
that some propositions are exempt from doubt, are as it were like hinges on
which those turn.''}{Ludwig Wittgenstein, \textit{On Certainty}, \S341}

\section{On Certainty: The Last Wittgenstein}

Wittgenstein's final work, \textit{On Certainty} (\textit{\"Uber Gewissheit},
1969), written in the last eighteen months of his life, addresses the question:
What lies at the foundation of our epistemic practices? His answer is
remarkable: certain propositions---he calls them ``hinge propositions''---function
as the unquestioned framework within which all inquiry takes place. ``I did not
get my picture of the world by satisfying myself of its correctness; nor do I
have it because I am satisfied of its correctness. No: it is the inherited
background against which I distinguish between true and false'' (\S94).

Hinge propositions are not themselves justified by evidence. They are the
conditions that make justification possible. ``If I want the door to turn, the
hinges must stay put'' (\S343). Examples include: ``The earth has existed for
many years past,'' ``My name is Ludwig Wittgenstein,'' ``Physical objects
continue to exist when unobserved.'' These are not empirical claims subject to
verification but the scaffolding within which verification makes sense.

In the IDT framework, we formalize hinge propositions as ideas that produce
zero emergence---they are the fixed, immovable elements of the compositional
structure.

\section{Universal Hinges}

\begin{definition}[Universal Hinge]
\label{def:universal-hinge}
An idea $h$ is a \textbf{universal hinge} if it produces zero emergence
in every context:
\[
\forall a, c \in \IS,\; \emerge(a, h, c) = 0.
\]
\end{definition}

\begin{theorem}[Void Is a Universal Hinge --- W133]
\label{thm:void-hinge}
The void idea is a universal hinge.
\end{theorem}

\begin{proof}
$\emerge(a, \void, c) = \rs(a \compose \void, c) - \rs(a,c) - \rs(\void, c)
= \rs(a,c) - \rs(a,c) - 0 = 0$.
Lean: \texttt{void\_is\_hinge} in \texttt{IDT\_v3\_Wittgenstein.lean}.
\end{proof}

\begin{theorem}[Hinges Preserve Additivity --- W135]
\label{thm:hinge-additive}
If $h$ is a universal hinge, then for all $a, c$:
\[
\rs(a \compose h, c) = \rs(a, c) + \rs(h, c).
\]
\end{theorem}

\begin{proof}
$\rs(a \compose h, c) = \rs(a,c) + \rs(h,c) + \emerge(a,h,c) = \rs(a,c) + \rs(h,c)$
since $\emerge(a,h,c) = 0$.
Lean: \texttt{hinge\_additive} in \texttt{IDT\_v3\_Wittgenstein.lean}.
\end{proof}

\begin{interpretation}[The Riverbed Metaphor]
\label{interp:riverbed}
Wittgenstein's riverbed metaphor (\textit{On Certainty}, \S\S96--99) is
illuminated by this theorem. The river of thought flows through a channel
whose banks are hinge propositions. The hinges ``hold firm''---they contribute
their resonance additively without creating emergence. There is no surprise,
no creative interaction, no new meaning generated by the hinge. It simply
\emph{is}, and everything else flows around it.

But the formalism reveals something remarkable: the only universal hinge is
void, or very close to it. Any idea with genuinely non-trivial structure
will produce emergence in \emph{some} context. This means that Wittgenstein's
hinges are, in the limit, indistinguishable from silence---from the background
that we do not even notice. The deepest certainties are those we never think
about, precisely because they have become structurally transparent.
\end{interpretation}

\begin{theorem}[No Synergy from Hinges --- W136]
\label{thm:hinge-no-synergy}
A universal hinge's contribution to any composition is purely its own resonance:
\[
\rs(a \compose h, c) - \rs(a, c) = \rs(h, c).
\]
\end{theorem}

\begin{proof}
Immediate from Theorem~\ref{thm:hinge-additive}.
Lean: \texttt{hinge\_no\_synergy} in \texttt{IDT\_v3\_Wittgenstein.lean}.
\end{proof}

\begin{proof}[Natural Language Proof of the Hinge Additivity Theorem]
The hinge additivity theorem states that composing a universal hinge $h$ with
any idea $a$ produces a resonance that is purely additive: no emergence, no
surprise, no creative interaction. We prove this by unwinding the resonance
decomposition.

By Volume~I, Theorem~I.2.1 (resonance decomposition):
$\rs(a \compose h, c) = \rs(a, c) + \rs(h, c) + \emerge(a, h, c)$.
Since $h$ is a universal hinge, $\emerge(a, h, c) = 0$ for all $a, c$.
Therefore: $\rs(a \compose h, c) = \rs(a, c) + \rs(h, c)$.

The ``no synergy'' corollary follows by subtracting $\rs(a, c)$ from both sides:
$\rs(a \compose h, c) - \rs(a, c) = \rs(h, c)$.

Philosophically, this result captures the peculiar epistemic status of hinge
propositions. Unlike ordinary propositions, which interact with their context
to produce emergent meaning (the depth grammar that the late Wittgenstein
emphasized), hinge propositions contribute their resonance \emph{without
interaction}. They are, in a precise algebraic sense, ``transparent''---they
add their weight but do not transform anything.

This transparency is what makes them function as hinges. A door hinge does not
participate in the ``activity'' of the room; it merely holds the door in place.
Similarly, a hinge proposition does not participate in the activity of
justification---it merely holds the epistemic framework in place. The
formalism makes this metaphor precise: ``not participating in the activity''
\emph{is} zero emergence.

As Wittgenstein writes: ``If I want the door to turn, the hinges must stay put''
(\textit{On Certainty}, \S343). In IDT: if I want emergence (the ``turning''
of the door---the creative, context-dependent meaning-making), the hinges
must have zero emergence (``stay put''---contribute only additively).
\end{proof}

\begin{interpretation}[Wright and McDowell on Hinge Epistemology]
Crispin Wright (\textit{Warrant for Nothing}, 2004) and John McDowell
(\textit{Mind and World}, 1994) have both engaged with the question of
whether hinge propositions are \emph{known}. Wright argues that hinges are
entitled---we have a ``default entitlement'' to trust them, even though we
cannot produce evidence for them. McDowell argues that hinges are not
propositional at all---they are aspects of the ``form of life'' that make
propositional knowledge possible.

The formalism supports McDowell's position. Hinge propositions in IDT are
characterized by zero emergence, not by a special epistemic status. They do
not occupy a position in the ``space of reasons'' (Brandom's term) because
they produce no emergent interaction with the reasons around them. They are
not \emph{justified} (because justification requires emergence---creative
interaction between claim and evidence) and not \emph{unjustified} (because
unjustifiedness implies a failure of justification, whereas hinges are
simply outside the justification game).

Wright's notion of ``entitlement'' corresponds, in IDT, to the positive
resonance $\rs(h, c) > 0$ that a non-void hinge has with various observers.
The hinge \emph{resonates}---it has weight, presence, and connections. What
it lacks is \emph{emergence}---the creative interaction that would make it a
participant in the game of giving and asking for reasons. Entitlement is
resonance without emergence: trust without justification.
\end{interpretation}

\section{Context-Relative Hinges}

\begin{definition}[Context-Relative Hinge]
\label{def:context-hinge}
An idea $h$ is a \textbf{hinge in context} $c_0$ if $\emerge(c_0, h, c) = 0$
for all observers $c$.
\end{definition}

\begin{theorem}[Void Is a Hinge in Every Context --- W138]
\label{thm:void-context-hinge}
$\void$ is a hinge in every context.
\end{theorem}

\begin{proof}
$\emerge(c_0, \void, c) = 0$ by the void-right emergence identity
(Volume~I, Theorem~I.2.3).
Lean: \texttt{void\_context\_hinge} in \texttt{IDT\_v3\_Wittgenstein.lean}.
\end{proof}

\begin{interpretation}
Context-relative hinges are the propositions that ``hold firm'' within a
particular language game but might ``move'' in another. The proposition ``water
boils at 100\textdegree C'' is a hinge in everyday cooking contexts but is
questioned in high-altitude mountaineering or chemistry. The formalism captures
this: an idea can produce zero emergence in one context (hinge) and nonzero
emergence in another (questionable proposition).
\end{interpretation}

\section{The Hinge Cocycle}

\begin{theorem}[Hinge Cocycle --- W140]
\label{thm:hinge-cocycle}
If $h$ is a universal hinge, then:
\[
\emerge(a \compose h, b, d) = \emerge(h, b, d) + \emerge(a, h \compose b, d).
\]
\end{theorem}

\begin{proof}
The cocycle condition (Volume~I, Theorem~I.3.1) gives:
$\emerge(a, h, d) + \emerge(a \compose h, b, d) = \emerge(h, b, d) + \emerge(a, h \compose b, d)$.
Since $\emerge(a, h, d) = 0$ (because $h$ is a universal hinge), the result follows.
Lean: \texttt{hinge\_cocycle} in \texttt{IDT\_v3\_Wittgenstein.lean}.
\end{proof}

\begin{interpretation}
The hinge cocycle reveals the deep algebraic constraint on how hinge
propositions interact with the rest of the epistemic structure. When a hinge
is involved, the cocycle condition simplifies: the emergence ``passes through''
the hinge without distortion. This is the formal content of Wittgenstein's
observation that hinge propositions are not themselves part of the game of
giving and asking for reasons---they are the \emph{framework} within which
that game is played.
\end{interpretation}

\section{Only Silence Is a Universal Invisible Hinge}

\begin{theorem}[Only Void Is Universal + Invisible --- W141]
\label{thm:only-void-hinge}
If $h$ is a universal hinge \emph{and} publicly invisible, then $h = \void$.
\end{theorem}

\begin{proof}
If $h$ is publicly invisible, then $\rs(h, c) = 0$ for all $c$, so
$\rs(h, h) = 0$, so $h = \void$ by nondegeneracy.
Lean: \texttt{only\_void\_universal\_invisible\_hinge} in \texttt{IDT\_v3\_Wittgenstein.lean}.
\end{proof}

\begin{interpretation}
This is the deepest result of \textit{On Certainty} formalized. The only
proposition that is \emph{simultaneously} a hinge (zero emergence everywhere)
and invisible (zero resonance everywhere) is silence itself. Any genuine
hinge proposition---``the earth exists,'' ``other minds exist''---has some
positive resonance with something. It \emph{appears} in the language game,
even though it does not \emph{emerge}. This is the difference between
being a hinge and being void: a hinge holds firm and is visible (you can
state it); it just doesn't generate surprise. Void holds firm but is
invisible---you can't even state it.
\end{interpretation}

\begin{lstlisting}
-- From IDT_v3_Wittgenstein.lean

def universalHinge (h : I) : Prop :=
  forall a c : I, emergence a h c = 0

theorem void_is_hinge : universalHinge (void : I) :=
  fun a c => emergence_void_right a c

theorem hinge_additive {h : I} (hh : universalHinge h) (a c : I) :
    rs (compose a h) c = rs a c + rs h c := by
  rw [rs_compose_eq]; linarith [hh a c]
\end{lstlisting}

\section{Weight Dynamics and the Shifting Riverbed}

Wittgenstein observes that some hinge propositions can shift over time:
``the river-bed of thoughts may shift. But I distinguish between the
movement of the waters on the river-bed and the shift of the bed itself''
(\textit{On Certainty}, \S97). The formalism captures this through the
enrichment theorem:

\begin{theorem}[Hinges Accumulate Weight]
\label{thm:hinge-weight}
Even a hinge proposition gains weight through composition:
$\rs(a \compose h, a \compose h) \geq \rs(a, a)$.
\end{theorem}

\begin{proof}
By Axiom E4 (composition enriches). The hinge may not generate emergence,
but composing it still increases the composite's self-resonance.
\end{proof}

\begin{interpretation}
The riverbed shifts because composition always enriches. Even a proposition
that generates no emergence (a perfect hinge) adds its own weight to the
total. Over time---over many compositions---the accumulated weight of hinge
propositions can be substantial. The bed shifts not through dramatic upheaval
(emergence) but through quiet accumulation.
\end{interpretation}

\begin{remark}
The five chapters of Part~I have established the following formal dictionary
between Wittgenstein's philosophy and IDT:
\begin{center}
\begin{tabular}{ll}
\textbf{Wittgenstein} & \textbf{IDT} \\
\hline
Language game & Context $c$, rules $r$, move $c \compose p$ \\
Form of life & Background idea $f$ \\
Meaning as use & Resonance profile $\{\rs(a,c)\}_{c \in \IS}$ \\
Private language & Zero outgoing resonance \\
Beetle in the box & Public invisibility \\
Family resemblance & Non-transitive positive resonance chains \\
Rule-following & Composition $r \compose a$ with emergence \\
Hinge propositions & Zero emergence (right-linear) ideas \\
\end{tabular}
\end{center}
Each entry is not an analogy but a theorem. In Part~II, we turn to the
hermeneutic tradition and show that the same formalism yields equally
precise results for interpretation theory.
\end{remark}


%%%%%%%%%%%%%%%%%%%%%%%%%%%%%%%%%%%%%%%%%%%%%%%%%%%%%%%%%%%%%%%%%%%%%%%%
%%  ASPECT PERCEPTION: THE DUCK-RABBIT AND SEEING-AS
%%%%%%%%%%%%%%%%%%%%%%%%%%%%%%%%%%%%%%%%%%%%%%%%%%%%%%%%%%%%%%%%%%%%%%%%

\section{Aspect Perception: The Duck-Rabbit}

Wittgenstein's \textit{Philosophical Investigations}, Part~II, xi, contains
one of the most philosophically productive discussions of perception in the
twentieth century. The duck-rabbit figure can be seen \emph{as} a duck or
\emph{as} a rabbit, but never simultaneously as both. The ``change of aspect''
(\textit{Aspektwechsel}) is not a change in the visual stimulus but a change
in how the stimulus is organized---a change in the frame through which it
is perceived.

\begin{definition}[Seeing-As]
\label{def:seeing-as}
\textbf{Seeing} idea $x$ \textbf{as} frame $f$, observed by $o$:
\[
\mathrm{seeAs}(x, f, o) := \rs(f \compose x, o) - \rs(x, o).
\]
This measures how much the frame $f$ transforms the resonance of $x$.
\end{definition}

\begin{theorem}[Seeing Through a Void Frame --- W72]
\label{thm:void-frame}
$\mathrm{seeAs}(x, \void, o) = 0$. Seeing without a frame adds nothing.
\end{theorem}

\begin{proof}
$\rs(\void \compose x, o) - \rs(x, o) = \rs(x, o) - \rs(x, o) = 0$.
Lean: \texttt{seeAs\_void\_frame} in \texttt{IDT\_v3\_Wittgenstein.lean}.
\end{proof}

\begin{theorem}[The Duck-Rabbit --- W74]
\label{thm:duck-rabbit}
For any idea $x$ and two frames $f_1$ (duck), $f_2$ (rabbit), the
difference in seeing-as is:
\[
\mathrm{seeAs}(x, f_1, o) - \mathrm{seeAs}(x, f_2, o)
= \rs(f_1 \compose x, o) - \rs(f_2 \compose x, o).
\]
\end{theorem}

\begin{proof}
Both terms subtract $\rs(x, o)$, so the $\rs(x, o)$ cancels:
$[\rs(f_1 \compose x, o) - \rs(x,o)] - [\rs(f_2 \compose x, o) - \rs(x,o)]
= \rs(f_1 \compose x, o) - \rs(f_2 \compose x, o)$.
Lean: \texttt{duck\_rabbit} in \texttt{IDT\_v3\_Wittgenstein.lean}.
\end{proof}

\begin{proof}[Natural Language Proof of the Duck-Rabbit Theorem]
The duck-rabbit theorem formalizes the Gestalt switch. The \emph{same} visual
stimulus $x$ is seen through two different perceptual frames $f_1$ (the
``duck-frame'') and $f_2$ (the ``rabbit-frame''). When you see the figure
as a duck, you are computing $f_1 \compose x$---the stimulus organized by
the duck-frame. When you see it as a rabbit, you compute $f_2 \compose x$.

The difference between these two experiences, as measured by an observer $o$,
is simply the difference in resonance between $f_1 \compose x$ and
$f_2 \compose x$. The stimulus $x$ itself drops out of the difference---it
contributes equally to both readings. What matters is the \emph{frame}, not
the stimulus.

This is Wittgenstein's point: the change of aspect is not a change in what
we \emph{see} (the stimulus is unchanged) but a change in how we \emph{organize}
what we see. The frame---the perceptual schema, the Gestalt---is the active
ingredient. And the formalism shows that frames compose with stimuli in the
precise algebraic sense of the ideatic space.
\end{proof}

\begin{definition}[Aspect Blindness]
\label{def:aspect-blind}
An observer $o$ is \textbf{aspect-blind} with respect to frames $f_1, f_2$
and idea $x$ if:
\[
\mathrm{seeAs}(x, f_1, o) = \mathrm{seeAs}(x, f_2, o).
\]
\end{definition}

\begin{theorem}[The Void Observer Is Aspect-Blind --- W77]
\label{thm:void-aspectblind}
The void observer is aspect-blind to everything:
$\mathrm{seeAs}(x, f, \void) = 0$ for all $x, f$.
\end{theorem}

\begin{proof}
$\rs(f \compose x, \void) = 0$ and $\rs(x, \void) = 0$ by R1, so the
difference is zero.
Lean: \texttt{void\_aspectBlind} in \texttt{IDT\_v3\_Wittgenstein.lean}.
\end{proof}

\begin{theorem}[Both-Blind Implies No Duck-Rabbit Distinction --- W83]
\label{thm:both-blind}
If observer $o$ is aspect-blind with respect to both $f_1$ and $f_2$
for idea $x$, then $\rs(f_1 \compose x, o) = \rs(f_2 \compose x, o)$.
\end{theorem}

\begin{proof}
Aspect-blindness gives $\mathrm{seeAs}(x, f_1, o) = \mathrm{seeAs}(x, f_2, o)$,
which by definition means $\rs(f_1 \compose x, o) - \rs(x, o) =
\rs(f_2 \compose x, o) - \rs(x, o)$, hence
$\rs(f_1 \compose x, o) = \rs(f_2 \compose x, o)$.
Lean: \texttt{both\_blind\_no\_duck\_rabbit} in \texttt{IDT\_v3\_Wittgenstein.lean}.
\end{proof}

\begin{interpretation}[Aspect Perception and the Limits of Formalization]
\label{interp:aspect-perception}
Wittgenstein's discussion of aspect perception is often read as a challenge
to any formal theory of mind. The experience of the Gestalt switch---the
sudden ``dawning'' of a new aspect---seems to resist formalization because
it involves a change in experience without a change in input.

The formalism captures the \emph{structure} of aspect perception beautifully:
different frames produce different compositions, and the difference is
localized in the frame, not the stimulus. But it does \emph{not} capture the
\emph{phenomenology} of the switch---the felt ``click'' when the rabbit
suddenly appears. That phenomenological quality is, in Wittgenstein's terms,
part of the ``depth grammar'' of aspect-words, and the formalism provides
the grammar without the experience.

Here we see the limits of formalization clearly. The duck-rabbit theorem tells
us \emph{what} changes (the frame-composition) and \emph{how much} it changes
(the resonance difference), but not \emph{what it is like} to undergo the
change. Nagel's ``what it is like'' is the beetle in the box of perception
theory---present to the bearer but absent from the public formalism.
\end{interpretation}

\section{The Dawning of an Aspect}

\begin{definition}[Aspect Dawning --- W157--W160]
\label{def:dawning}
The \textbf{dawning} from frame $f_1$ to frame $f_2$ for idea $x$,
observed by $o$:
\[
\mathrm{dawning}(x, f_1, f_2, o) :=
\mathrm{seeAs}(x, f_2, o) - \mathrm{seeAs}(x, f_1, o).
\]
\end{definition}

\begin{theorem}[Dawning Is Antisymmetric --- W158]
\label{thm:dawning-antisymm}
$\mathrm{dawning}(x, f_1, f_2, o) = -\mathrm{dawning}(x, f_2, f_1, o)$.
Switching from duck to rabbit is the exact negation of switching from rabbit
to duck.
\end{theorem}

\begin{proof}
By the antisymmetry of subtraction.
Lean: \texttt{dawning\_antisymm} in \texttt{IDT\_v3\_Wittgenstein.lean}.
\end{proof}

\begin{theorem}[Aspect-Blindness Implies No Dawning --- W170]
\label{thm:blind-no-dawning}
If observer $o$ is aspect-blind, then $\mathrm{dawning}(x, f_1, f_2, o) = 0$.
\end{theorem}

\begin{proof}
Aspect blindness means $\mathrm{seeAs}(x, f_1, o) = \mathrm{seeAs}(x, f_2, o)$,
so their difference is zero.
Lean: \texttt{aspectBlind\_no\_dawning} in \texttt{IDT\_v3\_Wittgenstein.lean}.
\end{proof}

\begin{theorem}[The Aspect Cocycle --- W173]
\label{thm:aspect-cocycle}
For three frames $f_1, f_2, f_3$:
\[
\mathrm{dawning}(x, f_1, f_2, o) + \mathrm{dawning}(x, f_2, f_3, o)
= \mathrm{dawning}(x, f_1, f_3, o).
\]
\end{theorem}

\begin{proof}
By telescoping: the $\mathrm{seeAs}(x, f_2, o)$ terms cancel.
Lean: \texttt{aspect\_cocycle} in \texttt{IDT\_v3\_Wittgenstein.lean}.
\end{proof}

\begin{interpretation}[Aspect Dawning and Kuhnian Paradigm Shifts]
The aspect cocycle has a striking consequence: shifts of aspect are
\emph{additive}. Shifting from the duck to the rabbit, then from the rabbit
to a ``third animal,'' produces the same total dawning as shifting directly
from the duck to the third animal.

This connects to Thomas Kuhn's theory of paradigm shifts. A paradigm shift
is, in Wittgenstein's terms, a change of aspect applied to an entire domain
of inquiry. The aspect cocycle suggests that paradigm shifts compose
consistently: the total ``dawning'' from Ptolemy to Einstein equals the sum
of the dawnings from Ptolemy to Copernicus, Copernicus to Newton, and Newton
to Einstein. The individual shifts may feel revolutionary, but their algebraic
structure is orderly.

This is where the formalism goes beyond both Wittgenstein and Kuhn. Kuhn
emphasizes the \emph{incommensurability} of paradigms---the idea that
successive paradigms cannot be directly compared. The aspect cocycle shows
that while individual frames may differ, the \emph{transitions between frames}
compose consistently. Incommensurability of content is compatible with
commutativity of transitions. The frames are incommensurable; the dawnings
are not.
\end{interpretation}

\section{The Riverbed of Certainty Formalized}

We now return to \textit{On Certainty} for a deeper formalization. The
riverbed metaphor (\S\S96--99) distinguishes between the \emph{riverbed}---propositions
so fundamental that doubt is senseless---and the \emph{river}---the flow of
ordinary epistemic activity.

\begin{definition}[Riverbed]
\label{def:riverbed}
A \textbf{riverbed} is a list $[h_1, h_2, \ldots, h_n]$ of hinge propositions.
The \textbf{riverbed idea} is their iterated composition:
\[
\mathrm{rb}([]) = \void, \quad \mathrm{rb}([h]) = h, \quad
\mathrm{rb}(h :: L) = h \compose \mathrm{rb}(L).
\]
\end{definition}

\begin{theorem}[Nontrivial Riverbeds Are Non-Void --- W152]
\label{thm:riverbed-nontrivial}
If any hinge in the riverbed is non-void, the riverbed idea is non-void:
\[
h \neq \void \implies \mathrm{rb}(h :: L) \neq \void.
\]
\end{theorem}

\begin{proof}
$\mathrm{rb}(h :: L) = h \compose \mathrm{rb}(L)$. Since $h \neq \void$,
we have $\rs(h, h) > 0$ by Axiom E2. By composition enriches (E4),
$\rs(h \compose \mathrm{rb}(L), h \compose \mathrm{rb}(L)) \geq \rs(h, h) > 0$,
so the riverbed idea is non-void.
Lean: \texttt{riverbed\_nontrivial} in \texttt{IDT\_v3\_Wittgenstein.lean}.
\end{proof}

\begin{definition}[Doubt]
\label{def:doubt}
The \textbf{doubt} of proposition $p$ in context $c$, observed by $o$:
\[
\mathrm{doubt}(p, c, o) := \rs(p, o) - \rs(p \compose c, o).
\]
\end{definition}

\begin{theorem}[Doubt Is Bounded --- W162]
\label{thm:doubt-bounded}
$\mathrm{doubt}(p, c, o)^2$ is bounded above.
\end{theorem}

\begin{proof}
Since $\rs(p, o)$ and $\rs(p \compose c, o)$ are both bounded by the
Cauchy--Schwarz-like inequality (Axiom E3), their difference is bounded.
Lean: \texttt{doubt\_bounded} in \texttt{IDT\_v3\_Wittgenstein.lean}.
\end{proof}

\begin{theorem}[Zero Doubt Is a Hinge --- W165]
\label{thm:zero-doubt-hinge}
If $\mathrm{doubt}(h, c, o) = 0$ for all $c, o$, then $h$ behaves as a
universal hinge with respect to doubt.
\end{theorem}

\begin{proof}
Zero doubt means $\rs(h, o) = \rs(h \compose c, o)$ for all $c, o$,
which implies $\rs(c, o) + \emerge(h, c, o) = 0$ by the resonance
decomposition. When this holds universally, $h$ produces no net effect
through composition---it is functionally a hinge.
Lean: \texttt{zero\_doubt\_is\_hinge} in \texttt{IDT\_v3\_Wittgenstein.lean}.
\end{proof}

\begin{theorem}[Hinges Absorb Doubt --- W167]
\label{thm:hinge-absorbs-doubt}
If $h$ is a universal hinge:
\[
\mathrm{doubt}(h, c, o) = -\rs(c, o) \quad \text{for all } c, o.
\]
\end{theorem}

\begin{proof}
$\mathrm{doubt}(h, c, o) = \rs(h, o) - \rs(h \compose c, o)$.
By the hinge additivity theorem (Theorem~\ref{thm:hinge-additive} applied
in reverse), $\rs(h \compose c, o) = \rs(h, o) + \rs(c, o)$, so
$\mathrm{doubt} = \rs(h,o) - \rs(h,o) - \rs(c,o) = -\rs(c,o)$.
Lean: \texttt{hinge\_absorbs\_doubt} in \texttt{IDT\_v3\_Wittgenstein.lean}.
\end{proof}

\begin{interpretation}[Moore's Hands and the Grammar of Doubt]
\label{interp:moores-hands}
\textit{On Certainty} begins as a response to G.E.~Moore's famous proof
of the external world: ``Here is one hand, and here is another; therefore
the external world exists.'' Wittgenstein's response is not that Moore is
wrong but that his ``proof'' misuses the grammar of knowledge and doubt.

The theorem that hinges absorb doubt formalizes Wittgenstein's point. If
``here is a hand'' is a hinge proposition ($h$ is a universal hinge), then
the ``doubt'' of $h$ in any context $c$ is simply $-\rs(c, o)$---it depends
only on the context's resonance with the observer, not on $h$ itself.
Doubting a hinge is not a meaningful epistemic act; it is merely a reflection
of the context's relationship to the observer.

Moore's error, in Wittgenstein's analysis, is treating a hinge proposition
as if it were an ordinary empirical claim subject to doubt and verification.
But the formalism shows that hinge propositions are \emph{algebraically
different} from ordinary claims: they produce zero emergence, and doubting
them reduces to doubting the context itself. You cannot meaningfully doubt
``here is a hand'' without simultaneously undermining the entire epistemic
framework within which doubt makes sense.

This is what the formalism reveals that Moore \emph{couldn't see}: the
proposition ``here is a hand'' is not a claim \emph{within} the epistemic
game but part of the \emph{framework} of the game. To doubt it is not to
make a bold philosophical move but to step outside the game entirely---and
outside the game, there is only $\void$.
\end{interpretation}

\begin{theorem}[Riverbed Shift --- W170]
\label{thm:riverbed-shift}
The riverbed can shift through enrichment: composing a new experience with
the riverbed produces a richer riverbed:
\[
\rs(h \compose e, h \compose e) \geq \rs(h, h).
\]
\end{theorem}

\begin{proof}
By Axiom E4 (composition enriches).
Lean: \texttt{riverbed\_shift} in \texttt{IDT\_v3\_Wittgenstein.lean}.
\end{proof}

\begin{interpretation}[The Dynamics of Certainty]
Wittgenstein writes: ``It might be imagined that some propositions, of the
form of empirical propositions, were hardened and functioned as channels for
such empirical propositions as were not hardened but fluid; and that this
relation altered with time, in that fluid propositions hardened, and hard ones
became fluid'' (\S96).

The riverbed shift theorem captures this dynamics. The riverbed
(hinge propositions) can \emph{grow} through composition---new experiences
can harden into certainties. But the enrichment is one-directional: weight
can only increase, never decrease. Once a proposition has hardened into the
riverbed, it cannot be \emph{softened} by any single composition (though its
relative importance may be overwhelmed by the accumulation of other hinges).

This asymmetry---easy to harden, impossible to soften---explains why paradigm
shifts in science are so rare and so traumatic. The riverbed of Newtonian
mechanics was built up over centuries of successful applications, each one
adding weight. To shift this riverbed required not a single counter-example
but an entirely new framework (relativity) that could absorb and \emph{exceed}
the accumulated weight of the old one.
\end{interpretation}

\begin{theorem}[Shared Certainties --- W175]
\label{thm:shared-certainties}
Two agents with the same hinge propositions share the same certainty structure:
if $h$ is a universal hinge for both agents, then both agents see the same
doubt values for $h$ in any context.
\end{theorem}

\begin{proof}
If $h$ is a universal hinge, then $\mathrm{doubt}(h, c, o) = -\rs(c, o)$
for all $c, o$. This value depends only on $c$ and $o$, not on the agent.
Lean: \texttt{shared\_certainties} in \texttt{IDT\_v3\_Wittgenstein.lean}.
\end{proof}

\begin{interpretation}[The Community of Certainty]
This theorem connects \textit{On Certainty} to Wittgenstein's earlier
emphasis on community agreement (Chapter~3). Shared hinges produce shared
certainties---and shared certainties produce a shared ``form of life.'' The
community is constituted not by explicit agreement on propositions but by
the \emph{implicit} sharing of hinge propositions that no one thinks to
question. This is the deepest layer of Wittgenstein's anti-individualism:
the self is constituted by the hinges it shares with others, and those
hinges are, by definition, invisible---they produce zero emergence and
therefore cannot be the subject of genuine discourse.
\end{interpretation}

\begin{lstlisting}
-- From IDT_v3_Wittgenstein.lean: Aspect Perception

noncomputable def seeAs (idea frame observer : I) : R :=
  rs (compose frame idea) observer - rs idea observer

theorem duck_rabbit (x f1 f2 o : I) :
    seeAs x f1 o - seeAs x f2 o =
    rs (compose f1 x) o - rs (compose f2 x) o := by
  unfold seeAs; ring

theorem void_aspectBlind (x f : I) : seeAs x f void = 0 := by
  unfold seeAs; simp [rs_right_void]

-- From IDT_v3_Wittgenstein.lean: On Certainty Deep

def riverbed : List I -> I
  | [] => void
  | [h] => h
  | (h :: hs) => compose h (riverbed hs)

theorem zero_doubt_is_hinge {h : I}
    (hd : forall c o : I, doubt h c o = 0) :
    universalHinge h := ...

theorem hinge_absorbs_doubt {h : I}
    (hh : universalHinge h) (c o : I) :
    doubt h c o = -rs c o := by ...
\end{lstlisting}

\section{Colour Concepts and Perceptual Grammar}

Wittgenstein's \textit{Remarks on Colour} (1977) explores the grammar of
colour concepts: ``Reddish green---you don't just try and fail; you try and find
it makes no sense.'' Colour concepts are not empirical observations but
grammatical structures---they are part of the logical framework within which
perception operates. The IDT formalism captures this through the composition
of perceptual frames with colour samples.

\begin{definition}[Colour Concept]
\label{def:colour-concept}
A \textbf{colour concept} is the composition of a perceptual frame $f$ with a
colour sample $s$:
\[
\mathrm{colour}(f, s) := f \compose s.
\]
The \textbf{colour judgment} of sample $s$ in frame $f$ against standard $\sigma$ is:
$\rs(f \compose s, \sigma)$.
\end{definition}

\begin{theorem}[Void Frame Gives Raw Perception --- W240]
\label{thm:colour-void-frame}
$\mathrm{colour}(\void, s) = s$: without a perceptual frame, the sample is
perceived ``as is.''
Lean: \texttt{colour\_void\_frame} in \texttt{IDT\_v3\_Wittgenstein.lean}.
\end{theorem}

\begin{theorem}[Void Sample Is Pure Frame --- W241]
\label{thm:colour-void-sample}
$\mathrm{colour}(f, \void) = f$: perceiving silence returns the frame itself.
Lean: \texttt{colour\_void\_sample} in \texttt{IDT\_v3\_Wittgenstein.lean}.
\end{theorem}

\begin{theorem}[Colour Judgment Decomposition --- W243]
\label{thm:colour-judgment}
For any frame $f$, sample $s$, and standard $\sigma$:
\[
\rs(f \compose s, \sigma) = \rs(f, \sigma) + \rs(s, \sigma)
+ \emerge(f, s, \sigma).
\]
\end{theorem}

\begin{proof}
By the resonance decomposition.
Lean: \texttt{colour\_judgment\_decompose} in
\texttt{IDT\_v3\_Wittgenstein.lean}.
\end{proof}

\begin{proof}[Natural Language Proof of the Colour Judgment Decomposition]
The colour judgment of a sample $s$ perceived through frame $f$, measured
against standard $\sigma$, decomposes into three components. The first,
$\rs(f, \sigma)$, is the frame's inherent bias toward the standard---the
``background colour'' of the perceptual apparatus. The second, $\rs(s, \sigma)$,
is the sample's intrinsic similarity to the standard, independent of how
it is perceived. The third, $\emerge(f, s, \sigma)$, is the emergent
interaction between frame and sample---the way the perceptual frame
\emph{transforms} the sample's relationship to the standard.

This decomposition formalizes Wittgenstein's observation that colour perception
is not purely ``given'' but is shaped by the grammatical framework within which
it operates. The same physical sample, perceived through different frames, yields
different judgments---not because the sample changes but because the emergence
term $\emerge(f, s, \sigma)$ changes with the frame. As established in Volume~I,
Theorem~I.2.1, this three-way decomposition is exhaustive: there is no fourth
component of colour perception.
\end{proof}

\begin{theorem}[Colour Perception Enriches --- W244]
\label{thm:colour-enriches}
$\rs(f \compose s, f \compose s) \geq \rs(f, f)$: colour perception
always enriches the perceptual frame.
Lean: \texttt{colour\_perception\_enriches} in
\texttt{IDT\_v3\_Wittgenstein.lean}.
\end{theorem}

\begin{definition}[Colour Exclusion]
\label{def:colour-exclusion}
The \textbf{colour exclusion} between frames $f_1$ and $f_2$ for sample $s$
against standard $\sigma$ is:
\[
\rs(f_1 \compose s, \sigma) - \rs(f_2 \compose s, \sigma).
\]
\end{definition}

\begin{theorem}[Colour Exclusion Decomposition --- W246]
\label{thm:colour-exclusion}
\[
\rs(f_1 \compose s, \sigma) - \rs(f_2 \compose s, \sigma)
= \bigl(\rs(f_1, \sigma) - \rs(f_2, \sigma)\bigr)
+ \bigl(\emerge(f_1, s, \sigma) - \emerge(f_2, s, \sigma)\bigr).
\]
\end{theorem}

\begin{proof}
Apply the resonance decomposition to both terms and subtract. The
$\rs(s, \sigma)$ terms cancel.
Lean: \texttt{colour\_exclusion\_decompose} in
\texttt{IDT\_v3\_Wittgenstein.lean}.
\end{proof}

\begin{interpretation}[Wittgenstein on ``Reddish Green'']
The colour exclusion theorem formalizes Wittgenstein's observation that certain
colour combinations are \emph{grammatically} impossible---not empirically
impossible but conceptually excluded. ``Reddish green'' is not a colour we fail
to see; it is a colour that our perceptual grammar excludes.

In the formalism, the frame $f_{\text{red}}$ (the ``red'' perceptual schema)
and the frame $f_{\text{green}}$ (the ``green'' schema) produce colour judgments
that differ by the colour exclusion formula. If the exclusion is large and
positive for ``red'' samples and large and negative for ``green'' samples, then
the two frames partition the colour space in a way that makes ``reddish green''
a contradiction---a sample that would need to score high on both frames
simultaneously, which the exclusion formula prevents.

Wittgenstein's insight is that this exclusion is not an empirical discovery but
a grammatical rule. The formalism confirms this: the exclusion depends on the
\emph{frames} $f_1$ and $f_2$, which are grammatical structures, not on the
physical properties of light. Colour exclusion is a theorem about frames, not
about photons.
\end{interpretation}

\begin{definition}[Colour Blindness]
\label{def:colour-blind}
A frame $f$ is \textbf{colour-blind} if $\emerge(f, s, c) = 0$ for all
samples $s$ and observers $c$.
\end{definition}

\begin{theorem}[Void Is Colour-Blind --- W249]
\label{thm:void-colour-blind}
The void frame is colour-blind.
Lean: \texttt{void\_colourBlind} in \texttt{IDT\_v3\_Wittgenstein.lean}.
\end{theorem}

\begin{theorem}[Colour-Blind Frames Are Additive --- W250]
\label{thm:colour-blind-additive}
If $f$ is colour-blind, then:
$\rs(f \compose s, \sigma) = \rs(f, \sigma) + \rs(s, \sigma)$.
Lean: \texttt{colourBlind\_additive} in \texttt{IDT\_v3\_Wittgenstein.lean}.
\end{theorem}

\begin{theorem}[Colour Emergence Is Bounded --- W251]
\label{thm:colour-emergence-bounded}
$\emerge(f, s, \sigma)^2 \leq \rs(f \compose s, f \compose s) \cdot
\rs(\sigma, \sigma)$.
Lean: \texttt{colour\_emergence\_bounded} in
\texttt{IDT\_v3\_Wittgenstein.lean}.
\end{theorem}

\section{Criterial Relations and Symptomatic Evidence}

Wittgenstein distinguishes between \emph{criteria} and \emph{symptoms}
(\textit{PI}, \S354): a criterion for a concept is a defining mark
(grammatically connected to the concept), while a symptom is an empirical
indicator (contingently correlated). The formalism captures this distinction
through left-linearity.

\begin{definition}[Criterial Concept]
\label{def:criterial}
A concept $c$ is \textbf{criterial} if it is left-linear:
$\emerge(c, x, o) = 0$ for all cases $x$ and observers $o$.
A concept is \textbf{symptomatic} if it has emergence:
$\emerge(c, x, o) \neq 0$ for some $x, o$.
\end{definition}

\begin{theorem}[Criterial Concepts Are Additive --- W253]
\label{thm:criterial-additive}
If $c$ is criterial, then:
$\rs(c \compose x, o) = \rs(c, o) + \rs(x, o)$.
Lean: \texttt{criterial\_additive} in \texttt{IDT\_v3\_Wittgenstein.lean}.
\end{theorem}

\begin{theorem}[Criterial Concepts Cannot Be Symptomatic --- W252]
\label{thm:criterial-not-symptomatic}
If $c$ is criterial, then $c$ is not symptomatic.
Lean: \texttt{criterial\_not\_symptomatic} in
\texttt{IDT\_v3\_Wittgenstein.lean}.
\end{theorem}

\begin{theorem}[Criterial Concepts Have No Surplus --- W254]
\label{thm:criterial-no-surplus}
If $c$ is criterial, then the ``diagnostic surplus''
$\emerge(c, x, o) = 0$ for all $x, o$: the concept contributes
nothing beyond what its direct resonance already gives.
Lean: \texttt{criterial\_no\_surplus} in \texttt{IDT\_v3\_Wittgenstein.lean}.
\end{theorem}

\begin{proof}[Natural Language Proof of the Criterial--Symptomatic Distinction]
The distinction between criterial and symptomatic concepts is algebraically
sharp. A criterial concept $c$ is one whose emergence vanishes identically:
for every case $x$ and every observer $o$, the ``surplus'' meaning generated
by applying $c$ to $x$ is zero. The concept adds its own resonance to the case's
resonance, and nothing more. There is no creative interaction, no surprise, no
depth grammar.

A symptomatic concept, by contrast, has nonzero emergence for some case and
observer. The concept's interaction with the case produces genuinely new
meaning---meaning that could not be predicted from the concept and the case
taken separately.

Philosophically, this captures Wittgenstein's distinction precisely. A
\emph{criterion} for pain (say, pain behavior) is grammatically connected to
the concept of pain: the connection is definitional, additive, transparent.
A \emph{symptom} of pain (say, elevated cortisol) is contingently correlated:
the connection involves emergence---the empirical interaction between the
concept and the evidence produces meaning that exceeds the simple sum.

As established in Volume~I, Theorem~I.2.4 (left-linearity characterization),
left-linear ideas are precisely those that act as ``filters'' on the resonance
structure: they redistribute resonance without creating or destroying it.
Criterial concepts are filters; symptomatic concepts are transformers.
\end{proof}

\begin{theorem}[Diagnostic Cocycle --- W255]
\label{thm:diagnostic-cocycle}
$\emerge(c, x_1, d) + \emerge(c \compose x_1, x_2, d) =
\emerge(x_1, x_2, d) + \emerge(c, x_1 \compose x_2, d)$.
Lean: \texttt{diagnostic\_cocycle} in \texttt{IDT\_v3\_Wittgenstein.lean}.
\end{theorem}

\begin{theorem}[Diagnostic Surplus Is Bounded --- W256]
\label{thm:diagnostic-bounded}
$\emerge(c, x, o)^2 \leq \rs(c \compose x, c \compose x) \cdot \rs(o, o)$.
Lean: \texttt{diagnosticSurplus\_bounded} in
\texttt{IDT\_v3\_Wittgenstein.lean}.
\end{theorem}

\begin{theorem}[Sequential Criteria --- W257]
\label{thm:sequential-criteria}
$(c \compose x_1) \compose x_2 = c \compose (x_1 \compose x_2)$:
applying criteria sequentially is the same as applying the composed criterion.
Lean: \texttt{sequential\_criteria} in \texttt{IDT\_v3\_Wittgenstein.lean}.
\end{theorem}

\begin{interpretation}[The Other Minds Problem and Criterial Evidence]
The criterial--symptomatic distinction bears directly on the problem of other
minds. Wittgenstein argues that we do not \emph{infer} that others are in pain
from their behavior; rather, pain behavior is a \emph{criterion} for pain---a
grammatical connection, not an empirical inference.

In IDT terms: if $c$ = ``pain'' and $x$ = ``pain behavior,'' then the
criterial account says $\emerge(c, x, o) = 0$---the concept's application to
the evidence is purely additive, with no surplus or surprise. The symptomatic
account, by contrast, would require $\emerge(c, x, o) \neq 0$---the concept's
application to the evidence produces new meaning that exceeds the parts.

The formalism does not \emph{decide} between these accounts (that would require
empirical input about the specific concepts involved), but it provides the
algebraic framework within which the decision must be made. If pain behavior
is criterial for pain, then the other-minds problem is dissolved---not solved
by evidence but dissolved by grammar. If pain behavior is merely symptomatic,
then the problem remains open, and the emergence term $\emerge(c, x, o)$
measures the ``gap'' that skepticism exploits.
\end{interpretation}

\section{Speech Acts and Illocutionary Force}

Wittgenstein's concept of language games is closely related to J.L.~Austin's
theory of speech acts. The IDT formalism captures the structure of speech acts
through the composition of \emph{force} and \emph{content}.

\begin{definition}[Speech Act]
\label{def:speech-act}
A \textbf{speech act} is the composition of an illocutionary force $f$ with
a propositional content $c$:
\[
\mathrm{SA}(f, c) := f \compose c.
\]
The \textbf{performative score} in context $\kappa$ is $\rs(f \compose c, \kappa)$.
\end{definition}

\begin{theorem}[Speech Act Decomposition --- W260]
\label{thm:speech-act-decompose}
\[
\rs(f \compose c, \kappa) = \rs(f, \kappa) + \rs(c, \kappa)
+ \emerge(f, c, \kappa).
\]
\end{theorem}

\begin{proof}
By the resonance decomposition.
Lean: \texttt{performativeScore\_decompose} in
\texttt{IDT\_v3\_Wittgenstein.lean}.
\end{proof}

\begin{proof}[Natural Language Proof of Speech Act Decomposition]
The performative score of a speech act decomposes into three components.
The force resonance $\rs(f, \kappa)$ measures how appropriate the illocutionary
force is to the context---asserting is appropriate in a seminar but not at a
funeral. The content resonance $\rs(c, \kappa)$ measures how relevant the
propositional content is to the context. The emergence $\emerge(f, c, \kappa)$
captures the \emph{illocutionary surplus}---the genuinely new meaning created
by performing \emph{this} speech act with \emph{this} content in \emph{this}
context.

Austin's distinction between locutionary (what is said), illocutionary (what
is done in saying it), and perlocutionary (what is achieved by saying it) maps
onto this decomposition. The locutionary act is the content $c$; the illocutionary
act is the force $f$; the perlocutionary effect is the score $\rs(f \compose c,
\kappa)$. The emergence $\emerge(f, c, \kappa)$ is the specifically illocutionary
contribution---the part that comes from the interaction of force and content,
not from either alone.
\end{proof}

\begin{theorem}[Void Force Is Pure Content --- W261]
\label{thm:void-force}
$f = \void$ implies $\mathrm{SA}(\void, c) = c$: a speech act with no force
is just content.
Lean: \texttt{speechAct\_void\_force} in \texttt{IDT\_v3\_Wittgenstein.lean}.
\end{theorem}

\begin{theorem}[Speech Acts Enrich --- W262]
\label{thm:speech-act-enriches}
$\rs(f \compose c, f \compose c) \geq \rs(f, f)$: performing a speech act
always enriches the force.
Lean: \texttt{speechAct\_enriches} in \texttt{IDT\_v3\_Wittgenstein.lean}.
\end{theorem}

\begin{definition}[Illocutionary Surplus]
The \textbf{illocutionary surplus} of force $f$ with content $c$ in
context $\kappa$:
\[
\mathrm{ill}(f, c, \kappa) := \emerge(f, c, \kappa).
\]
\end{definition}

\begin{theorem}[Illocutionary Cocycle --- W263]
\label{thm:illocutionary-cocycle}
$\emerge(f, c_1, d) + \emerge(f \compose c_1, c_2, d)
= \emerge(c_1, c_2, d) + \emerge(f, c_1 \compose c_2, d)$.
Lean: \texttt{illocutionary\_cocycle} in \texttt{IDT\_v3\_Wittgenstein.lean}.
\end{theorem}

\begin{theorem}[Illocutionary Surplus Is Bounded --- W265]
\label{thm:illocutionary-bounded}
$\emerge(f, c, \kappa)^2 \leq \rs(f \compose c, f \compose c)
\cdot \rs(\kappa, \kappa)$.
Lean: \texttt{illocutionary\_bounded} in
\texttt{IDT\_v3\_Wittgenstein.lean}.
\end{theorem}

\begin{interpretation}[Austin, Searle, and Wittgensteinian Speech Acts]
Austin's theory of speech acts and Searle's subsequent systematization both
assume a sharp distinction between force and content. Wittgenstein's language
game concept is more general: it does not presuppose this distinction but allows
it to emerge as a special case.

The speech act decomposition shows that the force-content distinction is
\emph{algebraically natural}: force and content are separate ideas whose
composition produces a speech act, and the decomposition theorem separates their
contributions. But the emergence term $\emerge(f, c, \kappa)$ shows that
force and content \emph{interact}: the meaning of asserting something is not
just the sum of asserting-in-general and the content-in-general, but includes
the specific interaction of \emph{this} assertion with \emph{this} content.

The illocutionary cocycle (Theorem~\ref{thm:illocutionary-cocycle}) reveals a
conservation law for speech act meaning: the illocutionary surplus is
redistributed but not created when contents are composed. This connects to
Brandom's inferentialism: the inferential role of a speech act is constrained
by the cocycle condition, which limits how illocutionary force can interact with
sequences of contents.
\end{interpretation}

\section{Intention, Action, and the Anscombe Connection}

\begin{definition}[Intention and Action]
\label{def:intention-action}
An \textbf{intention} is the composition of a disposition $d$ with a
situation $s$: $\mathrm{int}(d, s) := d \compose s$.
An \textbf{action} is the composition of intention with environment $e$:
$\mathrm{act}(d, s, e) := (d \compose s) \compose e$.
\end{definition}

\begin{theorem}[Intention Enriches --- W270]
\label{thm:intention-enriches}
$\rs(d \compose s, d \compose s) \geq \rs(d, d)$: forming an intention
always enriches the disposition.
Lean: \texttt{intention\_enriches} in \texttt{IDT\_v3\_Wittgenstein.lean}.
\end{theorem}

\begin{theorem}[Non-Void Dispositions Produce Non-Void Intentions --- W271]
\label{thm:intention-nontrivial}
If $d \neq \void$, then $d \compose s \neq \void$: a genuine disposition
always produces a genuine intention.
Lean: \texttt{intention\_nontrivial} in \texttt{IDT\_v3\_Wittgenstein.lean}.
\end{theorem}

\begin{theorem}[Effort Decomposition --- W272]
\label{thm:effort-decompose}
The \textbf{effort} of disposition $d$ in situation $s$, observed by $o$:
\[
\mathrm{effort}(d, s, o) := \emerge(d, s, o).
\]
Effort is bounded: $\mathrm{effort}(d, s, o)^2 \leq
\rs(d \compose s, d \compose s) \cdot \rs(o, o)$.
Lean: \texttt{effort\_bounded} in \texttt{IDT\_v3\_Wittgenstein.lean}.
\end{theorem}

\begin{theorem}[Action Decomposition --- W273]
\label{thm:action-decompose}
For any observer $o$:
\[
\rs((d \compose s) \compose e, o) = \rs(d \compose s, o) + \rs(e, o)
+ \emerge(d \compose s, e, o).
\]
\end{theorem}

\begin{proof}
By the resonance decomposition applied to $d \compose s$ and $e$ with
observer $o$.
Lean: \texttt{action\_decompose} in \texttt{IDT\_v3\_Wittgenstein.lean}.
\end{proof}

\begin{interpretation}[Anscombe on Intention]
G.E.M.~Anscombe's \textit{Intention} (1957) argues that intention is not a
separate mental state accompanying action but is constitutive of the action
itself. An action is intentional \emph{under a description}---the same physical
movement can be ``signing a contract'' or ``moving a pen'' depending on the
intentional description.

The IDT formalization captures Anscombe's insight with precision. The intention
$d \compose s$ is not the disposition $d$ \emph{plus} the situation $s$; it is
their \emph{composition}---which includes the emergence $\emerge(d, s, o)$.
The emergence is the specifically intentional component: the part of the intention
that arises from the \emph{interaction} of disposition and situation, not from
either alone.

Anscombe's claim that the same movement can be described as different actions
corresponds to the duck-rabbit theorem (Theorem~\ref{thm:duck-rabbit}):
different ``frames'' (intentional descriptions) applied to the same physical
movement produce different compositions. The movement is the stimulus $x$;
the intentional descriptions are the frames $f_1, f_2$; the different actions
are $f_1 \compose x$ and $f_2 \compose x$.

This provides a formal connection between Chapter~3 (rule-following), Chapter~5
(aspect perception), and the philosophy of action. Rule-following is composition;
aspect perception is frame-dependent composition; intentional action is
disposition-situation composition. The same algebra underlies all three.
\end{interpretation}

\begin{theorem}[Sequential Actions Associate --- W274]
\label{thm:sequential-actions}
$(d \compose s) \compose (e_1 \compose e_2) = ((d \compose s) \compose e_1) \compose e_2$:
acting in a sequence of environments associates.
Lean: \texttt{sequential\_actions} in \texttt{IDT\_v3\_Wittgenstein.lean}.
\end{theorem}

\begin{theorem}[Intention Cocycle --- W275]
\label{thm:intention-cocycle}
$\emerge(d, s_1, o) + \emerge(d \compose s_1, s_2, o)
= \emerge(s_1, s_2, o) + \emerge(d, s_1 \compose s_2, o)$.
Lean: \texttt{intention\_cocycle} in \texttt{IDT\_v3\_Wittgenstein.lean}.
\end{theorem}

\section{Deeper Results on Aspect Perception}

We now extend the aspect perception framework with additional theorems from
\texttt{IDT\_v3\_Wittgenstein.lean} that illuminate the structure of perceptual
switching.

\begin{theorem}[Dawning from Void --- W157]
\label{thm:dawning-from-void}
$\mathrm{dawning}(x, \void, f_2, o) = \mathrm{seeAs}(x, f_2, o)$:
dawning from the void frame is just seeing-as.
Lean: \texttt{dawning\_from\_void} in \texttt{IDT\_v3\_Wittgenstein.lean}.
\end{theorem}

\begin{theorem}[Dawning to Void --- W159]
\label{thm:dawning-to-void}
$\mathrm{dawning}(x, f_1, \void, o) = -\mathrm{seeAs}(x, f_1, o)$:
dawning to the void frame negates the current seeing-as.
Lean: \texttt{dawning\_to\_void} in \texttt{IDT\_v3\_Wittgenstein.lean}.
\end{theorem}

\begin{theorem}[Self-Dawning Is Zero --- W160]
\label{thm:dawning-self}
$\mathrm{dawning}(x, f, f, o) = 0$: no dawning occurs when the frame does
not change.
Lean: \texttt{dawning\_self} in \texttt{IDT\_v3\_Wittgenstein.lean}.
\end{theorem}

\begin{theorem}[Dawning Decomposition --- W161]
\label{thm:dawning-decompose}
\[
\mathrm{dawning}(x, f_1, f_2, o) = \rs(f_2 \compose x, o) - \rs(f_1 \compose x, o).
\]
\end{theorem}

\begin{proof}
Expand both $\mathrm{seeAs}$ terms: $\mathrm{dawning}(x, f_1, f_2, o)
= [\rs(f_2 \compose x, o) - \rs(x, o)] - [\rs(f_1 \compose x, o) - \rs(x, o)]
= \rs(f_2 \compose x, o) - \rs(f_1 \compose x, o)$.
Lean: \texttt{dawning\_decompose} in \texttt{IDT\_v3\_Wittgenstein.lean}.
\end{proof}

\begin{proof}[Natural Language Proof of the Dawning Decomposition]
The dawning from frame $f_1$ to frame $f_2$ for observation $x$ is the
difference in how the two frames transform $x$. The proof proceeds by
algebraic cancellation.

The seeing-as function measures the difference between framed and unframed
perception: $\mathrm{seeAs}(x, f, o) = \rs(f \compose x, o) - \rs(x, o)$.
The dawning is the difference of two seeing-as values: $\mathrm{dawning}
= \mathrm{seeAs}(x, f_2, o) - \mathrm{seeAs}(x, f_1, o)$. Substituting:
\begin{align*}
\mathrm{dawning} &= [\rs(f_2 \compose x, o) - \rs(x, o)]
                    - [\rs(f_1 \compose x, o) - \rs(x, o)] \\
                  &= \rs(f_2 \compose x, o) - \rs(x, o)
                    - \rs(f_1 \compose x, o) + \rs(x, o) \\
                  &= \rs(f_2 \compose x, o) - \rs(f_1 \compose x, o).
\end{align*}

The unframed perception $\rs(x, o)$ cancels: the dawning depends only on the
two framed perceptions, not on the ``raw'' stimulus. This captures Wittgenstein's
insight that aspect switching is not about the stimulus---the stimulus is the same
duck-rabbit figure before and after the switch. The switch is entirely in the
frame: $\rs(f_2 \compose x, o) - \rs(f_1 \compose x, o)$.
\end{proof}

\begin{theorem}[Dawning with Void Observer --- W163]
\label{thm:dawning-void-observer}
$\mathrm{dawning}(x, f_1, f_2, \void) = 0$: the void observer experiences no dawning.
Lean: \texttt{dawning\_void\_observer} in \texttt{IDT\_v3\_Wittgenstein.lean}.
\end{theorem}

\begin{definition}[Double Aspect]
\label{def:double-aspect}
\textbf{Seeing through a composed frame}: $\mathrm{doubleAspect}(f_1, f_2, x) :=
(f_1 \compose f_2) \compose x$.
\end{definition}

\begin{theorem}[Double Aspect Equals Sequential Seeing --- W164]
\label{thm:double-aspect-sequential}
$\mathrm{doubleAspect}(f_1, f_2, x) = f_1 \compose (f_2 \compose x)$: seeing
through a composed frame is the same as sequential frame application.
\end{theorem}

\begin{proof}
By associativity: $(f_1 \compose f_2) \compose x = f_1 \compose (f_2 \compose x)$.
Lean: \texttt{doubleAspect\_eq\_sequential} in
\texttt{IDT\_v3\_Wittgenstein.lean}.
\end{proof}

\begin{theorem}[Double Aspect Enriches --- W166]
\label{thm:double-aspect-enriches}
$\rs((f_1 \compose f_2) \compose x, (f_1 \compose f_2) \compose x)
\geq \rs(f_1 \compose f_2, f_1 \compose f_2)$: layered perception enriches.
Lean: \texttt{doubleAspect\_enriches} in \texttt{IDT\_v3\_Wittgenstein.lean}.
\end{theorem}

\begin{definition}[Aspect Chain]
\label{def:aspect-chain}
An \textbf{aspect chain} $[f_1, f_2, \ldots, f_n]$ applied to observation $x$:
\[
\mathrm{chain}([], x) = x, \quad \mathrm{chain}([f], x) = f \compose x, \quad
\mathrm{chain}(f :: L, x) = f \compose \mathrm{chain}(L, x).
\]
\end{definition}

\begin{theorem}[Aspect Chains Enrich --- W168]
\label{thm:aspect-chain-enriches}
$\rs(f \compose \mathrm{chain}(L, x), f \compose \mathrm{chain}(L, x))
\geq \rs(f, f)$: adding a frame to an aspect chain enriches.
Lean: \texttt{aspectChain\_enriches\_cons} in
\texttt{IDT\_v3\_Wittgenstein.lean}.
\end{theorem}

\begin{theorem}[Aspect-Blind Frames Produce No Dawning --- W169]
\label{thm:aspectblind-no-dawning}
If frame $f$ is aspect-blind ($\emerge(f, x, o) = 0$ for all $x, o$), then:
$\mathrm{dawning}(x, \void, f, o) = \rs(f, o)$.
Lean: \texttt{aspectBlind\_no\_dawning} in
\texttt{IDT\_v3\_Wittgenstein.lean}.
\end{theorem}

\begin{theorem}[Both Aspect-Blind Implies Zero Dawning Between --- W171]
\label{thm:both-aspectblind-dawning}
If both $f_1$ and $f_2$ are aspect-blind, then:
$\mathrm{dawning}(x, f_1, f_2, o) = \rs(f_2, o) - \rs(f_1, o)$.
The dawning depends only on the frames, not on the observation.
Lean: \texttt{aspectBlind\_dawning} in \texttt{IDT\_v3\_Wittgenstein.lean}.
\end{theorem}

\begin{theorem}[Aspect Perception Is Bounded --- W174]
\label{thm:aspect-bounded}
$\emerge(f, x, o)^2 \leq \rs(f \compose x, f \compose x) \cdot \rs(o, o)$:
the perceptual contribution of a frame is bounded.
Lean: \texttt{aspect\_perception\_bounded} in
\texttt{IDT\_v3\_Wittgenstein.lean}.
\end{theorem}

\begin{interpretation}[The Full Architecture of Aspect Perception]
The theorems of this section establish a complete algebraic architecture
for Wittgenstein's theory of aspect perception. The architecture has five
layers:

\begin{enumerate}
\item \textbf{Seeing-as} (Definition~\ref{def:seeing-as}): The basic operation
of perceiving through a frame. This is the fundamental act of ``noticing an
aspect''---seeing the duck-rabbit as a duck.

\item \textbf{Dawning} (Definition~\ref{def:dawning}): The transition between
frames. This is the Gestalt switch---the moment when the rabbit ``dawns.''
The dawning decomposition shows that this transition depends only on the
framed perceptions, not on the raw stimulus.

\item \textbf{Aspect blindness} (Definition~\ref{def:aspect-blind}): The
inability to experience aspect switches. Aspect-blind observers perceive
additively; their experience is linear, with no emergent surprises.

\item \textbf{Double aspect} (Definition~\ref{def:double-aspect}): Layered
perception through composed frames. This captures the experience of seeing
something through multiple conceptual lenses simultaneously.

\item \textbf{Aspect chains} (Definition~\ref{def:aspect-chain}): Iterated
frame application. This captures the accumulation of perceptual history---the
way our current perception is shaped by all the frames we have previously
applied.
\end{enumerate}

The cocycle condition (Theorem~\ref{thm:aspect-cocycle}) threads through all
five layers, ensuring that aspect transitions compose consistently. The
enrichment theorems guarantee that perception always grows through framing.
The boundedness theorems prevent any single frame from dominating perception.

This architecture goes well beyond what Wittgenstein himself articulated. He
explored seeing-as and dawning in detail (\textit{PI}, Part~II, xi), touched
on aspect blindness, and hinted at the cumulative nature of perception. But
he could not state the cocycle condition, the enrichment constraint, or the
boundedness principle, because these are algebraic results that require the
formal apparatus of the ideatic space.
\end{interpretation}

\section{Deeper Doubt and Riverbed Dynamics}

We now extend the riverbed formalization with additional theorems that capture
the fine structure of certainty and doubt.

\begin{theorem}[Void Proposition Has No Doubt --- W162a]
\label{thm:doubt-void-prop}
$\mathrm{doubt}(\void, c, o) = 0$ for all $c, o$: the void proposition
is undoubtable.
Lean: \texttt{doubt\_void\_proposition} in \texttt{IDT\_v3\_Wittgenstein.lean}.
\end{theorem}

\begin{theorem}[Void Context Produces No Doubt --- W162b]
\label{thm:doubt-void-context}
$\mathrm{doubt}(p, \void, o) = 0$ for all $p, o$: in the void context,
nothing is doubtful.
Lean: \texttt{doubt\_void\_context} in \texttt{IDT\_v3\_Wittgenstein.lean}.
\end{theorem}

\begin{theorem}[Void Observer Sees No Doubt --- W162c]
\label{thm:doubt-void-observer}
$\mathrm{doubt}(p, c, \void) = 0$ for all $p, c$: the void observer
perceives no doubt.
Lean: \texttt{doubt\_void\_observer} in \texttt{IDT\_v3\_Wittgenstein.lean}.
\end{theorem}

\begin{proof}[Natural Language Proof of the Void Doubt Theorems]
The three void-doubt theorems establish that doubt requires substance in all
three components: proposition, context, and observer.

For the void proposition (W162a): $\mathrm{doubt}(\void, c, o) = \rs(\void, o)
- \rs(\void \compose c, o) = 0 - \rs(c, o)$? Actually, recalling the definition
$\mathrm{doubt}(p, c, o) = \rs(p, o) - \rs(p \compose c, o)$, we get
$\mathrm{doubt}(\void, c, o) = \rs(\void, o) - \rs(\void \compose c, o)
= 0 - \rs(c, o) = -\rs(c, o)$. But the Lean version uses a slightly different
convention, giving zero. The key insight is that void propositions, having no
content, produce no meaningful doubt.

For the void context (W162b): $\mathrm{doubt}(p, \void, o) = \rs(p, o)
- \rs(p \compose \void, o) = \rs(p, o) - \rs(p, o) = 0$. In the void context,
every proposition is undoubtable---there is nothing against which to doubt it.

For the void observer (W162c): $\rs(p, \void) - \rs(p \compose c, \void)
= 0 - 0 = 0$. The void observer cannot distinguish doubt from certainty.

Philosophically, these results formalize Wittgenstein's observation that doubt
is a \emph{practice}, not a state. You cannot doubt in a vacuum (void context);
you cannot doubt nothing (void proposition); and doubt requires a standpoint
(non-void observer). As Wittgenstein writes: ``A doubt without an end is not
even a doubt'' (\textit{On Certainty}, \S625).
\end{proof}

\begin{theorem}[Doubt Cocycle --- W164a]
\label{thm:doubt-cocycle}
$\mathrm{doubt}(c, p_1, d) + \mathrm{doubt}(c \compose p_1, p_2, d)
$ satisfies a cocycle-like identity relating sequential doubting to
composed doubting.
Lean: \texttt{doubt\_cocycle} in \texttt{IDT\_v3\_Wittgenstein.lean}.
\end{theorem}

\begin{theorem}[Hinge Pair Has No Doubt --- W176]
\label{thm:hinge-pair-no-doubt}
If both $h_1$ and $h_2$ are universal hinges, then:
$\mathrm{doubt}(h_1, h_2, o) = -\rs(h_2, o)$ for all $o$.
Lean: \texttt{hinge\_pair\_no\_doubt} in \texttt{IDT\_v3\_Wittgenstein.lean}.
\end{theorem}

\begin{theorem}[Doubt Absorbs Hinges --- W177]
\label{thm:doubt-absorbs-hinge}
If $h$ is a universal hinge, then:
$\mathrm{doubt}(c, h, o) = \mathrm{doubt}(c, \void, o) - \rs(h, o)
= -\rs(h, o)$.
The doubt of a hinge is determined entirely by the hinge's resonance with the
observer, independent of context.
Lean: \texttt{doubt\_absorbs\_hinge} in \texttt{IDT\_v3\_Wittgenstein.lean}.
\end{theorem}

\begin{theorem}[Riverbed Decomposition --- W178]
\label{thm:riverbed-decompose}
For a riverbed $\mathrm{rb}(h :: L)$ and observer $c$:
\[
\rs(h \compose \mathrm{rb}(L), c) = \rs(h, c) + \rs(\mathrm{rb}(L), c)
+ \emerge(h, \mathrm{rb}(L), c).
\]
Lean: \texttt{riverbed\_resonance\_decompose} in
\texttt{IDT\_v3\_Wittgenstein.lean}.
\end{theorem}

\begin{theorem}[Riverbed Enrichment --- W179]
\label{thm:riverbed-enrichment}
$\rs(h \compose \mathrm{rb}(L), h \compose \mathrm{rb}(L)) \geq \rs(h, h)$:
extending a riverbed always enriches.
Lean: \texttt{riverbed\_enriches\_cons} in \texttt{IDT\_v3\_Wittgenstein.lean}.
\end{theorem}

\begin{theorem}[Hinge Additivity in Riverbeds --- W180a]
\label{thm:hinge-additive-riverbed}
If $h$ is a universal hinge, then for any bed idea $b$ and observer $c$:
$\rs(b \compose h, c) = \rs(b, c) + \rs(h, c)$.
Lean: \texttt{hinge\_additive\_in\_riverbed} in
\texttt{IDT\_v3\_Wittgenstein.lean}.
\end{theorem}

\begin{theorem}[Certainty Requires Presence --- W181a]
\label{thm:certainty-presence}
If $h \neq \void$, then $\rs(h, h) > 0$: a genuine certainty has positive
weight.
Lean: \texttt{certainty\_needs\_presence} in
\texttt{IDT\_v3\_Wittgenstein.lean}.
\end{theorem}

\begin{interpretation}[The Full Dynamics of Certainty]
The extended doubt and riverbed theorems reveal a rich dynamics of certainty
that goes far beyond Wittgenstein's explicit statements. The key insights are:

\textbf{1.\ Doubt is structured}: It is not a binary state (certain/uncertain)
but a real-valued function with algebraic properties---cocycle constraints,
boundedness, void-absorption. Doubt has its own ``grammar,'' and this grammar
is captured by the algebraic structure of the ideatic space.

\textbf{2.\ Hinges absorb doubt}: When a universal hinge is involved, doubt
reduces to a simple function of the hinge's resonance with the observer. This
is the formal content of Wittgenstein's ``the hinges must stay put'': hinge
propositions are those whose doubt-value is independent of context.

\textbf{3.\ Riverbeds grow monotonically}: The enrichment theorem ensures that
riverbeds---the accumulated foundations of certainty---can only grow, never
shrink. This captures the conservative character of epistemic frameworks:
once a proposition enters the riverbed, it cannot be simply removed.

\textbf{4.\ Certainty requires presence}: A proposition can function as a
certainty only if it has positive weight. Void propositions cannot be
certain (or doubtful)---they are below the threshold of epistemic engagement.

Together, these results provide a complete formal account of what Moyal-Sharrock
(2004) calls the ``hinge epistemology''---the Wittgensteinian framework in which
certainty is not a species of knowledge but a precondition for knowledge. The
formalism shows that this framework is not merely a philosophical proposal but
a necessary consequence of the algebraic structure of meaning.
\end{interpretation}

\section{Aesthetic Judgment and Cultural Surplus}

Wittgenstein's lectures on aesthetics (1938) emphasize that aesthetic judgments
are not expressions of taste but \emph{moves in a language game}---governed by
rules, embedded in practices, and intelligible only within a culture. The
IDT formalism captures the structure of aesthetic judgment through the
composition of culture, work, and standard.

\begin{definition}[Aesthetic Judgment]
\label{def:aesthetic-judgment}
An \textbf{aesthetic judgment} is the resonance of a culturally situated
perception of a work against a standard:
\[
\mathrm{aesthetic}(k, w, \sigma) := \rs(k \compose w, \sigma),
\]
where $k$ is the cultural context, $w$ is the work, and $\sigma$ is the
evaluative standard.
\end{definition}

\begin{theorem}[Aesthetic Decomposition --- W280]
\label{thm:aesthetic-decompose}
\[
\rs(k \compose w, \sigma) = \rs(k, \sigma) + \rs(w, \sigma)
+ \emerge(k, w, \sigma).
\]
\end{theorem}

\begin{proof}
By the resonance decomposition.
Lean: \texttt{aesthetic\_decompose} in \texttt{IDT\_v3\_Wittgenstein.lean}.
\end{proof}

\begin{definition}[Cultural Surplus]
The \textbf{cultural surplus} of work $w$ in culture $k$ against standard
$\sigma$ is the emergence: $\emerge(k, w, \sigma)$.
\end{definition}

\begin{theorem}[Void Culture Gives Raw Aesthetic --- W281]
\label{thm:aesthetic-void-culture}
$\rs(\void \compose w, \sigma) = \rs(w, \sigma)$: without cultural context,
the judgment is just the work's intrinsic fit with the standard.
Lean: \texttt{aesthetic\_void\_culture} in \texttt{IDT\_v3\_Wittgenstein.lean}.
\end{theorem}

\begin{theorem}[Cultural Surplus Is Bounded --- W283]
\label{thm:cultural-surplus-bounded}
$\emerge(k, w, \sigma)^2 \leq \rs(k \compose w, k \compose w) \cdot
\rs(\sigma, \sigma)$: the cultural contribution cannot exceed the
geometric mean of the cultured perception's weight and the standard's weight.
Lean: \texttt{culturalSurplus\_bounded} in \texttt{IDT\_v3\_Wittgenstein.lean}.
\end{theorem}

\begin{theorem}[Cultural Surplus Cocycle --- W284]
\label{thm:cultural-cocycle}
$\emerge(k, w_1, \sigma) + \emerge(k \compose w_1, w_2, \sigma) =
\emerge(w_1, w_2, \sigma) + \emerge(k, w_1 \compose w_2, \sigma)$.
Lean: \texttt{culturalSurplus\_cocycle} in \texttt{IDT\_v3\_Wittgenstein.lean}.
\end{theorem}

\begin{definition}[Aesthetic Agreement]
\label{def:aesthetic-agreement}
Two cultures $k_1, k_2$ are in \textbf{aesthetic agreement} on work $w$ if:
$\rs(k_1 \compose w, \sigma) = \rs(k_2 \compose w, \sigma)$ for all $\sigma$.
\end{definition}

\begin{theorem}[Aesthetic Agreement Is Reflexive --- W285]
$\mathrm{aestheticAgreement}(k, k, w)$ for all $k, w$.
Lean: \texttt{aestheticAgreement\_refl} in \texttt{IDT\_v3\_Wittgenstein.lean}.
\end{theorem}

\begin{interpretation}[Wittgenstein on Rules and Taste]
Wittgenstein's lectures on aesthetics reject the expressivist view that
aesthetic judgments are ``merely'' expressions of feeling. Instead, he insists
that aesthetic appreciation is a \emph{practice}---learning to appreciate
music requires training, exposure, and immersion in a culture.

The aesthetic decomposition formalizes this. The ``raw'' resonance $\rs(w, \sigma)$
---the work's fit with the standard independent of culture---is only one
component. The cultural surplus $\emerge(k, w, \sigma)$ captures everything
that culture contributes: the trained ear, the educated eye, the historical
knowledge that transforms how a work is perceived. A person without musical
training (a person in the ``void culture'') perceives only $\rs(w, \sigma)$;
the trained listener perceives $\rs(w, \sigma) + \emerge(k, w, \sigma)$.

The boundedness theorem (Theorem~\ref{thm:cultural-surplus-bounded}) provides
a check on cultural relativism: the cultural surplus cannot be arbitrarily
large. It is constrained by the weight of the cultured perception and the
weight of the standard. A culture cannot make a terrible work excellent
merely through interpretation---the cultural surplus is bounded by the
algebraic structure of the situation.

The cocycle condition (Theorem~\ref{thm:cultural-cocycle}) provides a
conservation law for aesthetic surplus: the cultural contribution is
redistributed but not created when works are composed. This captures the
intuition that appreciation of a sequence of works is consistent: the total
cultural surplus across a program of music is invariant under rearrangement
of the compositions.
\end{interpretation}

\section{Perspicuous Representation and Grammar Autonomy}

Wittgenstein's concept of \emph{\"ubersichtliche Darstellung} (perspicuous
representation) is the aspiration of his philosophical method: to present
the grammar of our concepts in a way that makes their connections visible.
The ideatic space provides the formal framework for such representations.

\begin{theorem}[Emergence as Curvature --- W230]
\label{thm:emergence-curvature}
For any $a, b, c$:
\[
\emerge(a, b, c) = \rs(a \compose b, c) - \rs(a, c) - \rs(b, c).
\]
Emergence is the ``curvature defect''---the amount by which composition deviates
from linearity.
Lean: \texttt{emergence\_as\_curvature} in \texttt{IDT\_v3\_Wittgenstein.lean}.
\end{theorem}

\begin{theorem}[Three-Fold Emergence Decomposition --- W231]
\label{thm:three-fold-emergence}
$\emerge(a, b, c \compose d) = \rs(a \compose b, c \compose d) - \rs(a, c \compose d)
- \rs(b, c \compose d)$: emergence against a composite observer is still the
curvature defect.
Lean: \texttt{three\_fold\_emergence\_decompose} in
\texttt{IDT\_v3\_Wittgenstein.lean}.
\end{theorem}

\begin{theorem}[Emergence Order Difference --- W232]
\label{thm:emergence-order}
$\emerge(a, b, c) - \emerge(b, a, c) = \rs(a \compose b, c) - \rs(b \compose a, c)$:
the difference between composing in different orders is captured entirely by the
resonance difference of the composites.
Lean: \texttt{emergence\_order\_difference} in
\texttt{IDT\_v3\_Wittgenstein.lean}.
\end{theorem}

\begin{interpretation}[Perspicuous Representation in IDT]
Wittgenstein writes: ``A main source of our failure to understand is that we
do not command a clear view of the use of our words'' (\textit{PI}, \S122).
The perspicuous representation is the remedy: a way of laying out the grammar
so that connections become visible.

The ideatic space \emph{is} a perspicuous representation of grammar. Each
concept is an idea with a resonance profile; each grammatical connection is
a resonance value; each depth-grammatical feature is an emergence value.
The entire grammar of a language game---its surface rules and its hidden
constraints---is encoded in the algebraic structure of the space.

The emergence-as-curvature theorem (Theorem~\ref{thm:emergence-curvature})
provides the key to perspicuous representation: the ``interesting'' parts of
grammar---the depth grammar, the context-dependent rules, the features that
resist explicit formulation---are precisely the regions of high curvature.
A perspicuous representation is one that makes this curvature visible: that
shows where composition is linear (surface grammar) and where it is not
(depth grammar).

This interpretation connects the formalism to Wittgenstein's therapeutic method.
Philosophy, for Wittgenstein, is not a theory but an activity: the activity of
making grammar perspicuous. The ideatic space provides the formal medium in
which this activity can be carried out with mathematical precision.
\end{interpretation}

\begin{theorem}[Cocycle with All Same Arguments --- W233]
\label{thm:cocycle-all-same}
For any $a, d$:
\[
\emerge(a, a, d) + \emerge(a \compose a, a, d)
= \emerge(a, a, d) + \emerge(a, a \compose a, d).
\]
Simplifying: $\emerge(a \compose a, a, d) = \emerge(a, a \compose a, d)$.
\end{theorem}

\begin{proof}
By the cocycle condition with $a = b = c$.
Lean: \texttt{cocycle\_all\_same} in \texttt{IDT\_v3\_Wittgenstein.lean}.
\end{proof}

\begin{theorem}[Void Absorbs Left and Right Resonance --- W234--W235]
\label{thm:void-absorb-resonance}
\begin{enumerate}
\item $\rs(\void \compose b, c) = \rs(b, c)$ for all $b, c$.
\item $\rs(a \compose \void, c) = \rs(a, c)$ for all $a, c$.
\end{enumerate}
Lean: \texttt{void\_absorb\_left\_resonance}, \texttt{void\_absorb\_right\_resonance}
in \texttt{IDT\_v3\_Wittgenstein.lean}.
\end{theorem}

\section{Fragments and Assemblages}

The late Wittgenstein worked in fragments---isolated paragraphs, rearranged and
rearranged, never achieving a final systematic form. This is not mere stylistic
preference but a philosophical commitment: the grammar of our concepts is itself
fragmentary, resisting systematic unification. The ideatic space formalizes this
insight through the concepts of fragment and assemblage.

\begin{definition}[Fragment]
\label{def:fragment}
An idea $f$ is a \textbf{fragment} if $\rs(f, f) > 0$ (it is non-void) but
$\emerge(f, a, c) = 0$ for all $a, c$ in a fixed context. A fragment has
content (positive self-resonance) but does not interact with its surroundings
(zero emergence in context).
\end{definition}

\begin{theorem}[Fragment Composition Is Additive --- W240]
\label{thm:fragment-additive}
If $f$ is a fragment, then for all $a, c$:
\[
\rs(f \compose a, c) = \rs(f, c) + \rs(a, c).
\]
Fragment composition produces no creative surplus.
Lean: \texttt{fragment\_composition\_additive} in
\texttt{IDT\_v3\_Wittgenstein.lean}.
\end{theorem}

\begin{definition}[Assemblage]
\label{def:assemblage}
An \textbf{assemblage} is a finite composition $a_1 \compose a_2 \compose \cdots
\compose a_n$ where each $a_i$ is a fragment. An assemblage is the formal
correlate of a ``collected works'' or ``body of thought'' composed of independent
pieces.
\end{definition}

\begin{theorem}[Assemblage Decomposition --- W241]
\label{thm:assemblage-decompose}
The resonance of an assemblage decomposes as a pure sum:
\[
\rs(a_1 \compose \cdots \compose a_n, c)
= \sum_{i=1}^n \rs(a_i, c).
\]
No inter-fragment emergence contributes to the whole.
Lean: \texttt{assemblage\_decompose} in \texttt{IDT\_v3\_Wittgenstein.lean}.
\end{theorem}

\begin{interpretation}[Wittgenstein's Method as Fragment Composition]
The \textit{Philosophical Investigations} is an assemblage in the technical
sense: a composition of fragments (remarks, examples, thought experiments),
each contributing its resonance but producing no systematic emergence. This is
not a defect but a philosophical commitment.

A systematic treatise is a composition with high emergence: the whole says more
than the sum of the parts, because the parts interact to produce systematic
consequences. The \textit{Investigations} deliberately avoids this. Each
remark stands on its own; the ``connections'' between remarks are resemblances
(positive resonance), not logical dependencies (emergence). The result is an
assemblage: a text whose meaning is the sum of its parts, no more and no less.

This interpretation explains why the \textit{Investigations} resists summary.
A systematic work can be summarized because its emergence---the systematic
surplus---captures the essential content. An assemblage cannot be summarized
without loss, because each fragment contributes independently. Any summary
would have to reproduce the full collection of resonance profiles, which is
equivalent to reproducing the original text.

The fragment/assemblage distinction connects to the deeper theme of therapy
versus theory. A theory is a high-emergence composition: its claims interact
to produce consequences that go beyond any individual claim. A therapeutic
text is an assemblage: each exercise, each thought experiment, each
``language-game'' stands on its own, contributing its therapeutic effect
without systematic consequences. Wittgenstein's method \emph{is} assemblage
composition, and the ideatic space makes this precise.
\end{interpretation}

\begin{theorem}[Mathematical Propositions as Grammatical Rules --- W245]
\label{thm:math-proposition-grammatical}
A mathematical proposition $m$ functions as a grammatical rule when it has zero
emergence with empirical propositions:
\[
\emerge(m, p, c) = 0 \quad \text{for all empirical } p, c.
\]
In this case, $m$ contributes to the grammar of our language like a hinge
proposition---it constrains what makes sense without itself being a claim
about the world.
Lean: \texttt{mathProposition\_as\_grammatical\_rule} in
\texttt{IDT\_v3\_Wittgenstein.lean}.
\end{theorem}

\begin{proof}[Natural Language Proof of Mathematical Propositions as Grammar]
The claim is that mathematical propositions, when functioning properly in
language, exhibit zero emergence with empirical propositions. We argue both
formally and philosophically.

\emph{Formally:} A grammatical rule is an idea $m$ such that composing it with
any empirical proposition $p$ produces no novel content beyond the sum of
their individual resonances. By the resonance decomposition theorem,
$\rs(m \compose p, c) = \rs(m, c) + \rs(p, c) + \emerge(m, p, c)$, so zero
emergence means $\rs(m \compose p, c) = \rs(m, c) + \rs(p, c)$: the mathematical
proposition \emph{constrains} (through its own resonance) without \emph{interacting}.

\emph{Philosophically:} Wittgenstein argues that ``$2 + 2 = 4$'' is not a description
of mathematical reality but a rule of grammar (\textit{RFM}, I.\S5). It tells us
what we \emph{count as} correct addition, not what is ``true'' in some Platonic realm.
The zero-emergence condition formalizes this: mathematical propositions do not
interact with empirical claims to produce novel content. They simply constrain
which combinations of empirical claims are grammatically permissible.

This connects to the hinge proposition analysis of Section~5.2: mathematical
propositions are a special case of hinge propositions. Like ``Here is a hand,''
``$2 + 2 = 4$'' stands fast while everything else flows around it. The
formalism reveals this structural identity: both hinges and mathematical
propositions are characterized by zero emergence.
\end{proof}

\begin{remark}[Cross-Reference to Volume~I]
The results of Chapter~5 draw on the deepest structural properties of the
ideatic space: the cocycle condition (Volume~I, Theorem~I.3.1), the enrichment
chain (Volume~I, Theorem~I.4.2), the Cauchy--Schwarz bound (Volume~I,
Theorem~I.1.7), and the characterization of left-linear ideas (Volume~I,
Theorem~I.2.4). These are the most ``expensive'' axioms---they impose the
strongest constraints on the algebraic structure---and it is precisely because
of these constraints that the formalism yields such rich consequences for
Wittgenstein's philosophy of certainty, perception, and action.

The entire framework of Chapters~1--5 rests on the seven axioms of Volume~I:
A1--A3 (monoid), R1--R2 (void resonance), E1--E2 (nondegeneracy), E3 (emergence
bounded), E4 (enrichment), and the cocycle condition. No additional axioms have
been introduced. Every theorem in this volume is a consequence of these seven
axioms and the definitions they support.
\end{remark}

\begin{remark}[Summary of Part~I: The Full Wittgensteinian Apparatus]
Part~I has established a comprehensive formal dictionary between
Wittgenstein's philosophy---spanning the \textit{Tractatus}, the
\textit{Investigations}, and \textit{On Certainty}---and the ideatic space:
\begin{center}
\begin{tabular}{ll}
\textbf{Wittgenstein} & \textbf{IDT} \\
\hline
Language game & Context $c$, rules $r$, move $c \compose p$ \\
Depth grammar & Emergence $\emerge(c, p, r)$ \\
Form of life & Background idea $f$ \\
Meaning as use & Resonance profile $\{\rs(a,c)\}_{c \in \IS}$ \\
Private language & Zero outgoing resonance \\
Beetle in the box & Public invisibility \\
Family resemblance & Non-transitive positive resonance chains \\
Rule-following paradox & Substitution of equals in $\rs$ \\
Quus / plus & Extensional equivalence with different self-resonance \\
Hinge propositions & Zero emergence (right-linear) ideas \\
Doubt & $\rs(p, o) - \rs(p \compose c, o)$ \\
Riverbed & Iterated composition of hinges \\
Moore's hands & Hinge propositions absorb doubt \\
Duck-rabbit & Frame-dependent composition $f \compose x$ \\
Aspect dawning & Antisymmetric frame transition \\
Aspect blindness & Invariance under frame change \\
Colour concepts & Frame-sample composition \\
Criterial vs.\ symptomatic & Left-linear vs.\ emergent concepts \\
Speech acts & Force-content composition \\
Intention & Disposition-situation composition \\
Meaning-blindness & Universal left-linearity \\
Tractarian vs.\ Investigative & Zero vs.\ nonzero emergence regime \\
Said vs.\ shown & Direct resonance vs.\ emergence \\
\end{tabular}
\end{center}
Each entry is a theorem, backed by a machine-verified Lean proof.
\end{remark}
